% PLane Transformations — Further Mathematics Upper Sixth — Cours signé OnBuch+
% Official MINESEC syllabus. Compiler avec : tectonic FMA03-plane-transformations.tex
\newcommand{\DOCMATIERE}{Further Mathematics}
\newcommand{\DOCNIVEAU}{Upper Sixth}
\newcommand{\DOCMODULE}{Upper Sixth — Further mathematics}
\newcommand{\DOCLECON}{3}
\newcommand{\DOCTITRE}{PLane Transformations}
% OnBuch+ — préambule « cours signés » ENRICHI (compilé avec tectonic / XeLaTeX)
% Système visuel « Pop » d'OnBuch+ : fond crème (jamais blanc), bordures encre
% épaisses, ombres « dures » décalées, titres Archivo Black en capitales.
% Version MATHÉMATIQUES : graphes (pgfplots/tikz), boîtes pédagogiques
% (définition, propriété, méthode, attention, le savais-tu, exercice résolu,
% exercice, TP, bilan), page de garde et signature OnBuch+ ancrée en bas.
%
% Macros attendues dans main.tex AVANT \input{preamble} :
%   \DOCMATIERE  \DOCNIVEAU  \DOCMODULE  \DOCLECON (numéro)  \DOCTITRE
\documentclass[11pt]{article}

\usepackage{fontspec}
\usepackage[english]{babel}
\usepackage[a4paper,margin=2cm,top=3.4cm,bottom=2.8cm,headheight=1.4cm,footskip=1.3cm]{geometry}
\usepackage{amsmath,amssymb,amsfonts}
\usepackage{mathtools} % \xmapsto, \coloneqq... souvent utilisés par les modèles
\usepackage{cancel} % simplifications barrées (bilans de forces)
% \abs{z} (module, valeur absolue) et \norm{u} (norme), utilisés par les modèles
\providecommand{\abs}{}\let\abs\relax\DeclarePairedDelimiter{\abs}{\lvert}{\rvert}
\providecommand{\norm}{}\let\norm\relax\DeclarePairedDelimiter{\norm}{\lVert}{\rVert}
\usepackage[table]{xcolor} % [table] : fournit aussi \rowcolors (alternance de lignes)
\usepackage{pagecolor}
\usepackage{tikz}
\usetikzlibrary{arrows.meta,positioning,calc,shapes.geometric,shapes.misc,shapes.arrows,decorations.pathmorphing,decorations.pathreplacing,decorations.markings,patterns,shadows,mindmap,trees,fit,backgrounds,matrix,intersections,angles,quotes}
\usepackage{pgfplots}
\usepackage[version=4]{mhchem} % équations nucléaires \ce{^{235}_{92}U}
\usepackage[european,siunitx]{circuitikz} % schémas électriques
\pgfplotsset{compat=1.18}
\usepgfplotslibrary{fillbetween,groupplots} % aires entre courbes (calcul intégral)
\usepackage{enumitem}
\usepackage{fancyhdr}
% [most] charge toutes les bibliothèques tcolorbox (listings, minted, raster,
% documentation...) et gonfle inutilement la mémoire TeX sur un document long
% (~300 boîtes) : on ne charge que ce qui est réellement utilisé.
\usepackage[skins,breakable]{tcolorbox}
\usepackage{graphicx}
\usepackage{colortbl}
\usepackage{booktabs}
\usepackage{tabularx}
\usepackage{diagbox} % case d'en-tête diagonale des tableaux à double entrée
% Colonnes extensibles centrée / gauche / droite (C, L, R), souvent utilisées par les modèles
\newcolumntype{C}{>{\centering\arraybackslash}X}
\newcolumntype{L}{>{\raggedright\arraybackslash}X}
\newcolumntype{R}{>{\raggedleft\arraybackslash}X}
\usepackage{array}
\usepackage{multicol}
\usepackage{siunitx}
% parse-numbers=false : accepte aussi \SI{2,5 \times 10^{-3}}{...}
\sisetup{locale=UK,output-decimal-marker={.},group-minimum-digits=4,parse-numbers=false}
\DeclareSIUnit\ohm{\ensuremath{\symup{\Omega}}} % U+2126 absent de la police
\DeclareSIUnit\division{div} % réglages d'oscilloscope (ms/div, V/div)
\DeclareSIUnit\tr{tr} % tours (vitesse de rotation, ex. \SI{1500}{\tr\per\minute})
\usepackage{qrcode}
% \degree : commande fréquemment utilisée par les modèles pour un angle ou une
% température en dehors de \SI{}{} (non fournie par les paquets chargés).
\providecommand{\degree}{^{\circ}}
% \qed (fin de démonstration), utilisé par les modèles sans amsthm
\providecommand{\qed}{\hfill$\square$}
% \barre : double barre des valeurs interdites dans un tableau de variations (tkz-tab)
\providecommand{\barre}{\ensuremath{\|}}
% environnement proof (amsthm non chargé) utilisé par les modèles
\ifdefined\proof\else\newenvironment{proof}[1][Proof]{\par\noindent\textit{#1.}\ }{\qed\par}\fi
\usepackage{lastpage}
\usepackage{hyperref}
\hypersetup{hidelinks}

% ============================================================
% Palette « Pop » OnBuch+ — mêmes hex que la classe P (plus_ui.dart)
% ============================================================
\definecolor{poporange}{HTML}{F59321}
\definecolor{popdark}{HTML}{E07A0C}
\definecolor{popcream}{HTML}{FFF6E8}
\definecolor{popink}{HTML}{1C1714}
\definecolor{poppurple}{HTML}{7A5AE0}
\definecolor{popgreen}{HTML}{2E9E6A}
\definecolor{poppink}{HTML}{E0457B}
\definecolor{popblue}{HTML}{2F7DE1}
\definecolor{popgold}{HTML}{C98A12}
\definecolor{popmuted}{HTML}{8A7F76}
\definecolor{popline}{HTML}{EADFD2}
\definecolor{popgray}{HTML}{8A7F76}
\colorlet{popgrey}{popgray}
% Alias tolérants (le modèle nomme parfois les couleurs différemment)
\colorlet{popwhite}{white}
\colorlet{popblack}{popink}
\colorlet{popred}{poppink}
\colorlet{popyellow}{popgold}
% Teintes claires pour fonds de tableaux / zones
\colorlet{popblueL}{popblue!10!white}
\colorlet{poporangeL}{poporange!14!white}
\colorlet{popgreenL}{popgreen!12!white}
\colorlet{poppurpleL}{poppurple!12!white}
\colorlet{poppinkL}{poppink!12!white}
\colorlet{popgoldL}{popgold!14!white}

\pagecolor{popcream}

% ============================================================
% Typographie
% ============================================================
\defaultfontfeatures{Ligatures=TeX}
\setmainfont{PlusJakartaSans-Medium.ttf}[Path=../../../fonts/,BoldFont=PlusJakartaSans-Bold.ttf]
% Maths en sans-serif (Fira Math) avec chiffres et lettres droites de Plus
% Jakarta, pour que les formules se fondent dans le texte.
\usepackage[math-style=ISO,bold-style=ISO]{unicode-math}
\setmathfont{FiraMath-Regular.otf}[Path=../../../fonts/]
\setmathfont{PlusJakartaSans-Medium.ttf}[Path=../../../fonts/,range={up/num,up/{Latin,latin},\mathup}]
% (icomma retiré : décimales avec point en anglais)
% Caractères mathématiques Unicode absents des polices de texte (Plus Jakarta,
% Archivo Black) : rendus par leur équivalent mathématique, en texte comme en
% formule — sinon ils s'affichent en carré vide (ex. « Arithmétique dans ℤ »).
\usepackage{newunicodechar}
\newunicodechar{ℝ}{\ensuremath{\mathbb{R}}}
\newunicodechar{ℤ}{\ensuremath{\mathbb{Z}}}
\newunicodechar{ℂ}{\ensuremath{\mathbb{C}}}
\newunicodechar{ℕ}{\ensuremath{\mathbb{N}}}
\newunicodechar{ℚ}{\ensuremath{\mathbb{Q}}}
\newunicodechar{𝔻}{\ensuremath{\mathbb{D}}}
\newunicodechar{α}{\ensuremath{\alpha}}
\newunicodechar{β}{\ensuremath{\beta}}
\newunicodechar{γ}{\ensuremath{\gamma}}
\newunicodechar{δ}{\ensuremath{\delta}}
\newunicodechar{ε}{\ensuremath{\varepsilon}}
\newunicodechar{θ}{\ensuremath{\theta}}
\newunicodechar{λ}{\ensuremath{\lambda}}
\newunicodechar{μ}{\ensuremath{\mu}}
\newunicodechar{σ}{\ensuremath{\sigma}}
\newunicodechar{τ}{\ensuremath{\tau}}
\newunicodechar{φ}{\ensuremath{\varphi}}
\newunicodechar{ψ}{\ensuremath{\psi}}
\newunicodechar{ω}{\ensuremath{\omega}}
\newunicodechar{Σ}{\ensuremath{\Sigma}}
% majuscules produites par \MakeUppercase dans les titres (α-aminés -> Α)
\newunicodechar{Α}{\ensuremath{\alpha}}
\newunicodechar{Β}{\ensuremath{\beta}}
\newunicodechar{Γ}{\ensuremath{\Gamma}}
\newunicodechar{Δ}{\ensuremath{\Delta}}
\newunicodechar{Ε}{\ensuremath{\varepsilon}}
\newunicodechar{Θ}{\ensuremath{\Theta}}
\newunicodechar{Λ}{\ensuremath{\Lambda}}
\newunicodechar{Μ}{\ensuremath{\mu}}
\newunicodechar{Τ}{\ensuremath{\tau}}
\newunicodechar{Φ}{\ensuremath{\Phi}}
\newunicodechar{Ψ}{\ensuremath{\Psi}}
\newunicodechar{Ω}{\ensuremath{\Omega}}
\newunicodechar{↦}{\ensuremath{\mapsto}}
\newunicodechar{∈}{\ensuremath{\in}}
\newunicodechar{∉}{\ensuremath{\notin}}
\newunicodechar{∧}{\ensuremath{\wedge}}
\newunicodechar{⇔}{\ensuremath{\Leftrightarrow}}
\newunicodechar{⟺}{\ensuremath{\Longleftrightarrow}}
\newunicodechar{⇒}{\ensuremath{\Rightarrow}}
\newunicodechar{⟹}{\ensuremath{\Longrightarrow}}
\newunicodechar{∘}{\ensuremath{\circ}}
\newunicodechar{∪}{\ensuremath{\cup}}
\newunicodechar{∩}{\ensuremath{\cap}}
\newunicodechar{≡}{\ensuremath{\equiv}}
\newunicodechar{⊂}{\ensuremath{\subset}}
\newunicodechar{⊥}{\ensuremath{\perp}}
\newunicodechar{∃}{\ensuremath{\exists}}
\newunicodechar{∀}{\ensuremath{\forall}}
\newunicodechar{∛}{\ensuremath{\sqrt[3]{\,}}}
\newunicodechar{∜}{\ensuremath{\sqrt[4]{\,}}}
\newunicodechar{⃗}{\ensuremath{{}^{\to}}}
\newunicodechar{ⁿ}{\ifmmode^{n}\else\textsuperscript{\ensuremath{n}}\fi}
\newunicodechar{ˣ}{\ifmmode^{x}\else\textsuperscript{\ensuremath{x}}\fi}
\newunicodechar{ᵇ}{\ifmmode^{b}\else\textsuperscript{\ensuremath{b}}\fi}
\newunicodechar{ʳ}{\ifmmode^{r}\else\textsuperscript{\ensuremath{r}}\fi}
\newunicodechar{ᵘ}{\ifmmode^{u}\else\textsuperscript{\ensuremath{u}}\fi}
\newunicodechar{ʸ}{\ifmmode^{y}\else\textsuperscript{\ensuremath{y}}\fi}
\newunicodechar{ᵃ}{\ifmmode^{a}\else\textsuperscript{\ensuremath{a}}\fi}
\newunicodechar{ᵉ}{\ifmmode^{e}\else\textsuperscript{\ensuremath{e}}\fi}
\newunicodechar{ᵏ}{\ifmmode^{k}\else\textsuperscript{\ensuremath{k}}\fi}
\newunicodechar{ᵖ}{\ifmmode^{p}\else\textsuperscript{\ensuremath{p}}\fi}
\newunicodechar{ᴺ}{\ifmmode^{N}\else\textsuperscript{\ensuremath{N}}\fi}
\newunicodechar{ᶜ}{\ifmmode^{c}\else\textsuperscript{\ensuremath{c}}\fi}
\newunicodechar{ʰ}{\ifmmode^{h}\else\textsuperscript{\ensuremath{h}}\fi}
\newunicodechar{ᵠ}{\ifmmode^{q}\else\textsuperscript{\ensuremath{q}}\fi}
\newunicodechar{ₙ}{\ifmmode_{n}\else\textsubscript{\ensuremath{n}}\fi}
\newunicodechar{ₐ}{\ifmmode_{a}\else\textsubscript{\ensuremath{a}}\fi}
\newunicodechar{ᵢ}{\ifmmode_{i}\else\textsubscript{\ensuremath{i}}\fi}
\newunicodechar{ᵣ}{\ifmmode_{r}\else\textsubscript{\ensuremath{r}}\fi}
\newunicodechar{ₖ}{\ifmmode_{k}\else\textsubscript{\ensuremath{k}}\fi}
\newunicodechar{ᵦ}{\ifmmode_{\beta}\else\textsubscript{\ensuremath{\beta}}\fi}
\newunicodechar{ₓ}{\ifmmode_{x}\else\textsubscript{\ensuremath{x}}\fi}
\newfontfamily\popdisplay{ArchivoBlack-Regular.ttf}[Path=../../../fonts/]
\newfontfamily\poplogo{SpaceGrotesk-Bold.ttf}[Path=../../../fonts/]
\newfontfamily\popsemi{PlusJakartaSans-SemiBold.ttf}[Path=../../../fonts/]
\newfontfamily\popxbold{PlusJakartaSans-ExtraBold.ttf}[Path=../../../fonts/]
\newfontfamily\popxboldls{PlusJakartaSans-ExtraBold.ttf}[Path=../../../fonts/,LetterSpace=3.0]

\setlist[enumerate]{leftmargin=1.7em,itemsep=5pt,topsep=4pt}
\setlist[itemize]{leftmargin=1.7em,itemsep=4pt,topsep=4pt,label={\color{poporange}\textbullet}}
\setlength{\parindent}{0pt}
\setlength{\parskip}{5pt}
\renewcommand{\arraystretch}{1.25}

% pgfplots : style Pop par défaut
\pgfplotsset{
  every axis/.append style={axis line style={popink,line width=1pt},tick style={popink},
    grid style={popline,line width=0.5pt},label style={font=\small},tick label style={font=\footnotesize}},
  popcurve/.style={poporange,line width=1.8pt,no markers,smooth},
}
% Même style disponible dans un \draw TikZ simple (hors pgfplots)
\tikzset{popcurve/.style={poporange,line width=1.8pt}}
\tikzset{popgrid/.style={popline,very thin}}
\colorlet{popcurve}{poporange} % le nom du style est parfois utilisé comme couleur

% ============================================================
% Pill / squiggle
% ============================================================
\newcommand{\poppill}[3]{%
  \tikz[baseline=-0.55ex]\node[draw=popink,line width=1.2pt,rounded corners=2.6pt,%
    fill=#2,inner xsep=6.5pt,inner ysep=3pt]{%
    \popxboldls\fontsize{7}{7}\selectfont\color{#3}\MakeUppercase{#1}};%
}
\newcommand{\popsquiggle}[1]{%
  \tikz[baseline=-0.4ex]{
    \draw[#1,line width=2.6pt,line cap=round]
      plot [smooth] coordinates {(0,0.09) (0.34,0.01) (0.68,0.21) (1.02,0.01) (1.36,0.10)};
  }%
}
% Mot-clé surligné (marqueur orange sous le texte)
\newcommand{\cle}[1]{{\popxbold\color{popdark}#1}}

% ============================================================
% En-tête / pied
% ============================================================
\pagestyle{fancy}
\fancyhf{}
\renewcommand{\headrulewidth}{0pt}
\renewcommand{\footrulewidth}{0pt}
\lhead{\raisebox{-0.15ex}{\poplogo\fontsize{13}{13}\selectfont\color{popink}OnBuch\color{poporange}+}}
\rhead{\poppill{\DOCMATIERE\ \textbullet\ \DOCNIVEAU\ \textbullet\ Chapter \DOCLECON}{white}{popink}}
\lfoot{\popxbold\fontsize{7.6}{7.6}\selectfont\color{popmuted}\MakeUppercase{Signed course OnBuch+}\ \textbullet\ {\popsemi\DOCTITRE}}
\rfoot{\tikz[baseline=-0.6ex]\node[rounded corners=8.5pt,fill=popink,minimum height=17pt,minimum width=17pt,inner xsep=5pt,inner sep=0]{\popxbold\fontsize{8}{8}\selectfont\color{popcream}\thepage\,/\,\pageref*{LastPage}};}

% ============================================================
% Titres
% ============================================================
\newcommand{\chaptitre}[2]{%
  \begin{tcolorbox}[enhanced,colback=poporange,colframe=popink,boxrule=2.6pt,arc=11pt,
      drop shadow=popink,left=15pt,right=15pt,top=12pt,bottom=11pt,before skip=3pt,after skip=18pt]
    \poppill{#2}{popink}{popcream}\par\vspace{9pt}
    {\raggedright\hyphenpenalty=10000\exhyphenpenalty=10000\popdisplay\fontsize{22}{24}\selectfont\color{popcream}\MakeUppercase{#1}\par}
  \end{tcolorbox}
}
% Section numérotée : pastille orange avec le numéro + titre Archivo
\newcounter{popsec}
\newcommand{\coursec}[1]{%
  \stepcounter{popsec}\setcounter{popsub}{0}%
  \par\vspace{16pt}\noindent
  \begin{minipage}{\linewidth}
  \tikz[baseline=-0.7ex]\node[draw=popink,line width=1.6pt,fill=poporange,rounded corners=4pt,minimum size=22pt,inner sep=0]{\popdisplay\fontsize{12}{12}\selectfont\color{popcream}\thepopsec};%
  \hspace{8pt}{\popdisplay\fontsize{15}{17}\selectfont\color{popink}\MakeUppercase{#1}}\par
  \vspace{4pt}\noindent{\color{poporange}\rule{36pt}{3.6pt}}
  \end{minipage}\par\nopagebreak\vspace{8pt}
}
\newcounter{popsub}
\newcommand{\courssub}[1]{%
  \stepcounter{popsub}%
  \par\vspace{9pt}\noindent{\popxbold\fontsize{11.5}{13}\selectfont\color{popink}{\color{poporange}\thepopsec.\thepopsub}\hspace{6pt}#1}\par\nopagebreak\vspace{4pt}
}

% ============================================================
% Boîtes pédagogiques — même recette : fond blanc, bordure épaisse colorée,
% ombre dure encre, badge plein. Titre optionnel [#1].
% ============================================================
\tcbset{popbase/.style={breakable,enhanced,colback=white,boxrule=2.3pt,arc=9pt,
  shadow={3pt}{-3pt}{0pt}{popink},left=14pt,right=14pt,top=11pt,bottom=11pt,before skip=11pt,after skip=15pt,
  fontupper=\popsemi\fontsize{10.6}{14.5}\selectfont\color{popink},
  fontlower=\popsemi\fontsize{10.6}{14.5}\selectfont\color{popink}}}
% Coche dessinée (le glyphe ✓ n'existe pas dans les polices du document)
\newcommand{\popcheck}[1]{\tikz[baseline=-0.5ex]\draw[#1,line width=1.6pt,line cap=round,line join=round](-0.13,0.0)--(-0.04,-0.09)--(0.13,0.1);}
\providecommand{\coche}{\popcheck{popgreen}}
\providecommand{\resultat}[1]{\par\smallskip\noindent\fcolorbox{popgreen}{popgreenL}{\parbox{\dimexpr\linewidth-2\fboxsep-2\fboxrule\relax}{\textbf{Result.} #1}}\par\smallskip}
\newcommand{\popboxhead}[3]{\poppill{#1}{#2}{white}\ifx\relax#3\relax\else\hspace{6pt}{\popxbold\fontsize{10.6}{12}\selectfont\color{popink}#3}\fi\par\vspace{7pt}}

\newtcolorbox{aretenir}[1][]{popbase,colframe=popgreen,before upper={\popboxhead{Key points}{popgreen}{#1}}}
\newtcolorbox{exemplebox}[1][]{popbase,colframe=poporange,before upper={\popboxhead{Exemple}{poporange}{#1}}}
\newtcolorbox{definition}[1][]{popbase,colframe=popblue,before upper={\popboxhead{Definition}{popblue}{#1}}}
\newtcolorbox{propriete}[1][]{popbase,colframe=poppurple,before upper={\popboxhead{Property}{poppurple}{#1}}}
\newtcolorbox{methode}[1][]{popbase,colframe=popgold,before upper={\popboxhead{Method}{popgold}{#1}}}
\newtcolorbox{attention}[1][]{popbase,colframe=poppink,before upper={\popboxhead{Warning}{poppink}{#1}}}
\newtcolorbox{experience}[1][]{popbase,colframe=popblue,colback=popblueL,before upper={\popboxhead{Experiment}{popblue}{#1}}}
% « Le savais-tu ? » : carte violette pleine (contexte, culture, Cameroun)
\newtcolorbox{savaistu}[1][]{popbase,colback=poppurple,colframe=popink,
  fontupper=\popsemi\fontsize{10.4}{14.2}\selectfont\color{white},
  before upper={\poppill{Did you know?}{white}{poppurple}\ifx\relax#1\relax\else\hspace{6pt}{\popxbold\color{white}#1}\fi\par\vspace{7pt}}}
% Exercice résolu : énoncé en haut, \tcblower puis la solution sur fond crème
\newcounter{popexo}\newcounter{popexr}
\newtcolorbox{exoresolu}[1][]{popbase,colframe=poporange,colbacklower=popcream,
  segmentation style={popink,line width=1.2pt,dashed},
  before upper={\stepcounter{popexr}\popboxhead{Worked example \thepopexr}{poporange}{#1}},
  before lower={\poppill{Solution}{popink}{popcream}\par\vspace{6pt}}}
% Exercice (non résolu en place) — la correction est dans la partie Corrigés
\newtcolorbox{exercice}[1][]{popbase,colframe=popink,
  before upper={\stepcounter{popexo}\popboxhead{Exercise \thepopexo}{popink}{#1}}}
% Corrigé
\newtcolorbox{corrige}[1][]{popbase,colframe=popgreen,colback=popgreenL,
  before upper={\popboxhead{Solution}{popgreen}{#1}}}
% Objectifs / prérequis / bilan
\newtcolorbox{objectifs}{popbase,colframe=popink,colback=poporangeL,
  before upper={\popboxhead{Chapter objectives}{popink}{}}}
\newtcolorbox{prerequis}{popbase,colframe=popink,colback=popblueL,
  before upper={\popboxhead{Prerequisites}{popblue}{}}}
\newtcolorbox{bilan}{popbase,colframe=popink,colback=white,boxrule=2.8pt,
  before upper={\popboxhead{Fiche bilan}{poporange}{L'essentiel à connaître pour l'examen}}}
% Figure encadrée avec légende
\newtcolorbox{popfigure}[1][]{enhanced,breakable=false,colback=white,colframe=popline,boxrule=1.4pt,arc=8pt,
  left=10pt,right=10pt,top=10pt,bottom=8pt,before skip=10pt,after skip=12pt,halign=center,
  lower separated=false,halign lower=center,
  fontlower=\popsemi\fontsize{9}{11.5}\selectfont\color{popmuted},
  #1}
\newcommand{\legende}[1]{\par\vspace{6pt}{\popsemi\fontsize{9}{11.5}\selectfont\color{popmuted}#1}}

% ============================================================
% Signature OnBuch+
% ============================================================
\newcommand{\poplogoinline}[1]{{\poplogo\fontsize{#1}{#1}\selectfont\color{popink}OnBuch\color{poporange}+}}
\newcommand{\signatureonbuch}{%
  \begin{tcolorbox}[enhanced,colback=popink,colframe=popink,boxrule=2.2pt,arc=10pt,
      drop shadow=poporange,left=14pt,right=14pt,top=10pt,bottom=10pt,before skip=0pt,after skip=0pt]
    \tikz[baseline=-0.7ex]\node[circle,fill=popgreen,draw=popcream,line width=1.3pt,minimum size=22pt,inner sep=0]{\popcheck{white}};%
    \hspace{9pt}\begin{minipage}[c]{0.62\linewidth}
      {\popxbold\fontsize{12}{13}\selectfont\color{popcream}Signed\ \textbullet\ {\poplogo OnBuch\color{poporange}+}}\par\vspace{2pt}
      {\raggedright\popsemi\fontsize{8}{10.5}\selectfont\color{popline}\MakeUppercase{OnBuch+ teaching team}\\\MakeUppercase{Follows the official MINESEC syllabus}\par}
    \end{minipage}\hfill
    {\poplogo\fontsize{22}{22}\selectfont\color{popcream}OnBuch\color{poporange}+}
  \end{tcolorbox}%
}

% ============================================================
% Page de garde
% #1 titre · #2 matière · #3 niveau · #4 module · #5 n° leçon · #6 durée ·
% #7 description / objectifs (texte ou itemize)
% ============================================================
\newcommand{\pagegarde}[7]{%
  \thispagestyle{empty}
  \newgeometry{a4paper,margin=1.7cm,top=1.6cm,bottom=1.4cm}%
  \begin{tikzpicture}[remember picture,overlay]
    \fill[poporange!18] ([xshift=-3.2cm,yshift=-3.6cm]current page.north east) circle (4.6cm);
    \fill[poppurple!14] ([xshift=2.4cm,yshift=7.6cm]current page.south west) circle (3.8cm);
  \end{tikzpicture}%
  {\poplogo\fontsize{30}{30}\selectfont\color{popink}OnBuch\color{poporange}+}\hfill
  \raisebox{0.6ex}{\poppill{Signed course \textbullet\ Premium}{popink}{popcream}}\par
  \vspace{1pt}\popsquiggle{poppurple}\par
  \vspace{8pt}
  \begin{tcolorbox}[enhanced,colback=poporange,colframe=popink,boxrule=2.8pt,arc=13pt,
      drop shadow=popink,left=18pt,right=18pt,top=12pt,bottom=13pt,after skip=10pt]
    \poppill{#2 \textbullet\ #3}{popink}{popcream}\hspace{5pt}\poppill{#4}{white}{popink}\par\vspace{8pt}
    {\popdisplay\fontsize{14}{14}\selectfont\color{popink}CHAPTER #5}\par\vspace{4pt}
    {\raggedright\hyphenpenalty=10000\exhyphenpenalty=10000\popdisplay\fontsize{25}{27}\selectfont\color{popcream}\MakeUppercase{#1}\par}
    \vspace{8pt}
    {\popxbold\fontsize{9.5}{9.5}\selectfont\color{popink}\MakeUppercase{Suggested time: #6}}
  \end{tcolorbox}
  \begin{tcolorbox}[enhanced,colback=white,colframe=popink,boxrule=2.2pt,arc=10pt,
      drop shadow=popink,left=16pt,right=16pt,top=10pt,bottom=10pt,after skip=8pt]
    \poppill{In this chapter}{poppurple}{white}\par\vspace{5pt}
    \popsemi\fontsize{9.3}{12.6}\selectfont\color{popink}#7
  \end{tcolorbox}
  \noindent\begin{minipage}[t]{0.487\linewidth}
    \begin{tcolorbox}[enhanced,colback=popgreen,colframe=popink,boxrule=2.2pt,arc=10pt,
        drop shadow=popink,left=11pt,right=11pt,top=10pt,bottom=10pt,height=3.1cm,valign=center]
      \tcbox[colback=white,colframe=popink,boxrule=1.3pt,arc=5pt,boxsep=0pt,nobeforeafter,
        left=3pt,right=3pt,top=3pt,bottom=3pt,valign=center]{\qrcode[height=1.9cm]{https://play.google.com/store/apps/details?id=com.luvvix.onbuch&hl=en}}%
      \hspace{8pt}\begin{minipage}[c]{0.5\linewidth}
        {\popdisplay\fontsize{11}{12.5}\selectfont\color{white}\MakeUppercase{Download OnBuch}}\par\vspace{4pt}
        {\raggedright\popsemi\fontsize{8.4}{11}\selectfont\color{white}Free \textbullet\ Google Play\\AI tutor, past papers, results\par}
      \end{minipage}
    \end{tcolorbox}
  \end{minipage}\hfill
  \begin{minipage}[t]{0.487\linewidth}
    \begin{tcolorbox}[enhanced,colback=poppurple,colframe=popink,boxrule=2.2pt,arc=10pt,
        drop shadow=popink,left=11pt,right=11pt,top=10pt,bottom=10pt,height=3.1cm,valign=center]
      \tikz[baseline=-0.7ex]\node[circle,fill=white,draw=popink,line width=1.3pt,minimum size=1.6cm,inner sep=0]{\popdisplay\fontsize{24}{24}\selectfont\color{poppurple}+};%
      \hspace{8pt}\begin{minipage}[c]{0.6\linewidth}
        {\popdisplay\fontsize{11}{12.5}\selectfont\color{white}\MakeUppercase{Go OnBuch+}}\par\vspace{4pt}
        {\raggedright\popsemi\fontsize{8.4}{11}\selectfont\color{white}All signed courses, unlimited AI tutor, PDF sheets — OnBuch+ tab in the app\par}
      \end{minipage}
    \end{tcolorbox}
  \end{minipage}
  \vspace{8pt}
  \popcontenu
  \vfill
  \signatureonbuch
  \restoregeometry
  \newpage
}

% Rangée « Ce cours contient » (page de garde)
\newcommand{\poptile}[3]{% #1 couleur #2 gros texte #3 légende
  \begin{tcolorbox}[enhanced,colback=#1,colframe=popink,boxrule=2pt,arc=9pt,shadow={2.5pt}{-2.5pt}{0pt}{popink},
      left=6pt,right=6pt,top=9pt,bottom=9pt,halign=center,valign=center,height=1.75cm,width=\linewidth,nobeforeafter]
    {\popdisplay\fontsize{12.5}{13}\selectfont\color{white}\MakeUppercase{#2}}\par\vspace{3pt}
    {\centering\popsemi\fontsize{7.8}{9.5}\selectfont\color{white}#3\par}
  \end{tcolorbox}}
\newcommand{\popcontenu}{%
  \par\noindent{\popxboldls\fontsize{7.5}{7.5}\selectfont\color{popmuted}THIS SIGNED COURSE CONTAINS}\par\vspace{7pt}
  \noindent\begin{minipage}[t]{0.235\linewidth}\poptile{popblue}{Course}{complete, illustrated, step by step}\end{minipage}\hfill
  \begin{minipage}[t]{0.235\linewidth}\poptile{poporange}{Methods}{and worked examples}\end{minipage}\hfill
  \begin{minipage}[t]{0.235\linewidth}\poptile{poppink}{Exercises}{exam style + solutions}\end{minipage}\hfill
  \begin{minipage}[t]{0.235\linewidth}\poptile{popgreen}{Summary sheet}{the essentials to remember}\end{minipage}\par}

% Sommaire
\newtcolorbox{sommaire}{enhanced,colback=white,colframe=popink,boxrule=2.2pt,arc=10pt,
  drop shadow=popink,left=18pt,right=18pt,top=13pt,bottom=13pt,before skip=6pt,after skip=20pt,
  before upper={\poppill{Contents}{poppurple}{white}\par\vspace{9pt}\popsemi\fontsize{10.6}{14.5}\selectfont\color{popink}}}

% Signature de fin de cours — poussée tout en bas de la dernière page
\newcommand{\findecours}{%
  \par\vfill\nopagebreak
  \begin{center}\popsquiggle{poporange}\end{center}\vspace{4pt}
  \signatureonbuch
}
\AtBeginDocument{\def\llbracket{\mathopen{[\mkern-3.5mu[}}\def\rrbracket{\mathclose{]\mkern-3.5mu]}}} % intervalles d'entiers (glyphe ⟦ absent de la police)
% Glyphes absents de Fira Math : redéfinis à partir de glyphes disponibles
\AtBeginDocument{%
  \def\setminus{\mathbin{\backslash}}\let\smallsetminus\setminus
  \def\vdots{\mathord{\vbox{\baselineskip3.2pt\lineskiplimit0pt\kern4pt\hbox{.}\hbox{.}\hbox{.}}}}%
  \def\ddots{\mathinner{\mkern1mu\raise6.5pt\hbox{.}\mkern2mu\raise3.6pt\hbox{.}\mkern2mu\raise0.7pt\hbox{.}\mkern1mu}}%
  \def\triangle{\mathord{\scalebox{0.9}{$\symup{\Delta}$}}}%
  \def\nabla{\mathord{\raisebox{\depth}{\scalebox{1}[-1]{$\symup{\Delta}$}}}}%
  \def\checkmark{\popcheck{popdark}}%
  \def\star{\mathbin{\ast}}\def\bigstar{\mathord{\ast}}%
  \def\diamond{\mathbin{\scalebox{0.6}{$\Diamond$}}}%
  \def\bigcup{\mathop{\mathchoice{\raisebox{-0.3em}{\scalebox{1.7}{$\cup$}}}{\scalebox{1.25}{$\cup$}}{\cup}{\cup}}\displaylimits}%
  \def\bigcap{\mathop{\mathchoice{\raisebox{-0.3em}{\scalebox{1.7}{$\cap$}}}{\scalebox{1.25}{$\cap$}}{\cap}{\cap}}\displaylimits}%
  \def\lesseqqgtr{\mathrel{\lessgtr}}\def\bigoplus{\mathop{\oplus}\displaylimits}%
}
\providecommand{\ssi}{if and only if} % abréviation fréquente
\AtBeginDocument{\providecommand{\wideparen}{\overparen}} % arc de cercle

%% FIGFIX-BEGIN v3 — correctifs globaux des figures (docs/onbuch-generation/tools/figfix)
\usepackage{adjustbox}\usepackage{etoolbox}
% toute figure TikZ/pgfplots plus large que la ligne est réduite à la largeur disponible
\BeforeBeginEnvironment{tikzpicture}{\begin{adjustbox}{max width=\linewidth}}
\AfterEndEnvironment{tikzpicture}{\end{adjustbox}}
% légendes pgfplots lisibles par-dessus les courbes
\pgfplotsset{every axis/.append style={legend style={fill=white,fill opacity=0.92,text opacity=1,draw=black!15,inner sep=3pt}}}
% étiquettes d'arêtes (syntaxe edge["..."]) avec fond blanc
\tikzset{every edge quotes/.append style={fill=white,inner sep=1.2pt}}
% bulles de cartes mentales : texte à la ligne dans la bulle, bulle un peu plus grande
\tikzset{every concept/.append style={text width=2.0cm,align=center,minimum size=2.3cm,inner sep=2pt}}
%% FIGFIX-END
\begin{document}

\pagegarde{PLane Transformations}{Further Mathematics}{Upper Sixth}{Upper Sixth — Further mathematics}{3}{18 periods}{This chapter studies the composition of plane transformations — reflections, translations, rotations, glide reflections and homotheties — leading to the classification of plane isometries and similarities. Students learn to identify, construct and combine these transformations geometrically, analytically and with complex numbers, in preparation for the GCE Advanced Level examination.
\vspace{4pt}
\begin{multicols}{2}\small
\begin{itemize}
\item By the end of this chapter I can determine the nature and characteristic elements of the composite of two orthogonal symmetries (parallel or intersecting axes).
\item By the end of this chapter I can compose two translations and a translation with a reflection, and recognise a glide reflection.
\item By the end of this chapter I can compose two rotations of the same or different centres and identify the resulting transformation.
\item By the end of this chapter I can define a plane isometry and classify any isometry as a translation, rotation, reflection or glide reflection.
\item By the end of this chapter I can construct the image of a point or figure under a homothety and state its basic and advanced properties.
\item By the end of this chapter I can distinguish direct from indirect homothety and use the sign of the ratio correctly.
\end{itemize}\end{multicols}}

\begin{sommaire}
\begin{multicols}{2}
\begin{enumerate}
\item Composing Reflections and Translations
\item Composing Rotations; Introduction to Isometries
\item Classification of Plane Isometries and Glide Reflections
\item Homothety: Definition, Construction and Basic Properties
\item Analytical and Complex Representation of Homothety
\item Composition of Homotheties and Advanced Properties
\item Plane Similarities: Concept and Characteristic Elements
\item Analytical Study of Similarities and Applications
\item Integration activity
\item Methods and worked examples
\item Exercises
\item Solutions to the exercises
\item Summary sheet
\end{enumerate}
\end{multicols}
\end{sommaire}

\begin{objectifs}
\begin{itemize}
\item By the end of this chapter I can determine the nature and characteristic elements of the composite of two orthogonal symmetries (parallel or intersecting axes).
\item By the end of this chapter I can compose two translations and a translation with a reflection, and recognise a glide reflection.
\item By the end of this chapter I can compose two rotations of the same or different centres and identify the resulting transformation.
\item By the end of this chapter I can define a plane isometry and classify any isometry as a translation, rotation, reflection or glide reflection.
\item By the end of this chapter I can construct the image of a point or figure under a homothety and state its basic and advanced properties.
\item By the end of this chapter I can distinguish direct from indirect homothety and use the sign of the ratio correctly.
\item By the end of this chapter I can represent a homothety analytically and by a complex number or vector expression.
\item By the end of this chapter I can compose two homotheties with the same or distinct centres and find the centre and ratio of the result.
\item By the end of this chapter I can define a plane similarity, give its analytical form, and solve problems on triangles, circles and spiral similarities.
\end{itemize}
\end{objectifs}

\begin{prerequis}
\begin{itemize}
\item Translations, reflections (orthogonal symmetries) and rotations as studied in Lower Sixth: definitions and basic properties
\item Vectors in the plane: addition, scalar multiplication, coordinates and collinearity
\item Complex numbers: modulus, argument, and the geometric meaning of z' = az + b
\item Cartesian coordinates, equations of lines and circles in the plane
\item Congruent and similar triangles, Thales' theorem and ratios of lengths
\end{itemize}
\end{prerequis}

\newpage

\coursec{Composing Reflections and Translations}

When a tailor in Bamenda folds a wax-print fabric along two parallel lines before cutting, he is unknowingly composing two reflections — and the result is exactly a translation. When a mechanic rotates a wheel nut with a spanner, or when a photographer in Douala enlarges a passport photo without distorting it, plane transformations are silently at work. What single transformation is hidden behind the composition of two symmetries, two rotations, or two enlargements? This chapter gives you the algebraic and geometric tools to answer such questions — and they appear every year in the GCE Advanced Level Further Mathematics paper.

In this first section we master the two simplest building blocks — \cle{reflections} and \cle{translations} — and we discover the beautiful fact that composing two reflections produces either a translation or a rotation, while composing a reflection with a translation produces a new creature: the \cle{glide reflection}.

\courssub{Recall: Orthogonal Symmetry (Reflection) Across a Line}

\begin{definition}[Orthogonal symmetry (reflection) across a line]
Let $d$ be a line in the plane. The \cle{orthogonal symmetry} (or \cle{reflection}) across $d$, denoted $s_d$, is the transformation of the plane defined by:
\begin{itemize}
\item if $M \in d$, then $s_d(M) = M$;
\item if $M \notin d$, then $s_d(M) = M'$, where $d$ is the perpendicular bisector of the segment $[MM']$.
\end{itemize}
The line $d$ is called the \cle{axis} of the reflection.
\end{definition}

\begin{aretenir}[Essential facts about a reflection]
\begin{itemize}
\item The \textbf{invariant (fixed) points} of $s_d$ are exactly the points of the axis $d$.
\item $s_d$ is an \textbf{involution}: $s_d \circ s_d = \mathrm{Id}$ (reflecting twice brings every point back).
\item $s_d$ \textbf{preserves distances} (it is an isometry) but \textbf{reverses orientation}: a clockwise triangle becomes anticlockwise.
\item To construct the image of $M$: drop the perpendicular from $M$ to $d$, meeting $d$ at $H$, then mark $M'$ on the same line with $HM' = HM$ on the other side.
\end{itemize}
\end{aretenir}

\courssub{Composition of Two Reflections with Parallel Axes}

Let $d_1$ and $d_2$ be two \textbf{parallel} lines. What is the nature of $s_{d_2} \circ s_{d_1}$? Take a point $M$, reflect it in $d_1$ to get $M_1$, then reflect $M_1$ in $d_2$ to get $M'$. Whatever the position of $M$, the segment $[MM']$ is perpendicular to both axes, and its length is always \emph{twice the distance between the axes}. This suggests a translation.

\begin{propriete}[Two reflections with parallel axes give a translation — proof examinable]
Let $d_1 \parallel d_2$, and let $\vec{u}$ be the vector going from $d_1$ to $d_2$ along a common perpendicular (so $\|\vec{u}\|$ is the distance between the axes). Then
\[ s_{d_2} \circ s_{d_1} = t_{2\vec{u}}, \]
the translation of vector $2\vec{u}$.
\end{propriete}

\textbf{Proof.} Let $M$ be any point, $M_1 = s_{d_1}(M)$ and $M' = s_{d_2}(M_1)$. Let $H_1$ and $H_2$ be the orthogonal projections of $M$ onto $d_1$ and $d_2$ respectively; since $d_1 \parallel d_2$, the points $M, M_1, M', H_1, H_2$ all lie on one line perpendicular to the axes. Using directed distances along this line:
\[ \overline{MM'} = \overline{MM_1} + \overline{M_1M'} = 2\overline{H_1M_1} + 2\overline{M_1H_2} = 2\overline{H_1H_2}. \]
But $\overrightarrow{H_1H_2} = \vec{u}$, independent of $M$. Hence $\overrightarrow{MM'} = 2\vec{u}$ for every point $M$, which is exactly the definition of the translation $t_{2\vec{u}}$. $\blacksquare$

\begin{popfigure}
\begin{center}
\begin{tikzpicture}[scale=1.0, >=stealth]
\draw[popblue, thick] (0,-2.2) -- (0,2.2) node[above] {$d_1$};
\draw[poporange, thick] (2.2,-2.2) -- (2.2,2.2) node[above] {$d_2$};
\fill[popink] (-1.3,0.9) circle (2pt) node[left] {$M$};
\fill[popblue] (1.3,0.9) circle (2pt) node[above right] {$M_1$};
\fill[poporange] (3.1,0.9) circle (2pt) node[right] {$M'$};
\draw[dashed, gray] (-1.3,0.9) -- (3.1,0.9);
\draw[popdark, very thick, ->] (-1.3,-1.4) -- (3.1,-1.4) node[midway, below] {$2\vec{u}$};
\draw[popgreen, thick, ->] (0,-1.9) -- (2.2,-1.9) node[midway, below] {$\vec{u}$};
\draw[<->, popmuted] (0,1.9) -- (2.2,1.9) node[midway, above] {gap};
\end{tikzpicture}
\end{center}
\legende{$M_1 = s_{d_1}(M)$, $M' = s_{d_2}(M_1)$: the composite is the translation of vector $2\vec{u}$, twice the perpendicular gap between the axes.}
\end{popfigure}

\begin{attention}[Order matters!]
$s_{d_2} \circ s_{d_1} = t_{2\vec{u}}$ where $\vec{u}$ goes from $d_1$ to $d_2$. Reversing the order gives the \textbf{opposite} translation:
\[ s_{d_1} \circ s_{d_2} = t_{-2\vec{u}}. \]
Composition of transformations is \emph{not} commutative in general. Always read $s_{d_2} \circ s_{d_1}$ as: ``first $s_{d_1}$, then $s_{d_2}$''.
\end{attention}

\courssub{Composition of Two Reflections with Intersecting Axes}

Now suppose $d_1$ and $d_2$ meet at a point $O$, making an oriented angle $\theta = (\widehat{d_1, d_2})$. Reflecting a point $M$ successively in $d_1$ then $d_2$ keeps its distance to $O$ unchanged (each reflection does), but swings it around $O$. The total swing is exactly $2\theta$.

\begin{propriete}[Two reflections with intersecting axes give a rotation — proof examinable]
If $d_1 \cap d_2 = \{O\}$ and $\theta$ is the oriented angle from $d_1$ to $d_2$, then
\[ s_{d_2} \circ s_{d_1} = r(O, 2\theta), \]
the rotation of centre $O$ and angle $2\theta$.
\end{propriete}

\textbf{Proof.} Let $M \neq O$, $M_1 = s_{d_1}(M)$, $M' = s_{d_2}(M_1)$. Since reflections preserve distances and fix $O$, we have $OM = OM_1 = OM'$, so $M'$ lies on the circle of centre $O$ through $M$. For the angle, use oriented angles of lines. Reflection in a line $d$ sends any ray $[OM)$ to a ray symmetric with respect to $d$, so:
\[ (\widehat{\overrightarrow{OM}, \overrightarrow{OM'}}) = (\widehat{\overrightarrow{OM}, d_1}) + (\widehat{d_1, d_2}) + (\widehat{d_2, \overrightarrow{OM'}}). \]
By symmetry, $(\widehat{d_2, \overrightarrow{OM'}}) = (\widehat{\overrightarrow{OM_1}, d_2})$ and $(\widehat{\overrightarrow{OM}, d_1}) = (\widehat{d_1, \overrightarrow{OM_1}})$. Hence
\[ (\widehat{\overrightarrow{OM}, \overrightarrow{OM'}}) = (\widehat{d_1, \overrightarrow{OM_1}}) + (\widehat{\overrightarrow{OM_1}, d_2}) + (\widehat{d_1, d_2}) = 2(\widehat{d_1, d_2}) = 2\theta. \]
Also $O$ is fixed, and the map preserves distances, so $s_{d_2} \circ s_{d_1}$ is the rotation of centre $O$ and angle $2\theta$. $\blacksquare$

\begin{popfigure}
\begin{center}
\begin{tikzpicture}[scale=1.1, >=stealth]
\coordinate (O) at (0,0);
\draw[popblue, thick] (-2.6,-0.9) -- (2.6,0.9) node[right] {$d_1$};
\draw[poporange, thick] (-2.3,-1.7) -- (2.3,1.7) node[right] {$d_2$};
\fill (O) circle (2pt) node[below left] {$O$};
\fill[popink] (2.1,-0.5) circle (2pt) node[below] {$M$};
\fill[popblue] (1.9,1.0) circle (2pt) node[above] {$M_1$};
\fill[poporange] (0.5,2.1) circle (2pt) node[above] {$M'$};
\draw[popdark, thick, ->] (1.4,0.47) arc (18.4:36.8:1.5) node[midway, right] {$\theta$};
\draw[popgreen, very thick, ->] (2.1,-0.5) arc (-13.4:76.6:2.16) node[midway, right, xshift=2pt] {$2\theta$};
\draw[dashed, gray] (O) -- (2.1,-0.5);
\draw[dashed, gray] (O) -- (0.5,2.1);
\end{tikzpicture}
\end{center}
\legende{Axes $d_1, d_2$ meeting at $O$ with angle $\theta$: the composite $s_{d_2}\circ s_{d_1}$ is the rotation of centre $O$ and angle $2\theta$.}
\end{popfigure}

\begin{attention}[The factor 2 — the classic trap]
A very common mistake at the GCE is to forget the factor $2$:
\begin{itemize}
\item parallel axes: the translation vector is \textbf{twice} the gap between the axes, not the gap itself;
\item intersecting axes: the rotation angle is \textbf{twice} the angle between the axes, not the angle itself.
\end{itemize}
Also remember: reversing the order of composition reverses the translation vector (resp.\ the rotation angle).
\end{attention}

\courssub{Composition of Two Translations}

\begin{propriete}[Composite of two translations]
Let $t_{\vec{u}}$ and $t_{\vec{v}}$ be translations. Then
\[ t_{\vec{v}} \circ t_{\vec{u}} = t_{\vec{u} + \vec{v}}. \]
In particular, translations \textbf{commute}: $t_{\vec{v}} \circ t_{\vec{u}} = t_{\vec{u}} \circ t_{\vec{v}}$.
\end{propriete}

\textbf{Proof.} For any point $M$, write $M_1 = t_{\vec{u}}(M)$ and $M' = t_{\vec{v}}(M_1)$. Then
\[ \overrightarrow{MM'} = \overrightarrow{MM_1} + \overrightarrow{M_1M'} = \vec{u} + \vec{v}, \]
so the composite is $t_{\vec{u}+\vec{v}}$. Since vector addition is commutative, $\vec{u}+\vec{v} = \vec{v}+\vec{u}$, and the two orders coincide. $\blacksquare$

\begin{exemplebox}[Composing translations]
A delivery rider in Yaoundé moves a parcel by $\vec{u} = (3, -1)$ km, then by $\vec{v} = (-2, 4)$ km. The total displacement is $\vec{u} + \vec{v} = (1, 3)$ km: a single translation $t_{(1,3)}$ achieves the same result, in either order.
\end{exemplebox}

\courssub{Composition of a Translation and a Reflection: the Glide Reflection}

What happens if we compose a reflection $s_d$ with a translation $t_{\vec{v}}$? Split $\vec{v} = \vec{v}_{\parallel} + \vec{v}_{\perp}$ into components parallel and perpendicular to $d$. The perpendicular part can be absorbed by shifting the axis $d$ (using the theorem on parallel axes above), and only the parallel part survives. The result is a new transformation.

\begin{definition}[Glide reflection]
A \cle{glide reflection} is the composite of a reflection across a line $d$ and a translation of vector $\vec{v}$ \textbf{parallel to} $d$ (with $\vec{v} \neq \vec{0}$):
\[ g = t_{\vec{v}} \circ s_d = s_d \circ t_{\vec{v}}. \]
The line $d$ is the \cle{axis} and $\vec{v}$ the \cle{vector} of the glide reflection. (The two orders coincide here because $\vec{v} \parallel d$; indeed, reflecting a vector parallel to the axis leaves it unchanged.)
\end{definition}

Think of footprints in the sand along a straight path: each step is the mirror image of the previous one across the line of the path, shifted forward by one pace. That is exactly a glide reflection.

\begin{popfigure}
\begin{center}
\begin{tikzpicture}[scale=0.9, >=stealth]
\draw[popblue, thick] (-0.5,0) -- (7.5,0) node[right] {axis $d$};
\draw[poporange, very thick, ->] (0.4,-0.7) -- (2.4,-0.7) node[midway, below] {$\vec{v}$};
\foreach \x/\s in {0.4/1, 2.4/-1, 4.4/1, 6.4/-1}{
\begin{scope}[shift={(\x,0)}, yscale=\s]
\draw[popink, thick] (0,0.25) ellipse (0.16 and 0.34);
\draw[popink, thick] (0.05,0.75) ellipse (0.13 and 0.2);
\end{scope}}
\end{tikzpicture}
\end{center}
\legende{Footprints along a line: each print is obtained from the previous one by reflection across $d$ followed by a translation of vector $\vec{v}$ parallel to $d$ — a glide reflection.}
\end{popfigure}

\begin{aretenir}[Recognition criteria for a glide reflection]
A plane transformation $g$ is a glide reflection (with $\vec{v} \neq \vec{0}$) if and only if:
\begin{itemize}
\item $g$ \textbf{preserves distances} (it is an isometry);
\item $g$ \textbf{reverses orientation} (it is an \emph{indirect} isometry);
\item $g$ has \textbf{no fixed point}.
\end{itemize}
Its axis is the unique invariant line, and $g \circ g = t_{2\vec{v}}$.
\end{aretenir}

\begin{methode}[Identifying a composite $s_{d_2} \circ s_{d_1}$]
\begin{enumerate}
\item \textbf{Draw the axes.} Are they parallel or intersecting?
\item \textbf{Parallel axes:} the composite is a translation; measure the perpendicular gap from $d_1$ to $d_2$ and double it (mind the direction: $d_1 \to d_2$).
\item \textbf{Intersecting axes:} the composite is a rotation; the centre is the intersection point, the angle is twice the oriented angle from $d_1$ to $d_2$.
\item \textbf{Check with one test point:} pick a convenient point $M$, construct $M_1 = s_{d_1}(M)$ then $M' = s_{d_2}(M_1)$, and verify that $\overrightarrow{MM'}$ (or the angle $\widehat{MOM'}$) matches your answer.
\end{enumerate}
\end{methode}

\begin{exoresolu}[Parallel axes]
The lines $d_1: x = 1$ and $d_2: x = 4$ are given in an orthonormal frame. Determine the nature and characteristic elements of $f = s_{d_2} \circ s_{d_1}$, and find the image of $A(2, 3)$.
\tcblower
\textbf{Solution.} The axes are parallel (both vertical). The vector from $d_1$ to $d_2$ along a common perpendicular is $\vec{u} = (3, 0)$. By the theorem, $f = t_{2\vec{u}} = t_{(6,0)}$.
Check with $A(2,3)$: $s_{d_1}(A) = A_1(0, 3)$ (since $A$ is $1$ unit right of $d_1$), then $s_{d_2}(A_1) = A'(8, 3)$ (since $A_1$ is $4$ units left of $d_2$). Indeed $\overrightarrow{AA'} = (6, 0) = 2\vec{u}$. $\blacksquare$
\end{exoresolu}

\begin{exoresolu}[Intersecting axes]
$d_1$ and $d_2$ pass through the origin, with $d_1$ the $x$-axis and $d_2$ making an angle of $30°$ with $d_1$. Determine $g = s_{d_2} \circ s_{d_1}$ and the image of $B(2, 0)$.
\tcblower
\textbf{Solution.} The axes intersect at $O$ with oriented angle $\theta = 30°$. Hence $g = r(O, 2\theta) = r(O, 60°)$. The image of $B(2,0)$ is
\[ B' = \bigl(2\cos 60°,\ 2\sin 60°\bigr) = \bigl(1,\ \sqrt{3}\bigr). \]
Check: $s_{d_1}(B) = B$ (since $B \in d_1$), and reflecting $B$ across $d_2$ (line at $30°$) sends the ray $[OB)$ to the ray at angle $2 \times 30° = 60°$, giving $B'(1, \sqrt{3})$. $\blacksquare$
\end{exoresolu}

\begin{exercice}[Practice]
\begin{enumerate}
\item $d_1: y = 0$ and $d_2: y = 2$. Give the nature and vector of $s_{d_2} \circ s_{d_1}$, then of $s_{d_1} \circ s_{d_2}$.
\item Two lines through $O$ make an angle of $45°$. Show that $s_{d_2} \circ s_{d_1}$ is a quarter turn, and deduce that $(s_{d_2} \circ s_{d_1})^2$ is a half-turn.
\item Let $g = t_{\vec{v}} \circ s_d$ with $d: y = 0$ and $\vec{v} = (3, 0)$. Find the image of $P(1, 2)$ and show that $g$ has no fixed point.
\end{enumerate}
\end{exercice}

\begin{corrige}[Exercise 1]
\begin{enumerate}
\item $\vec{u} = (0,2)$ from $d_1$ to $d_2$, so $s_{d_2} \circ s_{d_1} = t_{(0,4)}$ and $s_{d_1} \circ s_{d_2} = t_{(0,-4)}$.
\item $2\theta = 90°$: a quarter turn about $O$. Its square has angle $180°$: a half-turn about $O$.
\item $s_d(P) = (1, -2)$, then $g(P) = (4, -2)$. A fixed point would need $(x+3, -y) = (x, y)$, i.e.\ $3 = 0$: impossible. So $g$ has no fixed point — consistent with a glide reflection.
\end{enumerate}
\end{corrige}

\begin{savaistu}[Why this matters]
The theorems of this section are the key to the \emph{classification of isometries} (Section 3): every plane isometry is a translation, a rotation, a reflection or a glide reflection, because every isometry can be written as a composite of at most three reflections. Glide reflections also describe the symmetry of footprints, of woven raffia patterns, and of many crystal structures studied in materials science.
\end{savaistu}

\coursec{Composing Rotations; Introduction to Isometries}

In the previous section we composed reflections and translations, and we discovered a powerful principle: composing two reflections with parallel axes gives a translation, while composing two reflections with intersecting axes gives a rotation. We now turn this principle around and use it as a tool. Since every rotation can be written as a composite of two reflections, we can study the composition of two rotations by cleverly decomposing them into reflections whose axes we choose ourselves. This beautiful technique, sometimes called the \emph{reflection-decomposition trick}, will allow us to determine completely the nature of any composite of rotations — and it will lead us naturally to the general notion of a \cle{plane isometry}.

\courssub{Recall: Rotations and their Decomposition into Reflections}

\begin{definition}[Rotation]
Let $O$ be a point of the plane and $\theta$ a real number (an angle, measured in degrees or radians). The \cle{rotation} of centre $O$ and angle $\theta$, denoted $r(O,\theta)$, is the transformation of the plane that fixes $O$ and sends every point $M \neq O$ to the point $M'$ such that
\[ OM' = OM \qquad \text{and} \qquad \left(\overrightarrow{OM}, \overrightarrow{OM'}\right) = \theta \pmod{360^\circ}. \]
\end{definition}

The key structural fact, established in the previous section, is the following.

\begin{propriete}[A rotation is a composite of two reflections]
Let $(d_1)$ and $(d_2)$ be two lines intersecting at $O$, with $\big((d_1),(d_2)\big) = \dfrac{\theta}{2} \pmod{180^\circ}$. Then
\[ r(O,\theta) = s_{(d_2)} \circ s_{(d_1)}. \]
Moreover, for a given rotation $r(O,\theta)$, one of the two axes may be chosen \emph{freely} through $O$; the second axis is then forced: it is the image of the first by the rotation $r\!\left(O, \tfrac{\theta}{2}\right)$.
\end{propriete}

\begin{attention}[The order matters]
In $s_{(d_2)} \circ s_{(d_1)}$ we reflect first in $(d_1)$, then in $(d_2)$, and the resulting rotation goes from $(d_1)$ towards $(d_2)$. Swapping the axes reverses the angle: $s_{(d_1)} \circ s_{(d_2)} = r(O, -\theta)$.
\end{attention}

The freedom to choose one axis arbitrarily is exactly what makes the composition of rotations tractable, as we now see.

\courssub{Composition of Two Rotations with the Same Centre}

Let us start with the simplest situation. Take a point $M$ in the plane and rotate it about $O$ through an angle $\alpha$, obtaining $M_1$; then rotate $M_1$ about the \emph{same} centre $O$ through an angle $\beta$, obtaining $M'$. Geometrically, the point has simply turned through $\alpha$ and then through $\beta$ about the same pivot: the total effect is a turn through $\alpha + \beta$.

\begin{propriete}[Same centre]
For any point $O$ and any angles $\alpha, \beta$,
\[ r(O,\beta) \circ r(O,\alpha) = r(O, \alpha + \beta). \]
In particular, two rotations with the same centre commute: $r(O,\beta)\circ r(O,\alpha) = r(O,\alpha)\circ r(O,\beta)$.
\end{propriete}

\textbf{Proof (examinable).} Write $r(O,\alpha) = s_{(d_2)} \circ s_{(d_1)}$ where $(d_1), (d_2)$ pass through $O$ with $\big((d_1),(d_2)\big) = \tfrac{\alpha}{2}$, and write $r(O,\beta) = s_{(d_3)} \circ s_{(d_2)}$, choosing the \emph{same} line $(d_2)$ as first axis of the second rotation (this is allowed by the freedom of choice above). Then
\[ r(O,\beta) \circ r(O,\alpha) = s_{(d_3)} \circ s_{(d_2)} \circ s_{(d_2)} \circ s_{(d_1)} = s_{(d_3)} \circ s_{(d_1)}, \]
because $s_{(d_2)} \circ s_{(d_2)} = \mathrm{Id}$. The composite $s_{(d_3)} \circ s_{(d_1)}$ is a rotation about $O$ of angle $2\left(\tfrac{\alpha}{2} + \tfrac{\beta}{2}\right) = \alpha + \beta$. \hfill$\blacksquare$

\begin{exemplebox}[Same centre]
$r(O, 70^\circ) \circ r(O, 50^\circ) = r(O, 120^\circ)$, and $r(O, -30^\circ) \circ r(O, 90^\circ) = r(O, 60^\circ)$. Also $r(O, 180^\circ) \circ r(O, 180^\circ) = r(O, 360^\circ) = \mathrm{Id}$, which confirms that a half-turn is its own inverse.
\end{exemplebox}

\courssub{Composition of Two Rotations with Different Centres}

Now suppose the centres differ: $r_1 = r(O_1, \alpha)$ and $r_2 = r(O_2, \beta)$ with $O_1 \neq O_2$. What is $r_2 \circ r_1$? The pivot changes midway, so the answer is no longer obvious. The reflection-decomposition trick settles the question completely: we decompose each rotation into two reflections, and we choose the \emph{common axis} to be the line $(O_1O_2)$, so that the two middle reflections cancel.

Concretely:
\begin{itemize}
\item Write $r_1 = s_{(O_1O_2)} \circ s_{(d_1)}$, where $(d_1)$ is the line through $O_1$ such that $\big((d_1), (O_1O_2)\big) = \tfrac{\alpha}{2}$.
\item Write $r_2 = s_{(d_2)} \circ s_{(O_1O_2)}$, where $(d_2)$ is the line through $O_2$ such that $\big((O_1O_2), (d_2)\big) = \tfrac{\beta}{2}$.
\end{itemize}
Then
\[ r_2 \circ r_1 = s_{(d_2)} \circ s_{(O_1O_2)} \circ s_{(O_1O_2)} \circ s_{(d_1)} = s_{(d_2)} \circ s_{(d_1)}. \]

Two cases arise, according to whether $(d_1)$ and $(d_2)$ intersect or are parallel.

\begin{propriete}[Composite of two rotations with distinct centres — proof examinable]
Let $r_1 = r(O_1, \alpha)$ and $r_2 = r(O_2, \beta)$ with $O_1 \neq O_2$.
\begin{itemize}
\item If $\alpha + \beta \not\equiv 0 \pmod{360^\circ}$, the lines $(d_1)$ and $(d_2)$ constructed above meet at a point $\Omega$, and
\[ r_2 \circ r_1 = r(\Omega, \alpha + \beta). \]
The centre $\Omega$ is characterised by the half-angle conditions: the line $(\Omega O_1)$ makes angle $-\tfrac{\alpha}{2}$ with $(O_1O_2)$ at $O_1$, and the line $(\Omega O_2)$ makes angle $+\tfrac{\beta}{2}$ with $(O_2O_1)$ at $O_2$ (on the appropriate sides).
\item If $\alpha + \beta \equiv 0 \pmod{360^\circ}$, then $(d_1) \parallel (d_2)$ and $r_2 \circ r_1$ is a \cle{translation}.
\end{itemize}
\end{propriete}

\textbf{Justification of the two cases.} The angle between $(d_1)$ and $(d_2)$ equals $\tfrac{\alpha}{2} + \tfrac{\beta}{2} = \tfrac{\alpha+\beta}{2}$ (as directed angles of lines, modulo $180^\circ$). If $\alpha + \beta \not\equiv 0 \pmod{360^\circ}$, this angle is not a multiple of $180^\circ$, so $(d_1)$ and $(d_2)$ intersect at a point $\Omega$, and $s_{(d_2)} \circ s_{(d_1)}$ is the rotation of centre $\Omega$ and angle $2 \cdot \tfrac{\alpha+\beta}{2} = \alpha + \beta$. If $\alpha + \beta \equiv 0 \pmod{360^\circ}$, then $\tfrac{\alpha+\beta}{2}$ is a multiple of $180^\circ$, so $(d_1)$ and $(d_2)$ are parallel (and distinct, since $O_1 \neq O_2$), and the composite of two reflections with parallel distinct axes is a translation. \hfill$\blacksquare$

\begin{methode}[Finding the centre of $r(O_2,\beta) \circ r(O_1,\alpha)$]
\begin{enumerate}
\item Check that $\alpha + \beta \not\equiv 0 \pmod{360^\circ}$; otherwise the composite is a translation.
\item Draw the line $(O_1O_2)$.
\item Through $O_1$, draw the ray obtained by rotating the ray $[O_1O_2)$ by $-\tfrac{\alpha}{2}$.
\item Through $O_2$, draw the ray obtained by rotating the ray $[O_2O_1)$ by $+\tfrac{\beta}{2}$.
\item The two rays meet at $\Omega$, the centre of the composite rotation; the angle of the composite is $\alpha + \beta$.
\end{enumerate}
Equivalently, the triangle $O_1O_2\Omega$ has base angles $\tfrac{\alpha}{2}$ at $O_1$ and $\tfrac{\beta}{2}$ at $O_2$ (taken with the appropriate orientation).
\end{methode}

\begin{popfigure}
\begin{center}
\begin{tikzpicture}[scale=1.0, >=stealth]
\coordinate (O1) at (0,0);
\coordinate (O2) at (4.2,0);
\coordinate (Om) at (2.1,2.6);
\coordinate (A) at (0.9,1.0);
\coordinate (B) at (1.7,0.7);
\coordinate (C) at (1.1,1.9);
\coordinate (A1) at (-0.55,1.35);
\coordinate (B1) at (-0.1,0.55);
\coordinate (C1) at (-1.0,1.15);
\coordinate (A2) at (3.15,3.35);
\coordinate (B2) at (3.6,2.6);
\coordinate (C2) at (2.6,2.75);
% triangle ABC
\draw[thick, popblue, fill=popblueL] (A)--(B)--(C)--cycle;
\node[popblue, font=\small] at (1.35,1.35) {$ABC$};
% intermediate triangle
\draw[thick, popmuted, dashed, fill=gray!10] (A1)--(B1)--(C1)--cycle;
\node[popmuted, font=\small] at (-0.55,1.75) {$A_1B_1C_1$};
% final triangle
\draw[thick, poporange, fill=popgold!30] (A2)--(B2)--(C2)--cycle;
\node[poporange, font=\small] at (3.15,2.55) {$A'B'C'$};
% centres
\fill (O1) circle (1.6pt); \node[below, font=\small] at (O1) {$O_1$};
\fill (O2) circle (1.6pt); \node[below, font=\small] at (O2) {$O_2$};
\fill[popdark] (Om) circle (1.8pt); \node[above, font=\small, popdark] at (Om) {$\Omega$};
\draw[popline] (O1)--(O2);
\draw[popline, dashed] (O1)--(Om) (O2)--(Om);
% rotation arcs
\draw[->, popblue] (1.5,0) arc (0:75:1.5);
\node[popblue, font=\small] at (1.15,-0.35) {$\alpha$};
\draw[->, poporange] (5.6,0) arc (0:95:1.4);
\node[poporange, font=\small] at (5.0,-0.4) {$\beta$};
\draw[->, popdark] (3.4,2.6) arc (180:120:1.3);
\node[popdark, font=\small] at (0.9,2.9) {$\alpha+\beta$};
% half angles
\draw[popgreen, ->] (1.1,0) arc (0:51:1.1);
\node[popgreen, font=\footnotesize] at (1.5,0.45) {$\tfrac{\alpha}{2}$};
\draw[popgreen, ->] (3.1,0) arc (180:118:1.1);
\node[popgreen, font=\footnotesize] at (2.75,0.5) {$\tfrac{\beta}{2}$};
\end{tikzpicture}
\end{center}
\legende{The triangle $ABC$ is rotated about $O_1$ through $\alpha$ (giving $A_1B_1C_1$), then about $O_2$ through $\beta$ (giving $A'B'C'$). The net effect is a single rotation about $\Omega$ through $\alpha+\beta$. The base angles of triangle $O_1O_2\Omega$ are $\tfrac{\alpha}{2}$ and $\tfrac{\beta}{2}$.}
\end{popfigure}

\begin{popfigure}
\begin{center}
\begin{tikzpicture}[scale=1.05, >=stealth]
\coordinate (O1) at (0,0);
\coordinate (O2) at (5,0);
\coordinate (Om) at (2.35,2.2);
% axes
\draw[thick, popblue] (-1.3,-0.55) -- (1.5,0.62) node[right, font=\small] {$(d_1)$};
\draw[thick, poporange] (3.6,0.75) -- (6.3,-0.7) node[right, font=\small] {$(d_2)$};
\draw[thick, popdark] (-0.6,0) -- (5.8,0) node[right, font=\small] {$(O_1O_2)$};
\draw[popline, dashed] (O1)--(Om) (O2)--(Om);
\fill (O1) circle (1.6pt); \node[below left, font=\small] at (O1) {$O_1$};
\fill (O2) circle (1.6pt); \node[below right, font=\small] at (O2) {$O_2$};
\fill[popdark] (Om) circle (1.8pt); \node[above, font=\small] at (Om) {$\Omega = (d_1)\cap(d_2)$};
\draw[popgreen, ->] (1.0,0) arc (0:24:1.0);
\node[popgreen, font=\footnotesize] at (1.35,0.28) {$\tfrac{\alpha}{2}$};
\draw[popgreen, ->] (4.0,0) arc (180:164:1.0);
\node[popgreen, font=\footnotesize] at (3.6,0.3) {$\tfrac{\beta}{2}$};
\node[font=\small, align=left] at (2.5,-1.3) {$r_1 = s_{(O_1O_2)}\circ s_{(d_1)}$, \quad $r_2 = s_{(d_2)}\circ s_{(O_1O_2)}$};
\end{tikzpicture}
\end{center}
\legende{The reflection-axes construction: both rotations share the middle axis $(O_1O_2)$, which cancels; the composite is $s_{(d_2)}\circ s_{(d_1)}$, a rotation about the intersection $\Omega$ of the two outer axes.}
\end{popfigure}

\begin{exoresolu}[Composite of two rotations with distinct centres]
Let $O_1$ and $O_2$ be two distinct points. Determine the nature and characteristic elements of $f = r(O_2, 60^\circ) \circ r(O_1, 90^\circ)$.
\tcblower
\textbf{Solution.} The sum of the angles is $90^\circ + 60^\circ = 150^\circ \not\equiv 0 \pmod{360^\circ}$, so $f$ is a \textbf{rotation of angle $150^\circ$}. To find its centre $\Omega$:
\begin{itemize}
\item Through $O_1$, draw the ray obtained by rotating $[O_1O_2)$ by $-\tfrac{90^\circ}{2} = -45^\circ$.
\item Through $O_2$, draw the ray obtained by rotating $[O_2O_1)$ by $+\tfrac{60^\circ}{2} = +30^\circ$.
\end{itemize}
These rays intersect at $\Omega$. Hence $f = r(\Omega, 150^\circ)$, where $\Omega$ is the point such that the triangle $O_1O_2\Omega$ has angles $45^\circ$ at $O_1$ and $30^\circ$ at $O_2$ (and therefore $105^\circ$ at $\Omega$).
\end{exoresolu}

\courssub{The Special Case: Angles Summing to a Multiple of $360^\circ$}

When $\alpha + \beta \equiv 0 \pmod{360^\circ}$, that is $\beta = -\alpha$, the two axes $(d_1)$ and $(d_2)$ are parallel and distinct, and
\[ r(O_2, -\alpha) \circ r(O_1, \alpha) = s_{(d_2)} \circ s_{(d_1)} = t_{\vec{u}}, \]
a translation. Recall that the composite of two reflections with parallel axes is the translation by the vector $2\overrightarrow{H_1H_2}$, where $H_1, H_2$ are the feet of a common perpendicular on $(d_1)$ and $(d_2)$. Hence the translation vector is perpendicular to $(d_1)$ and $(d_2)$, with magnitude twice the distance between the two axes.

\begin{methode}[Finding the translation vector when $\alpha + \beta \equiv 0 \pmod{360^\circ}$]
The most reliable exam technique is to compute the image of one well-chosen point. Since $f = r(O_2,-\alpha)\circ r(O_1,\alpha)$ is a translation $t_{\vec u}$, we have $\vec u = \overrightarrow{M\,f(M)}$ for \emph{any} point $M$. The clever choice is the point $M$ whose first image is the centre of the second rotation:
\begin{enumerate}
\item Choose $M$ such that $r(O_1,\alpha)(M) = O_2$, i.e.\ $M$ is the image of $O_2$ under $r(O_1,-\alpha)$.
\item Then $f(M) = r(O_2,-\alpha)(O_2) = O_2$, because the centre of a rotation is fixed.
\item Conclude: $\vec u = \overrightarrow{MO_2}$.
\end{enumerate}
The geometry of the triangle $O_1MO_2$ (with $O_1M = O_1O_2$ and angle $\alpha$ at $O_1$) then gives the magnitude and direction of $\vec u$ explicitly.
\end{methode}

\begin{exoresolu}[A composite that is a translation]
Let $O_1 \neq O_2$ and let $f = r(O_2, -90^\circ) \circ r(O_1, 90^\circ)$. Show that $f$ is a translation and determine its vector in terms of $\overrightarrow{O_1O_2}$.
\tcblower
\textbf{Solution.} The sum of the angles is $90^\circ + (-90^\circ) = 0$, so $f$ is a translation $t_{\vec u}$. To find $\vec u$, let $M$ be the point such that $r(O_1, 90^\circ)(M) = O_2$, i.e.\ $M$ is obtained from $O_2$ by rotating about $O_1$ through $-90^\circ$. Then
\[ f(M) = r(O_2,-90^\circ)\big(r(O_1,90^\circ)(M)\big) = r(O_2,-90^\circ)(O_2) = O_2, \]
since $O_2$, being the centre of the second rotation, is fixed by it. Hence
\[ \vec u = \overrightarrow{M O_2}. \]
Now $M$ satisfies $O_1M = O_1O_2$ and $\left(\overrightarrow{O_1M}, \overrightarrow{O_1O_2}\right) = 90^\circ$, so the triangle $O_1MO_2$ is right-angled isosceles at $O_1$. By the Pythagorean theorem, $MO_2 = O_1O_2\sqrt{2}$, and the vector $\overrightarrow{MO_2}$ makes an angle of $-45^\circ$ with $\overrightarrow{O_1O_2}$. Thus $f$ is the translation by the vector $\overrightarrow{MO_2}$, of magnitude $O_1O_2\sqrt{2}$, directed along the hypotenuse of the right isosceles triangle built on $[O_1O_2]$.
\end{exoresolu}

\begin{attention}[The composite of two rotations is not always a rotation!]
This is the most frequent trap at the GCE. Before announcing ``rotation of angle $\alpha+\beta$'', you \textbf{must} check whether $\alpha + \beta \equiv 0 \pmod{360^\circ}$. If it is, the composite is a translation (or the identity, when the centres coincide and the angles cancel). Writing ``$r(O_2,-60^\circ)\circ r(O_1,60^\circ) = r(\Omega, 0^\circ)$'' is a serious error: a rotation of angle $0$ is the identity, but here the composite is a genuine translation by a non-zero vector.
\end{attention}

\courssub{Composition of a Rotation and a Translation}

The same reflection-decomposition technique handles mixed composites. A translation $t_{\vec u}$ is the composite of two reflections with parallel axes perpendicular to $\vec u$; a rotation is the composite of two reflections with axes through its centre. Composing $t_{\vec u}$ with $r(O,\theta)$ (with $\theta \not\equiv 0$) and choosing the shared axis to be the line through $O$ perpendicular to $\vec u$, the middle reflections cancel and one obtains a composite of two reflections with \emph{intersecting} axes: a \textbf{rotation of angle $\theta$} about a new centre.

\begin{propriete}[Rotation composed with translation — admitted]
If $\theta \not\equiv 0 \pmod{360^\circ}$, then both $t_{\vec u} \circ r(O,\theta)$ and $r(O,\theta) \circ t_{\vec u}$ are rotations of angle $\theta$, with (generally different) new centres that can be constructed by the shared-axis method. If $\theta \equiv 0$, the composite is of course a translation.
\end{propriete}

\begin{exemplebox}[Quarter-turn followed by a translation]
Let $r = r(O, 90^\circ)$ and $t = t_{\vec u}$. Then $t \circ r$ is a rotation of angle $90^\circ$ whose centre $\Omega$ is obtained as follows: through $O$ draw the line $(d)$ perpendicular to $\vec u$; write $r = s_{(d)} \circ s_{(d_1)}$ and $t = s_{(d_2)} \circ s_{(d)}$ with $(d_2) \parallel (d)$; then $t \circ r = s_{(d_2)} \circ s_{(d_1)}$, a quarter-turn about $\Omega = (d_1) \cap (d_2)$.
\end{exemplebox}

\courssub{Introduction to Plane Isometric Transformations}

Let us step back. Translations, rotations and reflections all share one fundamental feature: they do not distort distances. A triangle carried by any of them lands on a congruent triangle. This common feature deserves a name.

\begin{definition}[Plane isometry]
A \cle{plane isometry} (or \cle{isometric transformation}) is a transformation $f$ of the plane such that for all points $M$ and $N$, with images $M' = f(M)$ and $N' = f(N)$,
\[ M'N' = MN. \]
In words: an isometry is a \emph{distance-preserving} transformation of the plane.
\end{definition}

\begin{propriete}[Conservation properties of isometries — proof examinable for collinearity and midpoints]
Every plane isometry preserves:
\begin{itemize}
\item \textbf{collinearity}: if $A, B, C$ are collinear, so are their images (indeed $B \in [AC] \iff AB + BC = AC$, an equality of distances preserved by $f$);
\item \textbf{midpoints}: the image of the midpoint of $[AB]$ is the midpoint of $[A'B']$ (it is the unique point of $[A'B']$ equidistant from $A'$ and $B'$);
\item \textbf{angles} (in absolute value): by the SSS congruence case, the triangle $A'B'C'$ is congruent to $ABC$, so corresponding angles are equal in magnitude;
\item \textbf{areas}: congruent triangles have equal areas, and every polygonal region decomposes into triangles.
\end{itemize}
Consequently, the image of a line is a line, of a segment a congruent segment, of a circle a circle of the same radius.
\end{propriete}

Translations, rotations, reflections and the identity are isometries; so is any composite of them, since a composite of distance-preserving maps still preserves distances. A natural question arises: are there \emph{other} isometries? The fixed-point analysis below answers this completely.

\courssub{Fixed-Point Analysis of an Isometry}

A \cle{fixed point} of $f$ is a point $M$ such that $f(M) = M$. The number and arrangement of fixed points almost completely determine an isometry.

\begin{propriete}[Fixed points determine the isometry]
Let $f$ be a plane isometry.
\begin{enumerate}
\item If $f$ fixes three non-collinear points, then $f = \mathrm{Id}$ (the identity).
\item If $f$ fixes exactly the points of one line $(d)$ (and no other point), then $f$ is the reflection $s_{(d)}$.
\item If $f$ fixes exactly one point $O$, then $f$ is a rotation of centre $O$ (with $\theta \not\equiv 0$).
\item If $f$ has no fixed point, then $f$ is a translation or a glide reflection.
\end{enumerate}
\end{propriete}

\textbf{Justification of the key step (examinable).} Suppose $f$ fixes three non-collinear points $A, B, C$, and let $M$ be any point, $M' = f(M)$. Since $f$ preserves distances, $M'A = MA$, $M'B = MB$, $M'C = MC$. A point of the plane is uniquely determined by its distances to three non-collinear points: the circles of centre $A$ and radius $MA$ and of centre $B$ and radius $MB$ meet in at most two points, symmetric about $(AB)$, and the third distance to $C$ (which lies off $(AB)$) selects exactly one of them. Hence $M' = M$ for every $M$: $f = \mathrm{Id}$. It follows that an isometry is entirely determined by the images of three non-collinear points, and that the isometries fixing exactly a line, exactly one point, or no point are precisely the reflections, the rotations, and the translations or glide reflections respectively. \hfill$\blacksquare$

\begin{popfigure}
\begin{center}
\begin{tikzpicture}[scale=0.92, >=stealth, font=\small]
\node[draw, rounded corners, fill=popblueL, align=center] (iso) at (0,0) {Plane isometry $f$};
\node[draw, rounded corners, fill=popgreenL, align=center] (three) at (-5.4,-2.2) {3 non-collinear\\fixed points};
\node[draw, rounded corners, fill=popgreenL, align=center] (line) at (-1.8,-2.2) {a whole line\\of fixed points};
\node[draw, rounded corners, fill=popgreenL, align=center] (one) at (1.8,-2.2) {exactly one\\fixed point};
\node[draw, rounded corners, fill=popgreenL, align=center] (none) at (5.4,-2.2) {no fixed\\point};
\node[draw, rounded corners, fill=popgold!40, align=center] (id) at (-5.4,-4.3) {Identity};
\node[draw, rounded corners, fill=popgold!40, align=center] (ref) at (-1.8,-4.3) {Reflection};
\node[draw, rounded corners, fill=popgold!40, align=center] (rot) at (1.8,-4.3) {Rotation};
\node[draw, rounded corners, fill=popgold!40, align=center] (tg) at (5.4,-4.3) {Translation or\\glide reflection};
\draw[->, thick, popdark] (iso) -- (three);
\draw[->, thick, popdark] (iso) -- (line);
\draw[->, thick, popdark] (iso) -- (one);
\draw[->, thick, popdark] (iso) -- (none);
\draw[->, thick, popblue] (three) -- (id);
\draw[->, thick, popblue] (line) -- (ref);
\draw[->, thick, popblue] (one) -- (rot);
\draw[->, thick, popblue] (none) -- (tg);
\end{tikzpicture}
\end{center}
\legende{Classification of plane isometries according to their fixed points.}
\end{popfigure}

\courssub{Direct and Indirect Isometries}

There is one last fundamental distinction. Draw the letter ``R'' on a sheet of paper. Apply a translation or a rotation: the letter still reads ``R''. Apply a reflection: you see a mirror image, reversed. Isometries fall into two families according to whether they preserve or reverse \emph{orientation}.

\begin{definition}[Direct and indirect isometries]
An isometry is called \cle{direct} (or \emph{positive}) if it preserves the orientation of angles: the oriented angle $\left(\overrightarrow{A'B'}, \overrightarrow{A'C'}\right)$ equals $\left(\overrightarrow{AB}, \overrightarrow{AC}\right)$. It is called \cle{indirect} (or \emph{opposite}, \emph{negative}) if it reverses orientation: oriented angles are changed into their opposites.
\end{definition}

\begin{aretenir}[The two families]
\begin{itemize}
\item \textbf{Direct isometries}: the identity, translations, rotations — and every composite of an \emph{even} number of reflections.
\item \textbf{Indirect isometries}: reflections, glide reflections — every composite of an \emph{odd} number of reflections.
\item The composite of two direct or of two indirect isometries is direct; the composite of a direct and an indirect isometry is indirect.
\end{itemize}
\end{aretenir}

\begin{exemplebox}[Orientation check]
The composite $r(O_2,\beta)\circ r(O_1,\alpha)$ of two rotations is a composite of four reflections — an even number — so it is direct, consistent with it being a rotation or a translation. A reflection followed by a rotation is indirect: it must be a reflection or a glide reflection.
\end{exemplebox}

\begin{methode}[GCE exam technique: identifying a composite transformation]
When asked to determine $g \circ f$:
\begin{enumerate}
\item \textbf{State the nature} first (identity, translation, rotation, reflection, glide reflection), using the angle-sum test for rotations and the parity of the number of reflections for orientation.
\item \textbf{Give all characteristic elements}: centre and angle for a rotation; vector for a translation; axis for a reflection; axis and vector for a glide reflection.
\item \textbf{Justify} with a fixed-point argument or by computing the image of one or two well-chosen points.
\end{enumerate}
An answer giving only the nature (e.g.\ ``it is a rotation'') without the centre and angle earns only part of the marks.
\end{methode}

\begin{exercice}[Practice]
Let $O_1, O_2$ be distinct points, $f = r(O_2, 120^\circ) \circ r(O_1, -60^\circ)$ and $g = r(O_2, 60^\circ) \circ r(O_1, -60^\circ)$.
\begin{enumerate}
\item Determine the nature and characteristic elements of $f$.
\item Show that $g$ is a translation, and describe a method to construct its vector.
\item Is $f$ a direct or an indirect isometry? Justify.
\end{enumerate}
\end{exercice}

\begin{corrige}[Exercise 1]
\begin{enumerate}
\item The angle sum is $-60^\circ + 120^\circ = 60^\circ \not\equiv 0 \pmod{360^\circ}$, so $f$ is a rotation of angle $60^\circ$. Its centre $\Omega$ is the intersection of the ray from $O_1$ making angle $-\tfrac{-60^\circ}{2} = +30^\circ$ with $[O_1O_2)$ and the ray from $O_2$ making angle $+\tfrac{120^\circ}{2} = +60^\circ$ with $[O_2O_1)$; the triangle $O_1O_2\Omega$ has base angles $30^\circ$ and $60^\circ$, hence a right angle at $\Omega$.
\item The angle sum is $-60^\circ + 60^\circ = 0$, so $g$ is a translation. Its vector is $\vec u = \overrightarrow{M\,g(M)}$ for any convenient point $M$; choosing $M$ with $r(O_1,-60^\circ)(M) = O_2$ gives $g(M) = O_2$, so $\vec u = \overrightarrow{MO_2}$.
\item $f$ is a composite of two rotations, i.e.\ of four reflections (an even number); equivalently it is a rotation. Hence $f$ is a \textbf{direct} isometry.
\end{enumerate}
\end{corrige}

\begin{savaistu}[Why this matters beyond the classroom]
The classification of plane isometries is the starting point of \emph{crystallography}: the decorative friezes on traditional Cameroonian woven textiles and carved door frames repeat motifs by translations, rotations, reflections and glide reflections, and mathematicians have proved that exactly seven types of frieze patterns and seventeen types of wallpaper patterns exist — no more, no less. Every repeating pattern you see, from embroidered garments to modern fabric prints in Douala markets, belongs to one of these seventeen types!
\end{savaistu}

In the next section we complete this picture by classifying \emph{all} plane isometries, including the glide reflection, the subtle composite of a reflection and a translation parallel to its axis.

\coursec{Classification of Plane Isometries and Glide Reflections}

In the previous section we composed rotations and discovered the family of \cle{plane isometries} — transformations of the plane that preserve distances. We met translations, rotations, reflections and a curious newcomer, the \cle{glide reflection}. The natural question is now: \emph{are there any others?} Could some clever composition produce a fifth, exotic type of isometry? The remarkable answer, which we prove in this section, is \textbf{no}: the four types we already know exhaust all possibilities. This is the \emph{classification theorem} of plane isometries, one of the most beautiful results of the Upper Sixth syllabus and a favourite of GCE examiners.

\courssub{Orientation: the key that separates isometries into two families}

Before classifying, we need a tool to distinguish isometries that ``slide and turn'' the plane from those that ``flip it over''. That tool is \cle{orientation}.

\begin{definition}[Direct and indirect isometries]
Let $f$ be a plane isometry.
\begin{itemize}
\item $f$ is a \cle{direct isometry} (or \emph{positive} isometry) if it \textbf{preserves the orientation of angles}: for any angle $(\vec{u},\vec{v})$, the image angle $(f(\vec{u}),f(\vec{v}))$ has the same measure (same sign) as $(\vec{u},\vec{v})$.
\item $f$ is an \cle{indirect isometry} (or \emph{opposite} isometry) if it \textbf{reverses the orientation of angles}: every oriented angle is sent to its opposite.
\end{itemize}
\end{definition}

Concretely, take a triangle $ABC$ labelled anticlockwise. Under a direct isometry, the image $A'B'C'$ is still labelled anticlockwise; under an indirect isometry, $A'B'C'$ is labelled clockwise — as if the triangle had been turned face down, like a page flipped over.

\begin{propriete}[Orientation of the known isometries]
\begin{itemize}
\item \textbf{Direct:} the identity, every translation, every rotation (and any composite of an even number of reflections).
\item \textbf{Indirect:} every reflection, every glide reflection (and any composite of an odd number of reflections).
\end{itemize}
This follows because a single reflection reverses orientation, and orientation behaviour is multiplicative under composition: reversing twice restores the original orientation.
\end{propriete}

\courssub{The four types of plane isometries}

\begin{propriete}[Classification theorem — proof examinable in outline]
Every plane isometry $f$ is exactly one of the following:
\begin{enumerate}
\item a \textbf{translation} (the identity being the translation by $\vec{0}$);
\item a \textbf{rotation} of angle $\theta \not\equiv 0 \pmod{2\pi}$;
\item a \textbf{reflection} across a line;
\item a \textbf{glide reflection} (reflection across a line $\Delta$ followed by translation by a non-zero vector parallel to $\Delta$).
\end{enumerate}
The first two are direct; the last two are indirect.
\end{propriete}

\textbf{Why is this true? Proof sketch.} The argument rests on two lemmas.

\emph{Lemma 1: an isometry is determined by three non-collinear points} (see the property box below). \emph{Lemma 2: every isometry is a composite of at most three reflections.} Indeed, let $f$ be an isometry and $A,B,C$ three non-collinear points with images $A',B',C'$. We move $A$ to $A'$ by the reflection $s_1$ across the perpendicular bisector of $[AA']$ (if $A=A'$, skip). The point $s_1(B)$ is at distance $AB = A'B'$ from $A'$, so a second reflection $s_2$ across the perpendicular bisector of $[s_1(B)\,B']$ — which passes through $A'$ — fixes $A'$ and sends $s_1(B)$ to $B'$. Finally $s_2\circ s_1(C)$ lies at the correct distances from both $A'$ and $B'$, so it equals either $C'$ or its mirror image across line $(A'B')$; one more reflection $s_3$ across $(A'B')$ finishes the job if needed. Then $f = s_3\circ s_2\circ s_1$ (with possibly fewer factors).

Now count the reflections:
\begin{itemize}
\item \textbf{One reflection:} a reflection (indirect).
\item \textbf{Two reflections:} parallel axes give a \emph{translation}; intersecting axes give a \emph{rotation} (Lessons 28 and 31). Both are direct.
\item \textbf{Three reflections:} one shows (by regrouping two of them) that the composite reduces to a reflection or to a reflection followed by a translation parallel to its axis — a \emph{glide reflection} (indirect).
\end{itemize}
Hence the four types, and no others. $\blacksquare$

\begin{propriete}[Determination by three points]
An isometry is completely determined by the images of three non-collinear points: if two isometries $f$ and $g$ agree on three non-collinear points $A,B,C$, then $f=g$.
\end{propriete}

\textbf{Justification.} If $M$ is any point, its distances to $A,B,C$ equal the distances from $f(M)$ to $A',B',C'$. But a point of the plane is uniquely located by its distances to three non-collinear points (the circles centred at $A',B',C'$ with those radii meet in exactly one point). Hence $f(M)$ is forced, so $f=g$.

\begin{popfigure}
\begin{tikzpicture}[scale=0.9, every node/.style={font=\small}]
% Translation
\begin{scope}[shift={(0,0)}]
\draw[popline] (-0.4,0) -- (2.6,0);
\fill[popblue] (0,0) circle(1.2pt) node[below left]{$A$};
\fill[popblue] (1.4,0) circle(1.2pt) node[below]{$B$};
\fill[popblue] (0.3,1.1) circle(1.2pt) node[above]{$C$};
\draw[popblue, thick] (0,0)--(1.4,0)--(0.3,1.1)--cycle;
\draw[->, popblue] (0.5,0.45) arc(120:40:0.45);
\begin{scope}[shift={(1.1,0.7)}]
\fill[poporange] (0,0) circle(1.2pt) node[below left]{$A'$};
\fill[poporange] (1.4,0) circle(1.2pt) node[below]{$B'$};
\fill[poporange] (0.3,1.1) circle(1.2pt) node[above]{$C'$};
\draw[poporange, thick] (0,0)--(1.4,0)--(0.3,1.1)--cycle;
\draw[->, poporange] (0.5,0.45) arc(120:40:0.45);
\end{scope}
\draw[->, popmuted, dashed] (0.3,1.1) -- (1.4,1.8);
\node[popdark] at (1.3,-0.9) {\textbf{Translation} (direct)};
\end{scope}
% Rotation
\begin{scope}[shift={(5.2,0)}]
\fill[popblue] (0,0) circle(1.2pt) node[below left]{$A$};
\fill[popblue] (1.4,0) circle(1.2pt) node[below]{$B$};
\fill[popblue] (0.3,1.1) circle(1.2pt) node[above]{$C$};
\draw[popblue, thick] (0,0)--(1.4,0)--(0.3,1.1)--cycle;
\draw[->, popblue] (0.5,0.45) arc(120:40:0.45);
\fill[popdark] (0.7,0.35) circle(1.4pt) node[below right]{$\Omega$};
\begin{scope}[shift={(0.7,0.35)}, rotate=100, shift={(-0.7,-0.35)}]
\fill[poporange] (0,0) circle(1.2pt) node[left]{$A'$};
\fill[poporange] (1.4,0) circle(1.2pt) node[above]{$B'$};
\fill[poporange] (0.3,1.1) circle(1.2pt) node[right]{$C'$};
\draw[poporange, thick] (0,0)--(1.4,0)--(0.3,1.1)--cycle;
\draw[->, poporange] (0.5,0.45) arc(120:40:0.45);
\end{scope}
\draw[->, popmuted, dashed] (1.6,0.15) arc(10:110:0.95);
\node[popdark] at (0.9,-0.9) {\textbf{Rotation} (direct)};
\end{scope}
% Reflection
\begin{scope}[shift={(9.6,0)}]
\draw[popgreen, thick, dashed] (1.15,-0.5) -- (1.15,1.7) node[above]{$\Delta$};
\fill[popblue] (0,0) circle(1.2pt) node[below left]{$A$};
\fill[popblue] (0.9,0) circle(1.2pt) node[below]{$B$};
\fill[popblue] (0.2,1.1) circle(1.2pt) node[above]{$C$};
\draw[popblue, thick] (0,0)--(0.9,0)--(0.2,1.1)--cycle;
\draw[->, popblue] (0.35,0.45) arc(120:40:0.4);
\fill[poporange] (2.3,0) circle(1.2pt) node[below right]{$A'$};
\fill[poporange] (1.4,0) circle(1.2pt) node[below]{$B'$};
\fill[poporange] (2.1,1.1) circle(1.2pt) node[above]{$C'$};
\draw[poporange, thick] (2.3,0)--(1.4,0)--(2.1,1.1)--cycle;
\draw[->, poporange] (1.95,0.45) arc(60:140:0.4);
\node[popdark] at (1.15,-0.9) {\textbf{Reflection} (indirect)};
\end{scope}
% Glide reflection
\begin{scope}[shift={(13.6,0)}]
\draw[popgreen, thick, dashed] (-0.3,0.55) -- (3.1,0.55) node[right]{$\Delta$};
\fill[popblue] (0,1.1) circle(1.2pt) node[above]{$A$};
\fill[popblue] (1.2,1.1) circle(1.2pt) node[above]{$B$};
\fill[popblue] (0.25,1.9) circle(1.2pt) node[above]{$C$};
\draw[popblue, thick] (0,1.1)--(1.2,1.1)--(0.25,1.9)--cycle;
\draw[->, popblue] (0.45,1.45) arc(120:40:0.4);
\fill[poporange] (1.6,0) circle(1.2pt) node[below]{$A'$};
\fill[poporange] (2.8,0) circle(1.2pt) node[below]{$B'$};
\fill[poporange] (2.55,0.8) circle(1.2pt) node[right]{$C'$};
\draw[poporange, thick] (1.6,0)--(2.8,0)--(2.55,0.8)--cycle;
\draw[->, poporange] (2.25,0.35) arc(60:140:0.4);
\draw[->, popmuted, dashed] (0,0) -- (1.6,0) node[midway, below]{$\vec{u}\parallel\Delta$};
\node[popdark] at (1.4,-0.9) {\textbf{Glide reflection} (indirect)};
\end{scope}
\end{tikzpicture}
\legende{The same triangle $ABC$ (blue, anticlockwise) and its images (orange) under the four types of isometry. Direct isometries keep the anticlockwise arrow; indirect ones reverse it.}
\end{popfigure}

\courssub{Analytical (complex) form of isometries}

Identifying the plane with $\mathbb{C}$, every isometry has a remarkably simple algebraic description.

\begin{propriete}[Complex form of isometries]
Every plane isometry has one of the following forms, where $a,b\in\mathbb{C}$ and $|a|=1$:
\begin{itemize}
\item \textbf{Direct isometry:} $z' = az + b$. If $a=1$: translation by the vector of affix $b$. If $a\neq 1$: rotation of angle $\arg a$ about the unique fixed point $\omega = \dfrac{b}{1-a}$.
\item \textbf{Indirect isometry:} $z' = a\bar{z} + b$. This is a reflection if it has fixed points (equivalently $a\bar{b}+b=0$), otherwise a glide reflection.
\end{itemize}
\end{propriete}

\textbf{Why $|a|=1$?} Distances are preserved, so $|z'_1 - z'_2| = |a|\,|z_1 - z_2| = |z_1 - z_2|$ for all $z_1,z_2$, forcing $|a|=1$. The conjugate $\bar z$ appears precisely in the indirect case because conjugation reverses the sign of every angle.

\begin{exemplebox}[Reading the nature from the formula]
$f:\ z' = iz + 2 - i$. Here $a = i$, $|a|=1$, $a\neq 1$: a \textbf{rotation} of angle $\arg i = \dfrac{\pi}{2}$, with centre $\omega = \dfrac{2-i}{1-i} = \dfrac{(2-i)(1+i)}{2} = \dfrac{3+i}{2}$.
\end{exemplebox}

\courssub{Recognition strategy: how to classify a given isometry}

\begin{methode}[Classifying an isometry $f$]
\begin{enumerate}
\item \textbf{Look for fixed points.} Solve $f(M)=M$.
   \begin{itemize}
   \item Exactly one fixed point $\Rightarrow$ \textbf{rotation}.
   \item A whole line of fixed points $\Rightarrow$ \textbf{reflection}.
   \item Every point fixed $\Rightarrow$ \textbf{identity}.
   \item No fixed point $\Rightarrow$ translation or glide reflection: go to step 2.
   \end{itemize}
\item \textbf{Test orientation} on one triangle $ABC$: compare the senses of $ABC$ and $A'B'C'$.
   \begin{itemize}
   \item Same sense $\Rightarrow$ direct $\Rightarrow$ \textbf{translation}.
   \item Opposite sense $\Rightarrow$ indirect $\Rightarrow$ \textbf{glide reflection}.
   \end{itemize}
\item \textbf{Conclude and give characteristic elements:} vector of the translation; centre and angle of the rotation; axis of the reflection; axis and glide vector of the glide reflection.
\end{enumerate}
\end{methode}

\begin{attention}[A glide reflection is not a translation]
A glide reflection has \textbf{no fixed point}, exactly like a (non-identity) translation. Fixed points alone cannot separate them! The \emph{orientation test} is decisive: a translation preserves the sense of every triangle, a glide reflection reverses it. Alternatively, apply $f$ twice: for a glide reflection, $f\circ f$ is a translation by $2\vec{u}$ where $\vec{u}$ is the glide vector, while the midpoint of $[M\,f(M)]$ always lies on the axis — never true for a translation.
\end{attention}

\courssub{Composition table of isometries}

Since orientation is multiplicative (reversing twice restores), we obtain:

\begin{propriete}[Parity rule for composition]
Let $f,g$ be isometries. Then $f\circ g$ is
\begin{itemize}
\item \textbf{direct} if $f$ and $g$ have the \emph{same} type (both direct or both indirect);
\item \textbf{indirect} if $f$ and $g$ have \emph{different} types.
\end{itemize}
In symbols: direct $\circ$ direct = direct; indirect $\circ$ indirect = direct; direct $\circ$ indirect = indirect $\circ$ direct = indirect.
\end{propriete}

\begin{popfigure}
\begin{tikzpicture}[every node/.style={font=\small}]
\draw[fill=popblueL] (0,0) rectangle (2.2,0.8);
\draw[fill=popblueL] (0,0.8) rectangle (2.2,1.6);
\draw[fill=popblueL] (2.2,1.6) rectangle (4.4,2.4);
\draw[fill=popblueL] (0,1.6) rectangle (2.2,2.4);
\node at (1.1,2.0) {$f\circ g$};
\node at (3.3,2.0) {$g$ direct};
\node at (5.5,2.0) {$g$ indirect};
\node at (1.1,1.2) {$f$ direct};
\node at (1.1,0.4) {$f$ indirect};
\draw[fill=popgreenL] (2.2,0.8) rectangle (4.4,1.6);
\draw[fill=popgreenL] (4.4,0) rectangle (6.6,0.8);
\node at (3.3,1.2) {direct};
\node at (5.5,1.2) {indirect};
\node at (3.3,0.4) {indirect};
\node at (5.5,0.4) {direct};
\draw[popdark, thick] (0,0) rectangle (6.6,2.4);
\draw[popdark] (2.2,0)--(2.2,2.4) (4.4,0)--(4.4,2.4) (0,0.8)--(6.6,0.8) (0,1.6)--(6.6,1.6);
\node[popmuted, align=center, text width=6cm] at (3.3,-0.8) {Green cells: same parity gives a direct isometry.};
\end{tikzpicture}
\legende{Composition table of direct and indirect isometries (parity rule).}
\end{popfigure}

Beware: the table gives the \emph{orientation class}, not the exact type. For instance, indirect $\circ$ indirect is direct, but it may be a translation \emph{or} a rotation — fixed points decide.

\begin{exoresolu}[Identifying reflection $\circ$ rotation]
Let $r$ be the rotation of centre $O$ and angle $\dfrac{\pi}{2}$, and $s$ the reflection across a line $\Delta$ through $O$. Determine the nature of $f = s\circ r$.
\tcblower
\textbf{Orientation:} $s$ is indirect, $r$ is direct, so $f$ is \textbf{indirect}: a reflection or a glide reflection.
\textbf{Fixed points:} $f(O) = s(r(O)) = s(O) = O$, so $O$ is fixed. An indirect isometry with a fixed point is a \textbf{reflection} (a glide reflection has none).
\textbf{Axis:} take $A\notin\Delta$, $A\neq O$. Then $r(A)=A'$ with $\widehat{AOA'}=90^\circ$, and $f(A)=s(A')=A''$. The axis of $f$ is the perpendicular bisector of $[AA'']$, which passes through $O$: it is the line through $O$ making angle $\dfrac{\pi}{4}$ with $\Delta$. Hence $f$ is the \textbf{reflection across the line through $O$ inclined at $\frac{\pi}{4}$ to $\Delta$}.
\end{exoresolu}

\begin{exoresolu}[Glide reflection $\circ$ translation]
Let $g$ be the glide reflection with axis $\Delta: y=0$ and glide vector $\vec{u}$ of affix $2$ (i.e. $g:\ z' = \bar z + 2$), and $t$ the translation by vector of affix $i$. Determine the nature of $h = g\circ t$.
\tcblower
\textbf{Complex form:} $t(z) = z+i$, so $h(z) = g(z+i) = \overline{z+i}+2 = \bar z + 2 - i$. This has the form $z' = a\bar z + b$ with $|a|=1$: $h$ is \textbf{indirect}.
\textbf{Fixed points?} $z = \bar z + 2 - i$. Writing $z=x+iy$: $x+iy = x - iy + 2 - i$ gives $0 = 2$ on the real parts — impossible. \textbf{No fixed point.}
An indirect isometry with no fixed point is a \textbf{glide reflection}. Its axis is found from midpoints: the midpoint of $[M\,h(M)]$ has affix $\dfrac{z+\bar z+2-i}{2} = x + 1 - \dfrac{i}{2}$, which describes the line $y = -\dfrac{1}{2}$. Squaring: $h\circ h(z) = \overline{\bar z + 2 - i}+2-i = z + 2 + i + 2 - i = z + 4$, a translation by affix $4$, so the glide vector has affix $\dfrac{4}{2} = 2$. \textbf{Conclusion:} $h$ is the glide reflection of axis $y=-\tfrac12$ and glide vector of affix $2$ (i.e. horizontal translation by $2$).
\end{exoresolu}

\begin{methode}[GCE exam technique — ``deduce the nature of $f\circ g$'']
\begin{enumerate}
\item Use the \textbf{parity rule} to decide direct/indirect immediately.
\item Compute the images of \textbf{two or three well-chosen points} (the centre of a rotation, a point of the axis, the origin) instead of heavy algebra.
\item Count \textbf{fixed points} among those images, then conclude with the classification table.
\item Finish by giving the \textbf{characteristic elements} (centre, angle, axis, vector).
\end{enumerate}
\end{methode}

\begin{attention}[Frequent traps at the GCE]
\begin{itemize}
\item Saying ``no fixed point $\Rightarrow$ translation'': false, a glide reflection also has none.
\item Forgetting that the identity counts as a translation (by $\vec 0$) and as a rotation (of angle $0$).
\item Claiming rotation $\circ$ rotation is always a rotation: if the angles sum to $0\pmod{2\pi}$, it is a \emph{translation}.
\item Giving the orientation class when the question asks for the \emph{precise} nature with characteristic elements.
\end{itemize}
\end{attention}

\begin{exercice}[Practice]
\begin{enumerate}
\item Classify $f:\ z' = -i\bar z + 1 + i$ (nature and characteristic elements).
\item $r_1$ and $r_2$ are rotations of angles $\dfrac{\pi}{3}$ and $-\dfrac{\pi}{3}$ about distinct centres. What is the nature of $r_1\circ r_2$?
\item $s_1,s_2$ are reflections across perpendicular lines meeting at $O$. Give the nature of $s_2\circ s_1$.
\end{enumerate}
\end{exercice}

\begin{corrige}[Exercise 1]
\textbf{1.} $|{-i}|=1$ with a conjugate: indirect. Fixed points: $z=-i\bar z+1+i$; with $z=x+iy$, $x+iy = -i(x-iy)+1+i = -y+1+i(1-x)$, so $x=1-y$ and $y=1-x$: the whole line $x+y=1$ is fixed. Hence $f$ is the \textbf{reflection} across the line $x+y=1$.
\textbf{2.} Angles sum to $0$: $r_1\circ r_2$ is a \textbf{translation}.
\textbf{3.} Two reflections across intersecting axes: a \textbf{rotation} about $O$ of angle $2\times 90^\circ = 180^\circ$, i.e. the \textbf{half-turn} (central symmetry) about $O$.
\end{corrige}

\begin{aretenir}[Key points of the section]
\begin{itemize}
\item Every plane isometry is a \textbf{translation, rotation, reflection or glide reflection} — and is a composite of \textbf{at most three reflections}.
\item \textbf{Direct} isometries (translation, rotation): $z'=az+b$, $|a|=1$, preserve orientation. \textbf{Indirect} (reflection, glide reflection): $z'=a\bar z+b$, $|a|=1$, reverse it.
\item An isometry is determined by the images of \textbf{three non-collinear points}.
\item To classify: fixed points first, then orientation, then characteristic elements.
\item Parity rule: same type $\circ$ same type = direct; mixed = indirect.
\end{itemize}
\end{aretenir}

\begin{savaistu}[Why four and only four?]
The classification of plane isometries is a model of mathematical economy: a single hypothesis (preservation of distances) forces every transformation of the plane into one of four rigid ``moves''. The same spirit drives modern crystallography, where the symmetries of tilings and crystals — think of the repetitive patterns on Ndop or Bamileke traditional fabrics — are classified into exactly 17 ``wallpaper groups'', all built from these same four isometries.
\end{savaistu}

\coursec{Homothety: Definition, Construction and Basic Properties}

In the previous section we completed the classification of plane isometries: every transformation that \emph{preserves} distances is a translation, a rotation, a reflection or a glide reflection. We now leave the world of isometries. Photocopiers in a Yaoundé business centre offer ``enlarge to $141\%$'' or ``reduce to $71\%$''; the map of Cameroon hanging in your classroom is a reduced copy of the country; the shadow cast by a lamp grows as you move your hand towards the bulb. In each case the shape is preserved but the \emph{size} changes, and all lengths are multiplied by the same number. The mathematical model of this situation is the \cle{homothety}, the simplest plane transformation that is not an isometry. It is the building block of all plane similarities, which we shall study later in this chapter.

\courssub{Definition of a Homothety}

\begin{experience}[From a lamp shadow to a formula]
Place a point source of light $O$ and a small object at a point $M$. The shadow of $M$ on a screen is a point $M'$, and $O$, $M$, $M'$ are always \emph{aligned}. If the screen is twice as far from $O$ as the object, then $OM' = 2\,OM$ and $M'$ lies on the same side of $O$ as $M$: in vector language $\vec{OM'} = 2\vec{OM}$. If we could place the screen on the \emph{other} side of $O$, we would obtain $\vec{OM'} = -2\vec{OM}$. The whole geometry of the situation is captured by one centre $O$ and one number $k$.
\end{experience}

\begin{definition}[Homothety of centre $O$ and ratio $k$]
Let $O$ be a fixed point of the plane and $k$ a non-zero real number. The \cle{homothety} of \cle{centre} $O$ and \cle{ratio} $k$, denoted $h(O,k)$, is the plane transformation $h$ which maps every point $M$ to the point $M' = h(M)$ defined by
\[
\vec{OM'} = k\,\vec{OM}.
\]
The centre $O$ is the unique fixed point when $k \neq 1$.
\end{definition}

Two special cases are already familiar:
\begin{itemize}
\item If $k = 1$, then $\vec{OM'} = \vec{OM}$, so $M' = M$ for every $M$: $h(O,1)$ is the \cle{identity}.
\item If $k = -1$, then $\vec{OM'} = -\vec{OM}$, so $O$ is the midpoint of $[MM']$: $h(O,-1)$ is the \cle{central symmetry} about $O$.
\end{itemize}

\begin{attention}[The ratio must be non-zero]
The ratio $k = 0$ is excluded: it would send every point of the plane onto $O$, which is not a transformation (it is not bijective). Whenever you write $h(O,k)$ in an examination, keep in mind $k \in \mathbb{R}^{*}$.
\end{attention}

\courssub{Construction of the Image of a Point}

\begin{methode}[Constructing $M' = h(M)$ given $O$, $k$ and $M$]
\begin{enumerate}
\item Draw the line $(OM)$ and the ray $[OM)$ starting at $O$ through $M$.
\item Compute the distance $OM' = |k| \cdot OM$.
\item If $k > 0$, place $M'$ on the ray $[OM)$ (same side as $M$). If $k < 0$, place $M'$ on the \emph{opposite} ray.
\item Check: $O$, $M$, $M'$ must be collinear, with $\vec{OM'} = k\vec{OM}$.
\end{enumerate}
\end{methode}

\begin{definition}[Direct and indirect homothety]
A homothety $h(O,k)$ is called \cle{direct} when $k > 0$: each point and its image lie on the same side of $O$. It is called \cle{indirect} when $k < 0$: each point and its image lie on opposite sides of $O$. An indirect homothety is the composite of the direct homothety $h(O,|k|)$ with the central symmetry about $O$:
\[
h(O,k) = h(O,-1) \circ h(O,|k|) \qquad (k<0),
\]
since $k\vec{OM} = -\bigl(|k|\vec{OM}\bigr)$.
\end{definition}

\begin{popfigure}
\begin{center}
\begin{tikzpicture}[scale=1.05]
% k = 2
\coordinate (O) at (0,0);
\coordinate (A) at (1.6,0.9);
\coordinate (B) at (1.0,-1.1);
\coordinate (A2) at (3.2,1.8);
\coordinate (B2) at (2.0,-2.2);
\draw[popmuted,->] (O) -- (3.6,2.03);
\draw[popmuted,->] (O) -- (2.25,-2.48);
\fill[popblue] (O) circle (1.6pt) node[below left] {$O$};
\fill[popink] (A) circle (1.6pt) node[above] {$A$};
\fill[popink] (B) circle (1.6pt) node[below] {$B$};
\fill[poporange] (A2) circle (1.6pt) node[above] {$A'$};
\fill[poporange] (B2) circle (1.6pt) node[below right] {$B'$};
\node[popdark] at (1.9,2.3) {$k=2$ (direct)};
% k = -1/2
\begin{scope}[xshift=8.2cm]
\coordinate (O2) at (0,0);
\coordinate (C) at (1.8,0.8);
\coordinate (D) at (1.2,-1.2);
\coordinate (C2) at (-0.9,-0.4);
\coordinate (D2) at (-0.6,0.6);
\draw[popmuted,->] (C2) -- (2.05,0.91);
\draw[popmuted,->] (D2) -- (1.4,-1.4);
\fill[popblue] (O2) circle (1.6pt) node[below right] {$O$};
\fill[popink] (C) circle (1.6pt) node[above] {$A$};
\fill[popink] (D) circle (1.6pt) node[below] {$B$};
\fill[poppurple] (C2) circle (1.6pt) node[above left] {$A'$};
\fill[poppurple] (D2) circle (1.6pt) node[below left] {$B'$};
\node[popdark] at (0.4,1.6) {$k=-\tfrac12$ (indirect)};
\end{scope}
\end{tikzpicture}
\end{center}
\legende{Images under $h(O,2)$ (same side of $O$) and under $h(O,-\tfrac12)$ (opposite side, halved distances).}
\end{popfigure}

\begin{attention}[The classic trap with $k<0$]
For a negative ratio the image lies on the \textbf{opposite ray}. Many candidates construct $M'$ on the same side as $M$ and lose all the construction marks. Always write the sign analysis first: ``$k = -\tfrac12 < 0$, so $M' \in$ ray opposite to $[OM)$, with $OM' = \tfrac12 OM$.''
\end{attention}

\courssub{Basic Properties: Distances, Areas, Collinearity, Angles}

\begin{propriete}[Multiplication of distances and areas]
Let $h = h(O,k)$ and let $A' = h(A)$, $B' = h(B)$. Then
\[
\vec{A'B'} = k\,\vec{AB}, \qquad A'B' = |k|\,AB,
\]
and every area is multiplied by $k^{2}$: $\mathcal{A}' = k^{2}\,\mathcal{A}$.
\end{propriete}

\textbf{Proof (examinable).} By the Chasles relation,
\[
\vec{A'B'} = \vec{A'O} + \vec{OB'} = -\vec{OA'} + \vec{OB'} = -k\vec{OA} + k\vec{OB} = k\bigl(\vec{OB} - \vec{OA}\bigr) = k\vec{AB}.
\]
Taking norms gives $A'B' = |k|\,AB$. For areas, note that a triangle $ABC$ has area $\mathcal{A} = \tfrac12\,AB \cdot AC\,|\sin\widehat{BAC}|$; since $A'B' = |k|AB$, $A'C' = |k|AC$ and the angle is preserved (see below), $\mathcal{A}' = \tfrac12 |k|^{2} AB \cdot AC |\sin\widehat{BAC}| = k^{2}\mathcal{A}$. Any polygonal region is a union of triangles, so the result extends. $\blacksquare$

\begin{propriete}[Conservation properties]
A homothety preserves:
\begin{itemize}
\item \cle{collinearity}: if $A, B, C$ are aligned, so are $A', B', C'$;
\item \cle{barycentres} (hence midpoints): if $G$ is the barycentre of $\{(A,\alpha),(B,\beta)\}$, then $G' = h(G)$ is the barycentre of $\{(A',\alpha),(B',\beta)\}$;
\item \cle{angles}: $\widehat{A'B'C'} = \widehat{ABC}$ (geometric angles; oriented angles are preserved because $\vec{A'B'} = k\vec{AB}$ rotates no vector);
\item \cle{parallelism} and \cle{orthogonality}.
\end{itemize}
\end{propriete}

Indeed, if $G$ is the midpoint of $[AB]$, then $\vec{OG} = \tfrac12(\vec{OA}+\vec{OB})$, so $\vec{OG'} = k\vec{OG} = \tfrac12(\vec{OA'}+\vec{OB'})$, which says exactly that $G'$ is the midpoint of $[A'B']$.

\courssub{Image of a Line and of a Circle}

\begin{propriete}[Image of a line]
The image of a line $(d)$ under $h(O,k)$ is a line $(d')$ \textbf{parallel} to $(d)$. If $O \in (d)$, then $(d') = (d)$: every line through the centre is globally invariant.
\end{propriete}

\textbf{Proof.} Take two distinct points $A, B$ on $(d)$. Any point $M \in (d)$ satisfies $\vec{AM} = t\vec{AB}$ for some $t \in \mathbb{R}$. Then $\vec{A'M'} = k\vec{AM} = t\,(k\vec{AB}) = t\,\vec{A'B'}$, so $M'$ lies on the line $(A'B')$, which is parallel to $(AB) = (d)$ since $\vec{A'B'} = k\vec{AB}$. If $O \in (d)$, then for $M \in (d)$ the vector $\vec{OM'} = k\vec{OM}$ is still along $(d)$, so $(d') = (d)$. $\blacksquare$

\begin{propriete}[Image of a circle]
The image under $h(O,k)$ of a circle $\mathcal{C}$ of centre $\Omega$ and radius $R$ is the circle $\mathcal{C}'$ of centre $\Omega' = h(\Omega)$ and radius $R' = |k|R$. In particular $O$, $\Omega$, $\Omega'$ are collinear, with $\vec{O\Omega'} = k\,\vec{O\Omega}$.
\end{propriete}

\textbf{Proof.} $M \in \mathcal{C} \iff \Omega M = R$. Since homothety multiplies all distances by $|k|$, $\Omega'M' = |k|\,\Omega M = |k|R$, so $M'$ describes the circle of centre $\Omega'$ and radius $|k|R$. $\blacksquare$

\begin{popfigure}
\begin{center}
\begin{tikzpicture}[scale=0.9]
\coordinate (O) at (0,0);
\coordinate (Om) at (2.2,0.9);
\coordinate (Om2) at (4.4,1.8);
\draw[popmuted] (O) -- (5.6,2.29);
\draw[popblue,thick] (Om) circle (0.8);
\draw[poporange,thick] (Om2) circle (1.6);
\fill[popink] (O) circle (1.6pt) node[below left] {$O$};
\fill[popblue] (Om) circle (1.6pt) node[below] {$\Omega$};
\fill[poporange] (Om2) circle (1.6pt) node[above right] {$\Omega'$};
\draw[popblue] (Om) -- ++(40:0.8) node[midway,left,scale=0.85] {$R$};
\draw[poporange] (Om2) -- ++(40:1.6) node[midway,right,scale=0.85] {$2R$};
\node[popdark] at (3.1,-1.3) {$h(O,2)$: $\vec{O\Omega'} = 2\vec{O\Omega}$, $R' = 2R$};
\end{tikzpicture}
\end{center}
\legende{Image of a circle under $h(O,2)$: aligned centres, doubled radius.}
\end{popfigure}

\begin{propriete}[Characteristic property]
A plane transformation that maps \emph{every} line to a parallel line is either a \cle{translation} or a \cle{homothety}. If moreover it has a fixed point $O$, it is the homothety of centre $O$ (or the identity).
\end{propriete}

This converse is the standard tool for \emph{recognising} a homothety in examination problems: show that every line is sent to a parallel line and exhibit a fixed point.

\courssub{Image of a Figure: Enlargement and Reduction}

\begin{definition}[Enlargement, reduction, scale factor]
Under $h(O,k)$, a figure $\mathcal{F}$ is mapped to a figure $\mathcal{F}'$ of the \emph{same shape} whose lengths are multiplied by $|k|$. We speak of an \cle{enlargement} when $|k| > 1$, of a \cle{reduction} when $|k| < 1$; $|k|$ is the \cle{scale factor}. When $|k| = 1$ ($k = \pm 1$) the homothety is an isometry (identity or central symmetry).
\end{definition}

\begin{popfigure}
\begin{center}
\begin{tikzpicture}[scale=1.1]
\coordinate (O) at (0,0);
\coordinate (A) at (2.0,0.4);
\coordinate (B) at (2.6,1.6);
\coordinate (C) at (1.6,1.3);
\coordinate (A2) at (3.0,0.6);
\coordinate (B2) at (3.9,2.4);
\coordinate (C2) at (2.4,1.95);
\draw[popmuted] (O) -- (A2); \draw[popmuted] (O) -- (B2); \draw[popmuted] (O) -- (C2);
\draw[popblue,thick] (A) -- (B) -- (C) -- cycle;
\draw[poporange,thick] (A2) -- (B2) -- (C2) -- cycle;
\fill[popink] (O) circle (1.6pt) node[below left] {$O$};
\node[popblue] at (2.05,1.05) {$ABC$};
\node[poporange] at (3.35,1.55) {$A'B'C'$};
\draw[poppurple,dashed] (A) -- (A2);
\draw[poppurple,dashed] (B) -- (B2);
\node[popdark] at (2.6,-0.5) {$h(O,\tfrac32)$: corresponding sides are parallel};
\end{tikzpicture}
\end{center}
\legende{Triangle $ABC$ and its image under $h(O,\tfrac32)$: $(A'B') \parallel (AB)$, etc., and all lengths are multiplied by $\tfrac32$.}
\end{popfigure}

\begin{exoresolu}[Distances, area and image of a circle]
Let $h = h(O,-\tfrac{3}{2})$. Given $AB = 4\ \mathrm{cm}$, a triangle $ABC$ of area $6\ \mathrm{cm^{2}}$, and a circle $\mathcal{C}(\Omega, R)$ with $R = 2\ \mathrm{cm}$ and $O\Omega = 5\ \mathrm{cm}$: find $A'B'$, the area of $A'B'C'$, the radius of $\mathcal{C}' = h(\mathcal{C})$ and the distance $O\Omega'$.
\tcblower
\textbf{Solution.} Here $k = -\tfrac32$, so $|k| = \tfrac32$ and $k^{2} = \tfrac94$.
\begin{align*}
A'B' &= |k|\,AB = \tfrac32 \times 4 = 6\ \mathrm{cm},\\
\mathcal{A}' &= k^{2}\,\mathcal{A} = \tfrac94 \times 6 = 13.5\ \mathrm{cm^{2}},\\
R' &= |k|R = \tfrac32 \times 2 = 3\ \mathrm{cm},\\
O\Omega' &= |k|\,O\Omega = \tfrac32 \times 5 = 7.5\ \mathrm{cm}.
\end{align*}
Because $k < 0$, $\Omega'$ lies on the ray \emph{opposite} to $[O\Omega)$, and $A'B'C'$ is an enlargement (scale factor $\tfrac32 > 1$) of $ABC$.
\end{exoresolu}

\begin{exemplebox}[Recognising a homothety]
Two circles $\mathcal{C}(\Omega, 2)$ and $\mathcal{C}'(\Omega', 6)$ are such that $O$, $\Omega$, $\Omega'$ are aligned with $\vec{O\Omega'} = 3\vec{O\Omega}$. Since the radii satisfy $6 = 3 \times 2$, the homothety $h(O,3)$ maps $\mathcal{C}$ onto $\mathcal{C}'$: the point $O$ is a \emph{centre of similitude} of the two circles. (There is a second one, giving ratio $-3$; this idea returns in the study of similarities.)
\end{exemplebox}

\begin{attention}[Areas are multiplied by $k^{2}$, not by $k$]
A very frequent error at the GCE: ``the ratio is $2$, so the area is doubled.'' Wrong — lengths are multiplied by $|k|$, areas by $k^{2}$ (and volumes by $|k|^{3}$ in space). Also remember $k^{2}$ is positive even for an indirect homothety: an area is never negative.
\end{attention}

\begin{aretenir}[Key points]
\begin{itemize}
\item $h(O,k)$: $\vec{OM'} = k\vec{OM}$, with $k \neq 0$; $O$ is the only fixed point if $k \neq 1$.
\item $k > 0$: direct homothety (same side); $k < 0$: indirect homothety (opposite side), equal to $h(O,|k|)$ composed with the central symmetry about $O$.
\item $h(O,1) = \mathrm{id}$; $h(O,-1) = $ central symmetry about $O$.
\item Distances are multiplied by $|k|$, areas by $k^{2}$; collinearity, midpoints, barycentres, angles and parallelism are preserved.
\item Image of a line: a parallel line (lines through $O$ are invariant). Image of a circle $(\Omega, R)$: circle $(\Omega', |k|R)$ with $\vec{O\Omega'} = k\vec{O\Omega}$.
\item $|k| > 1$: enlargement; $|k| < 1$: reduction. A transformation sending every line to a parallel line is a translation or a homothety.
\end{itemize}
\end{aretenir}

\begin{exercice}[Construction and properties]
Let $OAB$ be a triangle with $OA = 3\ \mathrm{cm}$, $OB = 4\ \mathrm{cm}$ and $\widehat{AOB} = 60^{\circ}$.
\begin{enumerate}
\item Construct $A' = h(A)$ and $B' = h(B)$ where $h = h(O,-2)$.
\item Compute $A'B'$ given that $AB = \sqrt{13}\ \mathrm{cm}$.
\item Show that $(A'B') \parallel (AB)$ and compute the area of $OA'B'$ knowing that the area of $OAB$ is $3\sqrt{3}\ \mathrm{cm^{2}}$.
\end{enumerate}
\end{exercice}

\begin{corrige}[Exercise 1]
\begin{enumerate}
\item Since $k = -2 < 0$, $A'$ is on the ray opposite to $[OA)$ with $OA' = 2 \times 3 = 6\ \mathrm{cm}$; similarly $B'$ is opposite to $[OB)$ with $OB' = 8\ \mathrm{cm}$.
\item $A'B' = |k|\,AB = 2\sqrt{13}\ \mathrm{cm}$.
\item $\vec{A'B'} = -2\vec{AB}$, so the vectors are collinear and $(A'B') \parallel (AB)$. The area is multiplied by $k^{2} = 4$: $\mathcal{A}' = 4 \times 3\sqrt{3} = 12\sqrt{3}\ \mathrm{cm^{2}}$.
\end{enumerate}
\end{corrige}

\begin{savaistu}[Homotheties around us]
The ``zoom'' of a camera lens, the scaling of a technical drawing at $1/50$ used by builders in Douala, and the pantograph — a hinged parallelogram invented in the 17th century to copy drawings at a different scale — are all physical realisations of homotheties. In the next sections we shall describe homotheties analytically and with complex numbers, and then compose them, which will lead us to the general notion of plane similarity.
\end{savaistu}

\coursec{Analytical and Complex Representation of Homothety}

In the previous sections we defined a homothety geometrically: $h(O,k)$ sends each point $M$ to the point $M'$ on the line $(OM)$ such that $\vec{OM'} = k\,\vec{OM}$. Geometry alone, however, becomes heavy when we must compute the images of many points, compose several homotheties, or recognise a homothety hidden inside a formula. The remedy is to translate the vector relation into \cle{coordinates} and into \cle{complex numbers}. This is exactly the spirit of the GCE Advanced Level paper: a transformation is given by a formula such as $z' = 2z - 3 + i$, and you must identify it, find its centre and ratio, and compute images quickly. This section gives you all the tools.

\courssub{Analytical Expression in Cartesian Coordinates}

\textbf{From the vector relation to coordinates.}\par
Let the plane be equipped with an orthonormal frame $(O_0;\vec{i},\vec{j})$. Suppose the centre of the homothety is $O(a,b)$ and the ratio is $k \in \mathbb{R}$, $k \neq 0$. Let $M(x,y)$ and $M'(x',y') = h(M)$. By definition,
\[ \vec{OM'} = k\,\vec{OM}. \]
In coordinates, $\vec{OM} = (x-a,\ y-b)$ and $\vec{OM'} = (x'-a,\ y'-b)$. The vector equality therefore gives
\[ x' - a = k(x-a) \qquad \text{and} \qquad y' - b = k(y-b). \]
Expanding and rearranging:

\begin{propriete}[Analytical expression of a homothety]
The homothety $h(O,k)$ with centre $O(a,b)$ and ratio $k$ sends $M(x,y)$ to $M'(x',y')$ where
\[ x' = kx + (1-k)a, \qquad y' = ky + (1-k)b. \]
Equivalently, in vector form with position vectors $\vec{p} = \vec{O_0M}$, $\vec{p}\,' = \vec{O_0M'}$ and $\vec{o} = \vec{O_0O}$:
\[ \vec{p}\,' = k\,\vec{p} + (1-k)\,\vec{o}. \]
\end{propriete}

Notice the shape of the formula: the coordinates of $M$ are multiplied by $k$, and a \emph{constant} vector $(1-k)\vec{o}$ is added. When $k = 1$ the constant vanishes and $M' = M$: the homothety of ratio $1$ is the \cle{identity}. When $k = -1$, we get $x' = -x + 2a$, $y' = -y + 2b$, which is exactly the \cle{central symmetry} of centre $O$.

\begin{exemplebox}[Computing an image with coordinates]
Let $h$ be the homothety of centre $O(1,1)$ and ratio $k = 2$. Find the image of $M(3,2)$.

We apply the formula with $a = b = 1$ and $1-k = -1$:
\[ x' = 2(3) + (-1)(1) = 5, \qquad y' = 2(2) + (-1)(1) = 3. \]
Hence $M'(5,3)$. Check: $\vec{OM} = (2,1)$ and $\vec{OM'} = (4,2) = 2\vec{OM}$. \checkmark
\end{exemplebox}

\begin{popfigure}
\begin{center}
\begin{tikzpicture}[scale=1.1]
\draw[gray!40, very thin] (-0.5,-0.5) grid (6.5,4.5);
\draw[->, popink, thick] (-0.5,0) -- (6.5,0) node[below left]{$x$};
\draw[->, popink, thick] (0,-0.5) -- (0,4.5) node[below left]{$y$};
\foreach \x in {1,2,3,4,5,6} \draw (\x,0.06) -- (\x,-0.06) node[below, font=\small]{\x};
\foreach \y in {1,2,3,4} \draw (0.06,\y) -- (-0.06,\y) node[left, font=\small]{\y};
\draw[popblue, thick, ->] (1,1) -- (3,2) node[midway, above left, font=\small]{$\vec{OM}$};
\draw[poporange, thick, ->] (1,1) -- (5,3) node[midway, below right, font=\small]{$\vec{OM'} = 2\vec{OM}$};
\draw[popmuted, dashed] (1,1) -- (6,3.5);
\fill[popdark] (1,1) circle (2pt) node[below left]{$O(1,1)$};
\fill[popblue] (3,2) circle (2pt) node[above right]{$M(3,2)$};
\fill[poporange] (5,3) circle (2pt) node[above right]{$M'(5,3)$};
\end{tikzpicture}
\legende{Homothety of centre $O(1,1)$ and ratio $k=2$: $M(3,2) \mapsto M'(5,3)$, with $O$, $M$, $M'$ collinear.}
\end{center}
\end{popfigure}

\courssub{Complex Representation of a Homothety}

\textbf{Why complex numbers?}\par
Identify the plane with the complex plane: each point $M(x,y)$ has affix $z = x + iy$. The two real equations $x' = kx + (1-k)a$ and $y' = ky + (1-k)b$ then merge into \emph{one single} complex equation. Indeed, writing $\omega = a + ib$ for the affix of the centre,
\[ z' = x' + iy' = k(x+iy) + (1-k)(a+ib) = kz + (1-k)\omega. \]

\begin{propriete}[Complex expression of a homothety — key formula]
The homothety $h(O,k)$, where $O$ has affix $\omega$ and $k \in \mathbb{R}$, $k \neq 0$, has complex expression
\[ z' = kz + (1-k)\omega. \]
Conversely, every transformation of the form $z' = az + b$ with $a \in \mathbb{R}$, $a \neq 1$, is a homothety of ratio $a$ whose centre has affix $\omega = \dfrac{b}{1-a}$. (If $a = 1$ and $b \neq 0$, it is a translation.)
\end{propriete}

The proof of the converse is instructive and examinable. Suppose $z' = az + b$ with $a$ real, $a \neq 1$. Look for a fixed point: $z = az + b$ gives $z(1-a) = b$, hence the unique fixed point $\omega = \dfrac{b}{1-a}$. Then
\[ z' - \omega = az + b - \omega = az + (1-a)\omega - \omega = a(z - \omega), \]
which says exactly $\vec{OM'} = a\,\vec{OM}$ in complex form: the transformation is the homothety of centre $\omega$ and ratio $a$. $\blacksquare$

\begin{aretenir}[The two faces of the formula]
\begin{itemize}
\item \textbf{To write} the expression of a known homothety: $z' = kz + (1-k)\omega$.
\item \textbf{To identify} a homothety from $z' = az + b$ ($a$ real, $\neq 1$): ratio $= a$, centre $= \dfrac{b}{1-a}$.
\end{itemize}
\end{aretenir}

\begin{methode}[Identifying a homothety from its complex expression]
Given $z' = az + b$:
\begin{enumerate}
\item Check that $a$ is \textbf{real}. If not, it is not a pure homothety (see the warning below).
\item If $a = 1$: the map is a translation by the vector of affix $b$ (or the identity if $b = 0$).
\item If $a \neq 1$: it is a homothety of ratio $k = a$.
\item Find the centre by solving the fixed-point equation $z = az + b$, i.e. $\omega = \dfrac{b}{1-a}$.
\end{enumerate}
\end{methode}

\begin{attention}[The coefficient of $z$ must be real]
The formula $z' = kz + (1-k)\omega$ is only valid when $k \in \mathbb{R}$. If the coefficient of $z$ is a non-real complex number, for example $z' = 2i\,z + 1$, the transformation also rotates the plane: it is a \cle{similarity} (here a rotation composed with a homothety), \emph{not} a pure homothety. Always check the nature of the coefficient before concluding.
\end{attention}

\courssub{Collinearity of $O$, $M$, $M'$ and Fast Computation of Images}

The complex form makes the fundamental geometric property transparent. From $z' - \omega = k(z - \omega)$, the vectors $\vec{OM}$ and $\vec{OM'}$ have affixes that are \emph{real multiples} of each other; hence they are collinear.

\begin{propriete}[Fundamental property]
For a homothety $h(O,k)$, the points $O$, $M$ and $M' = h(M)$ are always collinear, and $\vec{OM'} = k\,\vec{OM}$. In complex form: $z' - \omega = k(z - \omega)$, so
\[ \arg\left(\frac{z' - \omega}{z - \omega}\right) = 0 \ (\text{mod } \pi), \]
since the quotient is the real number $k$.
\end{propriete}

This gives a rapid computational strategy: to find the image of $M$ of affix $z$, compute $z' = \omega + k(z - \omega)$ — subtract the centre, multiply by the ratio, add the centre back.

\begin{exemplebox}[Fast image computation]
Let $h$ have centre $\omega = 1 + i$ and ratio $k = -\tfrac{1}{2}$. Image of $M$ with affix $z = 3 + 3i$:
\[ z' = \omega + k(z - \omega) = (1+i) - \tfrac{1}{2}\big((3+3i) - (1+i)\big) = (1+i) - \tfrac{1}{2}(2+2i) = 0. \]
So $M'$ has affix $0$ (the origin of the frame). Check: $\vec{OM}$ has affix $2+2i$, and $-\tfrac12(2+2i) = -1-i$, so $M'$ has affix $\omega + (-1-i) = 0$. \checkmark
\end{exemplebox}

\begin{popfigure}
\begin{center}
\begin{tikzpicture}[scale=1.0]
\draw[gray!40, very thin] (-2.5,-1.5) grid (5.5,4.5);
\draw[->, popink, thick] (-2.5,0) -- (5.5,0) node[below left]{Re};
\draw[->, popink, thick] (0,-1.5) -- (0,4.5) node[below left]{Im};
\draw[popmuted, dashed] (-1.5,-0.5) -- (5,3.83);
\draw[popblue, thick, ->] (1,1) -- (3,2);
\draw[poporange, thick, ->] (1,1) -- (4.6,2.8);
\fill[popdark] (1,1) circle (2pt) node[below left]{$\Omega:\ \omega$};
\fill[popblue] (3,2) circle (2pt) node[above right]{$M:\ z$};
\fill[poporange] (4.6,2.8) circle (2pt) node[above right]{$M':\ z'$};
\fill[poppurple] (-1,-0.5) circle (2pt) node[below]{$M'':\ z''$};
\draw[poppurple, thick, ->] (1,1) -- (-1,-0.5);
\node[popmuted, font=\small, align=center] at (3.9,0.6) {$z' - \omega = k(z-\omega)$};
\end{tikzpicture}
\legende{Argand diagram: $\Omega$, $M$, $M'$ are collinear since $z'-\omega = k(z-\omega)$ with $k$ real. Here $k=1.8$ for $M'$ and $k=-1$ for $M''$.}
\end{center}
\end{popfigure}

\courssub{Recovering Centre and Ratio: a GCE-Style Worked Problem}

\begin{exoresolu}[Identifying $z' = 2z - 3 + i$]
The transformation $T$ of the complex plane is defined by $z' = 2z - 3 + i$. Show that $T$ is a homothety and determine its centre and ratio.
\tcblower
\textbf{Solution.}

\textbf{Step 1: nature of the coefficient.} The expression has the form $z' = az + b$ with $a = 2 \in \mathbb{R}$ and $b = -3 + i$. Since $a$ is real and $a \neq 1$, $T$ is a homothety of ratio $k = 2$.

\textbf{Step 2: the centre is the unique fixed point.} Set $z' = z$:
\[ z = 2z - 3 + i \;\Longrightarrow\; -z = -3 + i \;\Longrightarrow\; z = 3 - i. \]
So the centre is $\Omega$ of affix $\omega = 3 - i$.

\textbf{Step 3: verification (recommended in the exam).} Compute
\[ z' - \omega = 2z - 3 + i - (3 - i) = 2z - 6 + 2i = 2(z - 3 + i) = 2(z - \omega). \]
Thus $z' - \omega = 2(z - \omega)$, i.e. $\vec{\Omega M'} = 2\,\vec{\Omega M}$, confirming that $T = h(\Omega, 2)$ with $\omega = 3 - i$. $\blacksquare$

\textbf{Remark.} Using the formula $\omega = \dfrac{b}{1-a}$ directly: $\omega = \dfrac{-3+i}{1-2} = \dfrac{-3+i}{-1} = 3 - i$. Same answer.
\end{exoresolu}

\begin{methode}[Exam tip (GCE): fixed point first]
The centre of a homothety is its \textbf{unique fixed point}. Rather than memorising $\omega = \dfrac{b}{1-a}$, simply set $z' = z$ in the given expression and solve. This is faster, works even when you have forgotten the formula, and the equation $z' - \omega = k(z-\omega)$ then gives the ratio for free.
\end{methode}

\courssub{Direct and Indirect Homotheties}

The sign of the ratio $k$ determines the relative position of $M$ and $M'$ with respect to the centre $O$.

\begin{propriete}[Direct and indirect homothety]
\begin{itemize}
\item If $k > 0$ ($k \neq 1$): the homothety is \cle{direct}. The points $M$ and $M'$ lie on the \textbf{same ray} starting at $O$. The orientation of any triangle is preserved.
\item If $k < 0$: the homothety is \cle{indirect}. The points $M$ and $M'$ lie on \textbf{opposite sides} of $O$. The orientation of any triangle is reversed.
\item The special case $k = -1$ is the central symmetry (half-turn), which is indirect.
\end{itemize}
\end{propriete}

\begin{exemplebox}[Direct vs indirect]
\begin{itemize}
\item $h(O, 3)$: direct. $M$ and $M'$ are on the same side of $O$, $OM' = 3\,OM$.
\item $h(O, -\tfrac{1}{2})$: indirect. $M'$ is on the opposite side of $O$ from $M$, and $OM' = \tfrac{1}{2}OM$.
\end{itemize}
\end{exemplebox}

\courssub{Composition of Two Homotheties (Complex Approach)}

The complex representation makes composition almost trivial. Let $h_1 = h(\omega_1, k_1)$ and $h_2 = h(\omega_2, k_2)$ with $k_1, k_2 \in \mathbb{R}^*$. Their complex forms are
\[ z_1 = k_1 z + (1-k_1)\omega_1, \qquad z' = k_2 z_1 + (1-k_2)\omega_2. \]
Substituting the first into the second:
\[ z' = k_2\bigl(k_1 z + (1-k_1)\omega_1\bigr) + (1-k_2)\omega_2 = (k_1k_2)z + \bigl(k_2(1-k_1)\omega_1 + (1-k_2)\omega_2\bigr). \]

\begin{propriete}[Composition of two homotheties]
Let $h_1 = h(\omega_1, k_1)$ and $h_2 = h(\omega_2, k_2)$ with $k_1, k_2 \in \mathbb{R}^*$. The composition $h_2 \circ h_1$ is:
\begin{itemize}
\item A \textbf{homothety} of ratio $k = k_1k_2$ if $k_1k_2 \neq 1$. Its centre $\omega$ satisfies
  \[ \omega = \frac{k_2(1-k_1)\omega_1 + (1-k_2)\omega_2}{1 - k_1k_2}. \]
\item A \textbf{translation} by the vector of affix $k_2(1-k_1)\omega_1 + (1-k_2)\omega_2$ if $k_1k_2 = 1$ (and this vector is non-zero).
\item The \textbf{identity} if $k_1k_2 = 1$ and the translation vector is zero.
\end{itemize}
\end{propriete}

\begin{exemplebox}[Composition: homothety then homothety]
Let $h_1 = h(0, 2)$ (centre $0$, ratio $2$) and $h_2 = h(1+i, 3)$ (centre $1+i$, ratio $3$).
Then $h_2 \circ h_1$ has ratio $k = 2 \times 3 = 6 \neq 1$, so it is a homothety.
Centre $\omega = \dfrac{3(1-2)\cdot 0 + (1-3)(1+i)}{1-6} = \dfrac{-2(1+i)}{-5} = \dfrac{2}{5}(1+i)$.
Check: $h_1(z) = 2z$, $h_2(z_1) = 3z_1 + (1-3)(1+i) = 3z_1 -2-2i$.
Composition: $z' = 3(2z) -2-2i = 6z -2-2i = 6z + (1-6)\cdot\frac{2}{5}(1+i)$. \checkmark
\end{exemplebox}

\begin{exemplebox}[Composition giving a translation]
Let $h_1 = h(0, 2)$ and $h_2 = h(1, \tfrac{1}{2})$. Then $k_1k_2 = 1$.
$h_1(z) = 2z$, $h_2(z_1) = \tfrac{1}{2}z_1 + (1-\tfrac{1}{2})\cdot 1 = \tfrac{1}{2}z_1 + \tfrac{1}{2}$.
Composition: $z' = \tfrac{1}{2}(2z) + \tfrac{1}{2} = z + \tfrac{1}{2}$. This is a translation by vector $\tfrac{1}{2}$ (affix $0.5$).
\end{exemplebox}

\courssub{Links with the Other Plane Transformations}

The complex formula $z' = kz + (1-k)\omega$ unifies several transformations you already know:

\begin{center}
\begin{tabularx}{\linewidth}{|l|X|X|}
\hline
\rowcolor{popblueL} \textbf{Value of $k$} & \textbf{Transformation} & \textbf{Complex form} \\
\hline
$k = 1$ & Identity & $z' = z$ \\
\hline
$k = -1$ & Central symmetry, centre $\omega$ & $z' = -z + 2\omega$ \\
\hline
$k > 0$, $k \neq 1$ & Direct homothety & $z' = kz + (1-k)\omega$ \\
\hline
$k < 0$, $k \neq -1$ & Indirect homothety & $z' = kz + (1-k)\omega$ \\
\hline
\end{tabularx}
\end{center}

What about \cle{translations}? A translation $z' = z + b$ has coefficient $1$ in front of $z$: it is the case $k = 1$ with a non-zero constant term. One may regard a translation as the \emph{limit} of homotheties whose centre $\omega$ goes to infinity while $k \to 1$ in such a way that $(1-k)\omega \to b$. More concretely, a translation is the companion of homotheties inside the family of maps $z' = az + b$: the homotheties are exactly those with a fixed point, the translations those without one. This point of view will be essential when we classify plane similarities in the next sections.

\begin{attention}[Common mistakes to avoid]
\begin{itemize}
\item Confusing $b$ with the centre: in $z' = 2z - 3 + i$, the centre is \emph{not} $-3+i$; you must divide by $1 - a$ (or solve $z' = z$).
\item Forgetting to check that the coefficient of $z$ is real before declaring "homothety".
\item Writing $z' = kz + (1-k)\omega$ with $\omega$ the affix of the \emph{origin} of the frame instead of the centre of the homothety.
\item Sign errors when $k < 0$: then $1 - k > 1$, and $M'$ lies on the opposite side of $O$ from $M$.
\item In composition, forgetting that the product of ratios $k_1k_2$ determines the nature of the result (homothety if $\neq 1$, translation if $=1$).
\end{itemize}
\end{attention}

\begin{exercice}[Practice]
\begin{enumerate}
\item Write the complex expression of the homothety of centre $\omega = -2 + i$ and ratio $k = 3$, then find the image of the point $A$ of affix $1 - i$.
\item The map $S$ is given by $z' = -\tfrac{1}{2}z + 6 - 3i$. Show that $S$ is a homothety; give its centre and ratio.
\item Given $z' = 5z + 8 + 4i$, find the fixed point and deduce the centre and ratio of the homothety.
\item Is the map $z' = (1+i)z + 2$ a homothety? Justify.
\item Let $h_1 = h(2i, -2)$ and $h_2 = h(1, 3)$. Find the complex expression of $h_2 \circ h_1$ and identify the resulting transformation (give its centre and ratio, or its translation vector).
\item A homothety $h$ maps $A(1+2i)$ to $A'(4+5i)$ and $B(3+i)$ to $B'(7+4i)$. Determine the centre and ratio of $h$ using complex numbers.
\end{enumerate}
\end{exercice}

\begin{corrige}[Exercise 1]
\textbf{1.} $z' = 3z + (1-3)(-2+i) = 3z + 4 - 2i$. Image of $A$: $z'_A = 3(1-i) + 4 - 2i = 7 - 5i$.

\textbf{2.} Coefficient $a = -\tfrac12 \in \mathbb{R}$, $a \neq 1$: homothety of ratio $k = -\tfrac12$. Fixed point: $z = -\tfrac12 z + 6 - 3i \Rightarrow \tfrac32 z = 6 - 3i \Rightarrow \omega = 4 - 2i$. Check: $z' - \omega = -\tfrac12 z + 6 - 3i - 4 + 2i = -\tfrac12 z + 2 - i = -\tfrac12(z - 4 + 2i) = -\tfrac12(z-\omega)$. \checkmark

\textbf{3.} Fixed point: $z = 5z + 8 + 4i \Rightarrow -4z = 8 + 4i \Rightarrow \omega = -2 - i$; ratio $k = 5$.

\textbf{4.} No: the coefficient $1+i$ is not real, so this is a similarity (rotation of angle $\tfrac{\pi}{4}$ composed with a homothety of ratio $\sqrt{2}$), not a pure homothety.

\textbf{5.} $h_1(z) = -2z + (1-(-2))(2i) = -2z + 6i$. $h_2(z_1) = 3z_1 + (1-3)(1) = 3z_1 - 2$.
Composition: $z' = 3(-2z + 6i) - 2 = -6z + 18i - 2$.
Ratio $k = -6 \neq 1$, so it is a homothety. Centre $\omega = \dfrac{-2+18i}{1-(-6)} = \dfrac{-2+18i}{7} = -\tfrac{2}{7} + \tfrac{18}{7}i$.

\textbf{6.} Let $h$ have centre $\omega$ and ratio $k$. Then $z'_A - \omega = k(z_A - \omega)$ and $z'_B - \omega = k(z_B - \omega)$.
Subtracting: $z'_A - z'_B = k(z_A - z_B)$.
$z_A = 1+2i$, $z'_A = 4+5i$, $z_B = 3+i$, $z'_B = 7+4i$.
$z'_A - z'_B = -3+i$, $z_A - z_B = -2+i$.
$k = \dfrac{-3+i}{-2+i} = \dfrac{(-3+i)(-2-i)}{(-2+i)(-2-i)} = \dfrac{6+3i-2i-i^2}{4+1} = \dfrac{7+i}{5} = 1.4 + 0.2i$.
Since $k$ is not real, \textbf{no such homothety exists}. (The data would correspond to a spiral similarity.)
\end{corrige}

\begin{aretenir}[Key points of the section]
\begin{itemize}
\item Analytical form: $x' = kx + (1-k)a$, $y' = ky + (1-k)b$; vector form: $\vec{p}\,' = k\vec{p} + (1-k)\vec{o}$.
\item Complex form: $z' = kz + (1-k)\omega$, with $k$ \textbf{real} and $\omega$ the affix of the centre.
\item To identify: $z' = az + b$, $a \in \mathbb{R}\setminus\{1\}$ $\Rightarrow$ homothety, ratio $a$, centre found by solving $z' = z$.
\item $O$, $M$, $M'$ are collinear because $z' - \omega = k(z - \omega)$ with $k$ real.
\item $k > 0$: direct homothety (preserves orientation); $k < 0$: indirect homothety (reverses orientation).
\item $k = 1$: identity; $k = -1$: central symmetry; coefficient $1$ with $b \neq 0$: translation.
\item Composition $h_2 \circ h_1$: ratio $k_1k_2$. If $k_1k_2 \neq 1$, homothety with centre $\omega = \frac{k_2(1-k_1)\omega_1 + (1-k_2)\omega_2}{1-k_1k_2}$. If $k_1k_2 = 1$, translation (or identity).
\end{itemize}
\end{aretenir}

\coursec{Composition of Homotheties and Advanced Properties}

In the previous section we learnt to describe a homothety analytically: $h(O,k)$ sends the point $M(z)$ to $M'(z')$ with $z' - \omega = k(z - \omega)$, where $\omega$ is the affix of the centre. This compact formula is a powerful tool: composing two homotheties now amounts to composing two affine maps of the complex plane, and the nature of the composite can be read directly from the algebra. That is exactly the programme of this section: we compose homotheties (same centre, then distinct centres), we handle the surprising special case where the product of the ratios equals $1$, we mix homotheties with translations, reflections and rotations, and we finish with two jewels of geometry — the centres of similitude of two circles and the Euler line of a triangle.

\courssub{Composition of Two Homotheties with the Same Centre}

\textbf{Observation.} Enlarge a photograph from centre $O$ by a factor $2$, then enlarge the result from the \emph{same} centre by a factor $3$. A point $M$ goes to $M_1$ with $\vec{OM_1} = 2\vec{OM}$, then to $M'$ with $\vec{OM'} = 3\vec{OM_1} = 6\vec{OM}$. The combined effect is simply the homothety $h(O, 6)$: the ratios multiply.

\begin{propriete}[Composite of two homotheties with the same centre]
Let $O$ be a point of the plane and $k_1, k_2$ two non-zero real numbers. Then
\[ h(O, k_2) \circ h(O, k_1) = h(O, k_1 k_2). \]
In particular the composition is \cle{commutative}: $h(O,k_2)\circ h(O,k_1) = h(O,k_1)\circ h(O,k_2)$. Moreover $h(O,k)^{-1} = h\!\left(O, \dfrac{1}{k}\right)$, so the homotheties of centre $O$ form a \cle{group} under composition.
\end{propriete}

\textbf{Proof (examinable).} For any point $M$, write $M_1 = h(O,k_1)(M)$ and $M' = h(O,k_2)(M_1)$. Then
\[ \vec{OM'} = k_2\,\vec{OM_1} = k_2\bigl(k_1\,\vec{OM}\bigr) = k_1 k_2\,\vec{OM}, \]
so $M' = h(O, k_1k_2)(M)$ for every $M$. $\blacksquare$

\begin{exemplebox}[Same centre]
$h\!\left(O, \tfrac{1}{2}\right) \circ h(O, -4) = h(O, -2)$: an indirect homothety of ratio $-2$. Note that $h(O,-1)\circ h(O,-1) = h(O,1) = \mathrm{id}$: the central symmetry about $O$ is its own inverse.
\end{exemplebox}

\courssub{Composition of Two Homotheties with Distinct Centres}

\textbf{Observation.} Now take two \emph{different} centres $O_1$ and $O_2$. Apply $h_1 = h(O_1, k_1)$ then $h_2 = h(O_2, k_2)$. Using complex affixes with $\omega_1, \omega_2$ the affixes of the centres:
\[ z \xrightarrow{\ h_1\ } z_1 = k_1 z + (1-k_1)\omega_1 \xrightarrow{\ h_2\ } z' = k_2 z_1 + (1-k_2)\omega_2. \]
Substituting,
\[ z' = k_1 k_2\, z + \underbrace{k_2(1-k_1)\omega_1 + (1-k_2)\omega_2}_{\text{a fixed complex number } b}. \]
So the composite has the form $z' = k_1k_2\, z + b$. Two cases arise, according to whether $k_1k_2 = 1$ or not.

\begin{propriete}[Composite of two homotheties with distinct centres — proof examinable]
Let $h_1 = h(O_1, k_1)$ and $h_2 = h(O_2, k_2)$ with $O_1 \neq O_2$, and let $f = h_2 \circ h_1$.
\begin{itemize}
\item If $k_1 k_2 \neq 1$, then $f$ is a \cle{homothety} of ratio $k_1 k_2$ whose centre $\Omega$ lies on the line $(O_1O_2)$.
\item If $k_1 k_2 = 1$, then $f$ is a \cle{translation} of vector $(1-k_2)\,\vec{O_1O_2}$ (equivalently $(k_1 - 1)\,\vec{O_2O_1}$).
\end{itemize}
\end{propriete}

\textbf{Proof.} From the computation above, $z' = k_1k_2\, z + b$ with $b = k_2(1-k_1)\omega_1 + (1-k_2)\omega_2$.

\emph{Case $k_1k_2 \neq 1$.} An affine map $z' = kz + b$ with $k \neq 1$ is the homothety of ratio $k$ and centre of affix $\omega = \dfrac{b}{1-k}$ (since $\omega = k\omega + b$). Here
\[ \omega = \frac{k_2(1-k_1)\omega_1 + (1-k_2)\omega_2}{1 - k_1k_2}, \]
a barycentre of $\omega_1$ and $\omega_2$ with real coefficients summing to $1 - k_1k_2 \neq 0$; hence $\Omega \in (O_1O_2)$.

\emph{Case $k_1k_2 = 1$.} Then $k_2 = 1/k_1$ and $z' = z + b$: a translation. Its vector is obtained from $b$: with $k_2(1-k_1) = k_2 - 1$,
\[ b = (k_2 - 1)\omega_1 + (1 - k_2)\omega_2 = (1-k_2)(\omega_2 - \omega_1), \]
so the translation vector is $(1-k_2)\,\vec{O_1O_2}$. $\blacksquare$

\begin{methode}[Finding the centre of the composite $h_2 \circ h_1$ when $k_1k_2 \neq 1$]
\begin{enumerate}
\item Compute the ratio: $k = k_1 k_2$.
\item Choose a convenient point — the best choice is $O_1$ itself. Compute $P = h_1(O_1) = O_1$, then $P' = h_2(O_1)$: since $h_1$ fixes $O_1$, we get $f(O_1) = h_2(O_1)$.
\item The centre $\Omega$ is the fixed point of $f$; it lies on the line through a point and its image, and also on $(O_1O_2)$. Concretely, $\Omega$ is the intersection of $(O_1O_2)$ with the line through $O_1$ and $f(O_1)$ — or solve $\vec{\Omega f(O_1)} = k\,\vec{\Omega O_1}$.
\item Check with a second point if time allows.
\end{enumerate}
\end{methode}

\begin{exoresolu}[Composite with distinct centres]
Let $O_1(0,0)$, $O_2(4,0)$, $h_1 = h(O_1, 2)$ and $h_2 = h(O_2, \tfrac{1}{3})$. Determine the nature and the characteristic elements of $f = h_2 \circ h_1$.
\tcblower
\textbf{Solution.} The product of the ratios is $k = 2 \times \tfrac13 = \tfrac23 \neq 1$, so $f$ is a homothety of ratio $\tfrac23$ with centre $\Omega$ on the $x$-axis (the line $(O_1O_2)$).

Image of $O_1$: $h_1(O_1) = O_1$, then $h_2(O_1) = P'$ with $\vec{O_2P'} = \tfrac13 \vec{O_2O_1}$, i.e. $P' = O_2 + \tfrac13(O_1 - O_2) = \left(4 - \tfrac43, 0\right) = \left(\tfrac83, 0\right)$.

The centre $\Omega(x,0)$ satisfies $\vec{\Omega P'} = \tfrac23\,\vec{\Omega O_1}$:
\[ \tfrac83 - x = \tfrac23(0 - x) \;\Longrightarrow\; \tfrac83 - x = -\tfrac23 x \;\Longrightarrow\; \tfrac83 = \tfrac13 x \;\Longrightarrow\; x = 8. \]
Hence $f = h\!\left(\Omega(8,0),\ \tfrac23\right)$. Check with $O_2$: $h_1(O_2) = (8,0)$, $h_2(8,0) = O_2 + \tfrac13(8-4, 0) = \left(\tfrac{16}{3}, 0\right)$; and $f(O_2)$ should be $\Omega + \tfrac23(O_2 - \Omega) = \left(8 - \tfrac83, 0\right) = \left(\tfrac{16}{3}, 0\right)$. \checkmark
\end{exoresolu}

\begin{popfigure}
\begin{center}
\begin{tikzpicture}[scale=0.85, >=stealth]
\coordinate (O1) at (0,0);
\coordinate (O2) at (4,0);
\coordinate (Om) at (8,0);
\coordinate (M) at (1,2);
\coordinate (M1) at (2,4);
\coordinate (M2) at (2.667,1.333);
\draw[popline, dashed] (-1,0) -- (9.5,0);
\foreach \p/\n/\pos in {O1/{$O_1$}/below, O2/{$O_2$}/below, Om/{$\Omega$}/below}
  \fill[popink] (\p) circle (2pt) node[\pos, popink]{\n};
\fill[popblue] (M) circle (2pt) node[above left, popblue]{$M$};
\fill[poporange] (M1) circle (2pt) node[above, poporange]{$M_1 = h_1(M)$};
\fill[poppurple] (M2) circle (2pt) node[right, poppurple]{$M' = h_2(M_1)$};
\draw[popblue, ->] (O1) -- (M);
\draw[poporange, ->] (O1) -- (M1);
\draw[poporange, dashed] (O2) -- (M1);
\draw[poppurple, dashed, ->] (O2) -- (M2);
\draw[popgreen, dashed] (Om) -- (M);
\draw[popgreen, dashed, ->] (Om) -- (M2);
\node[popgreen, align=center] at (6.4,2.6) {$\vec{\Omega M'} = \tfrac23 \vec{\Omega M}$};
\end{tikzpicture}
\end{center}
\legende{$f = h(O_2, \tfrac13) \circ h(O_1, 2) = h(\Omega, \tfrac23)$: the centre $\Omega$ lies on the line $(O_1O_2)$, and $\Omega$, $M$, $M'$ are collinear.}
\end{popfigure}

\begin{attention}[Composition is not commutative in general]
When the centres differ, $h_2 \circ h_1 \neq h_1 \circ h_2$ in general: both are homotheties of the \emph{same ratio} $k_1k_2$ but with \emph{different centres} (both on $(O_1O_2)$). Commutativity holds only when $O_1 = O_2$. Always respect the order: in $h_2 \circ h_1$, it is $h_1$ that acts \emph{first}.
\end{attention}

\courssub{The Special Case $k_1 k_2 = 1$: a Translation}

When $k_2 = \dfrac{1}{k_1}$, the composite $h(O_2, 1/k_1) \circ h(O_1, k_1)$ is a translation. Intuitively: the first homothety scales all distances by $|k_1|$, the second scales them back by $1/|k_1|$, so distances are preserved; and every line is mapped to a parallel line, so no point can be fixed — the only isometry with these features and no fixed point is a translation.

\begin{propriete}[Vector of the translation — exam tip]
If $k_1 k_2 = 1$ with $O_1 \neq O_2$, then
\[ h\!\left(O_2, \tfrac{1}{k_1}\right) \circ h(O_1, k_1) = t_{\vec{u}}, \qquad \vec{u} = \left(1 - \tfrac{1}{k_1}\right)\vec{O_1O_2} = (1 - k_2)\,\vec{O_1O_2}. \]
In particular $\vec{u}$ is parallel to $(O_1O_2)$, and $\vec{u} = \vec{0}$ only if $k_1 = k_2 = 1$.
\end{propriete}

\begin{methode}[GCE exam strategy when $k_1k_2 = 1$]
\begin{enumerate}
\item Do \textbf{not} look for a centre — there is none. State immediately that the composite is a translation.
\item To find the vector, compute the image of one easy point (usually $O_1$): $\vec{u} = \vec{O_1 f(O_1)}$.
\item Alternatively, prove directly that for every $M$, the vector $\vec{M M''}$ is constant (independent of $M$): this is the definition of a translation and earns full marks.
\end{enumerate}
\end{methode}

\begin{exoresolu}[The case $k_1k_2 = 1$]
Let $A$ and $B$ be two distinct points, $h_1 = h(A, 3)$ and $h_2 = h\!\left(B, \tfrac13\right)$. Determine $f = h_2 \circ h_1$.
\tcblower
\textbf{Solution.} Since $3 \times \tfrac13 = 1$, $f$ is a translation. Compute the image of $A$: $h_1(A) = A$, then $h_2(A) = A'$ with $\vec{BA'} = \tfrac13 \vec{BA}$. Hence
\[ \vec{u} = \vec{AA'} = \vec{AB} + \vec{BA'} = \vec{AB} + \tfrac13\vec{BA} = \tfrac23 \vec{AB}. \]
Therefore $f = t_{\vec{u}}$ with $\vec{u} = \dfrac{2}{3}\vec{AB}$, agreeing with the formula $(1 - \tfrac13)\vec{AB}$. Direct check for any $M$: $\vec{M M''} = \vec{MA} + \vec{AA'} + \vec{A'M''}$; since $\vec{A'M''} = \vec{AM_1}\cdot\tfrac13 - \dots$ one verifies $\vec{MM''} = \tfrac23\vec{AB}$, constant, as required.
\end{exoresolu}

\courssub{Composing a Homothety with a Translation, a Reflection or a Rotation}

Analytically, homotheties, translations and rotations all have the form $z' = az + b$ with $|a| = k$ (for a homothety $a = k \in \mathbb{R}$, for a rotation $a = e^{i\theta}$, for a translation $a = 1$); reflections have the form $z' = a\bar{z} + b$ with $|a| = 1$. Composing a homothety $h(O,k)$ with any of them multiplies all distances by $|k|$ while preserving angles (direct case) or reversing them (indirect case). The result is therefore never a plain isometry (unless $|k|=1$) but a \cle{similarity}.

\begin{aretenir}[Nature of the composite $h \circ g$ or $g \circ h$]
Let $h$ be a homothety of ratio $k$ ($|k| \neq 1$).
\begin{itemize}
\item $h \circ t_{\vec{u}}$ (or $t_{\vec{u}} \circ h$) is a \cle{homothety} of ratio $k$ (with a new centre). Indeed $z \mapsto kz + b$, $k \neq 1$, always has a fixed point.
\item $h \circ r$ (or $r \circ h$), where $r$ is a rotation of angle $\theta \not\equiv 0 \pmod{2\pi}$, is a \cle{direct similarity}: $z \mapsto k e^{i\theta} z + b$ — a \emph{spiral similarity} (rotation + homothety about the same fixed point).
\item $h \circ s$ (or $s \circ h$), where $s$ is a reflection, is an \cle{indirect similarity}: $z \mapsto k e^{i\theta}\bar{z} + b$.
\end{itemize}
In all cases distances are multiplied by $|k|$ and the composite has exactly one fixed point (its \emph{centre}).
\end{aretenir}

\begin{exemplebox}[Homothety and translation]
$h(O, 2) \circ t_{\vec{u}}$ sends $z$ to $2(z + u) = 2z + 2u$: a homothety of ratio $2$ whose centre $\Omega$ satisfies $\omega = 2\omega + 2u$, i.e. $\omega = -2u$. The centre has moved, the ratio has not.
\end{exemplebox}

\courssub{Advanced Properties: Barycentres, Homothetic Polygons}

\begin{propriete}[Homothety and barycentres]
Every homothety $h(O,k)$ preserves \cle{barycentres}: if $G$ is the barycentre of $\{(A_i, \alpha_i)\}$ with $\sum \alpha_i \neq 0$, then $h(G)$ is the barycentre of $\{(h(A_i), \alpha_i)\}$ with the same weights. In particular midpoints, centroids and centres of mass are mapped to their counterparts.
\end{propriete}

\textbf{Proof.} $\vec{OG'} = k\vec{OG} = \dfrac{k}{\sum\alpha_i}\sum \alpha_i \vec{OA_i} = \dfrac{1}{\sum\alpha_i}\sum \alpha_i\, \vec{OA_i'}$, so $G'$ is the barycentre of the images with the same weights. $\blacksquare$

Consequences: a homothety preserves the \emph{directions} of lines (a line maps to a parallel line), hence it preserves parallelism, orthogonality and angles; it maps a polygon to a \cle{homothetic polygon} whose corresponding sides are parallel and proportional with factor $|k|$.

\begin{propriete}[Concurrency of lines joining corresponding vertices]
If $\mathcal{P}' = h(\mathcal{P})$ is the image of a polygon $\mathcal{P}$ by $h(O,k)$, then all the lines joining corresponding vertices $(AA'), (BB'), (CC'), \dots$ pass through the centre $O$: they are \cle{concurrent}. Conversely, if two polygons have corresponding sides parallel and the lines joining corresponding vertices are concurrent at $O$, then one is the image of the other by a homothety of centre $O$.
\end{propriete}

This converse is a standard GCE tool: to prove two figures are homothetic, exhibit the concurrency point and check the constant ratio $\dfrac{OA'}{OA} = \dfrac{A'B'}{AB}$.

\courssub{Application 1: Two Circles Are Always Homothetic}

\begin{propriete}[Centres of similitude of two circles]
Let $\mathcal{C}(O, R)$ and $\mathcal{C}'(O', R')$ be two circles with $R \neq R'$. There are exactly \textbf{two} homotheties mapping $\mathcal{C}$ onto $\mathcal{C}'$:
\begin{itemize}
\item $h\!\left(I, \dfrac{R'}{R}\right)$, the \cle{external centre of similitude} $I$, dividing $[OO']$ externally: $\vec{IO'} = \dfrac{R'}{R}\vec{IO}$ with $I \notin [OO']$;
\item $h\!\left(J, -\dfrac{R'}{R}\right)$, the \cle{internal centre of similitude} $J$, dividing $[OO']$ internally: $\vec{JO'} = -\dfrac{R'}{R}\vec{JO}$ with $J \in [OO']$.
\end{itemize}
Both $I$ and $J$ lie on the line of centres $(OO')$. The common \emph{external} tangents meet at $I$; the common \emph{internal} tangents (when they exist) meet at $J$. If $R = R'$ and $O \neq O'$, only the internal centre $J$ (midpoint of $[OO']$, central symmetry) and a translation remain.
\end{propriete}

\textbf{Why it works.} Under $h(I, R'/R)$, $O \mapsto O'$ and radii scale by $R'/R$, so $\mathcal{C} \mapsto \mathcal{C}'$. A tangent from $I$ to $\mathcal{C}$ touches at $T$; its image is a line through $I$ tangent to $\mathcal{C}'$ at $T' = h(T)$ — hence the external tangents concur at $I$. The same argument with ratio $-R'/R$ gives $J$ and the internal tangents.

\begin{popfigure}
\begin{center}
\begin{tikzpicture}[scale=0.8]
\coordinate (O) at (0,0);
\coordinate (Op) at (5,0);
\coordinate (I) at (-3,0);
\coordinate (J) at (1.25,0);
\draw[popblue, thick] (O) circle (1.2);
\draw[poppurple, thick] (Op) circle (2);
\draw[popline] (-4.2,0) -- (7.6,0);
\fill[popink] (O) circle (1.8pt) node[below, popink]{$O$};
\fill[popink] (Op) circle (1.8pt) node[below, popink]{$O'$};
\fill[poporange] (I) circle (2pt) node[below, poporange]{$I$};
\fill[popgreen] (J) circle (2pt) node[below, popgreen]{$J$};
\coordinate (T1) at (0.48,1.1);
\coordinate (T2) at (5.8,1.833);
\draw[poporange] (I) -- (T2);
\draw[poporange] (I) -- (0.48,-1.1) -- (5.8,-1.833);
\coordinate (U1) at (-0.554,1.063);
\coordinate (U2) at (4.076,-1.294);
\draw[popgreen, dashed] (U1) -- (U2);
\coordinate (V1) at (-0.554,-1.063);
\coordinate (V2) at (4.076,1.294);
\draw[popgreen, dashed] (V1) -- (V2);
\fill[popblue] (T1) circle (1.5pt);
\fill[poppurple] (T2) circle (1.5pt);
\node[popmuted] at (6.9,-1.6) {external tangents};
\node[popmuted] at (2.6,1.9) {internal tangents};
\end{tikzpicture}
\end{center}
\legende{Two circles with $R' \neq R$: the external centre of similitude $I$ (external tangents) and the internal centre $J$ (internal tangents), both on the line of centres $(OO')$.}
\end{popfigure}

\begin{exemplebox}[Locating the centres of similitude]
Circles $\mathcal{C}(O, 2)$ and $\mathcal{C}'(O', 6)$ with $OO' = 12$. Then $J \in [OO']$ with $OJ : JO' = R : R' = 1 : 3$, so $OJ = 3$; and $I$ on line $(OO')$ outside $[OO']$ with $IO' = 3\,IO$, giving $I$ on the side of $O$ with $OI = 6$. The two homotheties are $h(I, 3)$ and $h(J, -3)$.
\end{exemplebox}

\courssub{Application 2: The Euler Line of a Triangle}

Let $ABC$ be a non-degenerate triangle, $G$ its centroid, $O$ its circumcentre, $H$ its orthocentre, and $\mathcal{C}$ its circumcircle. Consider
\[ h = h\!\left(G, -\tfrac12\right). \]

\begin{enumerate}
\item \textbf{Step 1: $h$ maps vertices to midpoints of the opposite sides.} Since $\vec{GA'} = -\tfrac12 \vec{GA}$ when $A'$ is the midpoint of $[BC]$ (a classical centroid relation, because $\vec{GA} + \vec{GB} + \vec{GC} = \vec{0}$ gives $\vec{GA} = -2\vec{GA'}$), we get $h(A) = A'$, $h(B) = B'$, $h(C) = C'$, the midpoints of the sides.
\item \textbf{Step 2: $h$ maps the circumcircle to the nine-point circle.} The image of $\mathcal{C}(O, R)$ is the circle of centre $O' = h(O)$ and radius $\tfrac{R}{2}$ passing through $A', B', C'$: this is precisely the \cle{nine-point circle} (circle through the three midpoints of the sides), whose centre $N = h(O)$ is the midpoint of $[OH]$.
\item \textbf{Step 3: the altitudes map to the perpendicular bisectors.} The altitude from $A$ is the line through $A$ perpendicular to $(BC)$; its image is the line through $A'$ parallel to it, i.e. perpendicular to $(BC)$ through its midpoint: the perpendicular bisector of $[BC]$. Since the altitudes concur at $H$ and a homothety preserves concurrency, their images concur at $h(H)$. But the perpendicular bisectors concur at $O$. Hence
\[ h(H) = O, \qquad\text{i.e.}\qquad \vec{GO} = -\tfrac12 \vec{GH} \quad\Longleftrightarrow\quad \vec{OH} = 3\,\vec{OG}. \]
\end{enumerate}

\begin{propriete}[Euler line]
In any non-equilateral triangle, the circumcentre $O$, the centroid $G$ and the orthocentre $H$ are collinear on the \cle{Euler line}, with $G$ between $O$ and $H$ and
\[ \vec{GH} = -2\,\vec{GO}, \qquad OH = 3\,OG. \]
The homothety $h(G, -\tfrac12)$ maps $H \mapsto O$ and the circumcircle onto the nine-point circle; equivalently $h(G, -2)$ maps $O \mapsto H$.
\end{propriete}

\begin{popfigure}
\begin{center}
\begin{tikzpicture}[scale=1.05]
\coordinate (A) at (0,3.2);
\coordinate (B) at (-2.6,-1.2);
\coordinate (C) at (2.8,-1.4);
\coordinate (G) at (0.067,0.2);
\coordinate (O) at (0.1,0.55);
\coordinate (H) at (0.033,-0.5);
\draw[popblue, thick] (A) -- (B) -- (C) -- cycle;
\draw[popline, dashed] (O) circle (2.72);
\draw[popmuted] (A) -- (0.033,-0.5);
\draw[popmuted] (B) -- (1.4,0.9);
\draw[popmuted] (C) -- (-1.3,1.0);
\draw[popgold, thick] (-0.3,-1.35) -- (0.35,2.6);
\fill[popblue] (A) circle (1.8pt) node[above, popblue]{$A$};
\fill[popblue] (B) circle (1.8pt) node[left, popblue]{$B$};
\fill[popblue] (C) circle (1.8pt) node[right, popblue]{$C$};
\fill[poporange] (G) circle (2pt) node[below right, poporange]{$G$};
\fill[popgreen] (O) circle (2pt) node[above right, popgreen]{$O$};
\fill[poppurple] (H) circle (2pt) node[below, poppurple]{$H$};
\node[popgold] at (0.75,2.35) {Euler line};
\node[popmuted, align=center] at (-3.4,2.2) {$\vec{GH} = -2\vec{GO}$};
\end{tikzpicture}
\end{center}
\legende{The Euler line: $h(G, -\tfrac12)$ sends $H$ to $O$ and the circumcircle to the nine-point circle, forcing $O$, $G$, $H$ to be collinear with $GH = 2\,GO$.}
\end{popfigure}

\begin{savaistu}[Leonhard Euler]
The collinearity of $O$, $G$, $H$ was published by Leonhard Euler in 1765. The nine-point circle — through the three midpoints of the sides, the three feet of the altitudes and the three midpoints of $[AH], [BH], [CH]$ — is a direct gift of the homothety $h(G, -\tfrac12)$ and of $h(H, \tfrac12)$. It is a beautiful example of how a single transformation can unify nine seemingly unrelated points.
\end{savaistu}

\begin{attention}[Common mistakes]
\begin{itemize}
\item Adding ratios instead of multiplying them: $h(O,k_2)\circ h(O,k_1)$ has ratio $k_1k_2$, never $k_1 + k_2$.
\item Forgetting the order: $h_2 \circ h_1$ means $h_1$ first; with distinct centres, reversing the order changes the centre.
\item Looking for a centre when $k_1k_2 = 1$: the composite is a translation — find its vector instead.
\item Claiming the centre of the composite is the midpoint of $[O_1O_2]$: it is a specific barycentre of $O_1, O_2$, generally not the midpoint.
\item For the Euler line, writing $\vec{GH} = 2\vec{GO}$: the correct relation is $\vec{GH} = -2\vec{GO}$ (the homothety ratio $-\tfrac12$ is negative).
\end{itemize}
\end{attention}

\begin{exercice}[Practice]
\begin{enumerate}
\item Let $O_1(1,0)$, $O_2(5,0)$, $h_1 = h(O_1, -2)$, $h_2 = h(O_2, 3)$. Determine the nature, ratio and centre of $h_2 \circ h_1$.
\item Let $A(0,0)$, $B(6,0)$. Show that $h(B, \tfrac12) \circ h(A, 2)$ is a translation and give its vector.
\item Two circles have radii $3$ and $5$ and centres $8\ \mathrm{cm}$ apart. Locate their two centres of similitude on the line of centres.
\item In a triangle, $G(1,1)$ and $O(4,3)$ are the centroid and circumcentre. Find the orthocentre $H$.
\end{enumerate}
\end{exercice}

\begin{corrige}[Exercise 1]
\textbf{1.} Ratio $k = -6 \neq 1$: a homothety. Image of $O_1$: $h_1(O_1) = O_1$, $h_2(O_1) = O_2 + 3(O_1 - O_2) = (5 - 12, 0) = (-7, 0)$. Centre $\Omega(x,0)$: $-7 - x = -6(1 - x)$ gives $-7 - x = -6 + 6x$, so $x = -\tfrac17$. Hence $h\!\left(\Omega\left(-\tfrac17, 0\right), -6\right)$.

\textbf{2.} Ratios $2 \times \tfrac12 = 1$: translation. Vector $(1 - \tfrac12)\vec{AB} = \tfrac12\vec{AB} = (3, 0)$.

\textbf{3.} Internal centre $J \in [OO']$ with $OJ : JO' = 3 : 5$, so $OJ = 3\ \mathrm{cm}$. External centre $I$ outside $[OO']$ on the side of $O$ with $IO' = \tfrac53 IO$: $IO' - IO = 8$ gives $\tfrac23 IO = 8$, $IO = 12\ \mathrm{cm}$.

\textbf{4.} $\vec{GH} = -2\vec{GO}$: $H = G - 2(O - G) = (1,1) - 2(3,2) = (-5, -3)$.
\end{corrige}

\begin{aretenir}[Key points of the section]
\begin{itemize}
\item $h(O,k_2) \circ h(O,k_1) = h(O, k_1k_2)$ (same centre, commutative group).
\item Distinct centres: $h_2 \circ h_1$ is a homothety of ratio $k_1k_2$ with centre on $(O_1O_2)$ if $k_1k_2 \neq 1$; a translation of vector $(1-k_2)\vec{O_1O_2}$ if $k_1k_2 = 1$.
\item Homothety composed with translation / rotation / reflection gives a similarity (homothety, spiral similarity, or indirect similarity respectively).
\item Homotheties preserve barycentres and directions; lines joining corresponding vertices of homothetic polygons concur at the centre.
\item Two circles of unequal radii admit exactly two centres of similitude $I$ (ratio $R'/R$) and $J$ (ratio $-R'/R$) on the line of centres.
\item $h(G, -\tfrac12)$ sends circumcircle to nine-point circle and $H$ to $O$: hence the Euler line with $\vec{GH} = -2\vec{GO}$.
\end{itemize}
\end{aretenir}

\coursec{Plane Similarities: Concept and Characteristic Elements}

In the previous section we composed homotheties and discovered a striking phenomenon: the composite of two homotheties is either another homothety or a translation. In both cases, the resulting transformation \emph{shrinks or stretches every distance by the same factor}. Photocopiers, zoom lenses, and the maps of Cameroon used by geographers in Yaound\'e all share this feature: shapes are preserved, sizes are not. Transformations with this behaviour are called \cle{similarities}, and they form the natural family unifying everything we have studied so far: isometries (ratio $1$), homotheties, and their composites. This section gives the concept, the characteristic elements, and the first classification results.

\courssub{Definition and First Examples}

\begin{experience}[From similar triangles to similar figures]
Draw a triangle $ABC$ with $AB = 3\ \mathrm{cm}$, $BC = 4\ \mathrm{cm}$, $CA = 5\ \mathrm{cm}$. Now draw a second triangle $A'B'C'$ with $A'B' = 4.5\ \mathrm{cm}$, $B'C' = 6\ \mathrm{cm}$, $C'A' = 7.5\ \mathrm{cm}$. Measure the angles of both triangles: they are equal in pairs, and every side of the second triangle is exactly $1.5$ times the corresponding side of the first. The two triangles have \emph{the same shape} but different sizes. A similarity is precisely the transformation that realises this passage from one figure to the other.
\end{experience}

\begin{definition}[Similarity of ratio $k$]
Let $k$ be a strictly positive real number. A \cle{similarity} of \cle{ratio} $k$ is a transformation $s$ of the plane such that for all points $M$ and $N$, with images $M' = s(M)$ and $N' = s(N)$,
\[
M'N' = k \cdot MN.
\]
A similarity of ratio $k = 1$ is an \cle{isometry}.
\end{definition}

\begin{exemplebox}[Recognising similarities]
\begin{itemize}
\item Every translation, rotation, reflection and glide reflection is a similarity of ratio $k = 1$.
\item Every homothety of ratio $\lambda \neq 0$ is a similarity of ratio $k = |\lambda|$.
\item The composite of a homothety of ratio $|\lambda| = 2$ with a rotation is a similarity of ratio $2$: the homothety multiplies distances by $2$ and the rotation leaves them unchanged.
\end{itemize}
\end{exemplebox}

\begin{propriete}[Elementary properties of a similarity]
Let $s$ be a similarity of ratio $k$. Then:
\begin{enumerate}
\item $s$ is a bijection of the plane;
\item $s$ preserves collinearity and midpoints: the image of a line is a line, and if $I$ is the midpoint of $[MN]$ then $s(I)$ is the midpoint of $[s(M)s(N)]$;
\item $s$ preserves angles (magnitude): for any three non-collinear points $A, B, C$, the angle $\widehat{s(A)s(B)s(C)}$ has the same measure as $\widehat{ABC}$;
\item the image of a circle of radius $R$ is a circle of radius $kR$;
\item areas are multiplied by $k^2$: if a domain $\mathcal{D}$ has area $\mathcal{A}$, its image has area $k^2 \mathcal{A}$.
\end{enumerate}
\end{propriete}

\textbf{Justification (sketch).} Points (1) and (2) follow from the fact that $s$ multiplies all distances by $k$, exactly as for homotheties. For (3), triangles $ABC$ and $A'B'C'$ have proportional sides with factor $k$, hence are similar in the classical sense, so their angles are equal. Point (4) is immediate since a circle is a locus defined by a distance. Point (5) holds for a triangle (base and height both multiplied by $k$) and extends to any domain by decomposition into triangles. \hfill$\blacksquare$

\begin{attention}[Shape is preserved, size is not]
A similarity never distorts a figure: a right-angled triangle stays right-angled, a square stays a square, a circle stays a circle (never an ellipse!). If a transformation sends a circle onto an ellipse, it is \emph{not} a similarity. Conversely, do not confuse ``same shape'' with ``same size'': unless $k = 1$, lengths genuinely change.
\end{attention}

\courssub{Decomposition: Homothety times Isometry}

The key structural result of the theory says that similarities are \emph{exactly} the composites of the two families we already know.

\begin{propriete}[Decomposition theorem]
Every similarity $s$ of ratio $k$ can be written as
\[
s = h \circ f = f' \circ h,
\]
where $h$ is a homothety of ratio $k$ (with arbitrary centre) and $f, f'$ are isometries. Conversely, every composite of a homothety of ratio $k$ and an isometry is a similarity of ratio $k$.
\end{propriete}

\textbf{Proof.} Let $s$ be a similarity of ratio $k$. Fix any point $\Omega$ and let $h$ be the homothety of centre $\Omega$ and ratio $\dfrac{1}{k}$. Consider $f = h \circ s$. For all points $M, N$:
\[
f(M)f(N) = \frac{1}{k}\, s(M)s(N) = \frac{1}{k} \cdot k \cdot MN = MN.
\]
Hence $f$ is an isometry, and $s = h^{-1} \circ f$ where $h^{-1}$ is the homothety of centre $\Omega$ and ratio $k$. The equality $s = f' \circ h$ is obtained by choosing $f' = s \circ h^{-1}$. The converse is immediate: an isometry preserves distances and the homothety multiplies them by $k$, so the composite multiplies them by $k$. \hfill$\blacksquare$

\begin{aretenir}[The big picture]
Similarities $=$ homotheties $\times$ isometries. Since isometries are translations, rotations, reflections and glide reflections, every similarity is built from these bricks together with one homothety. The ratio $k$ comes entirely from the homothety part.
\end{aretenir}

\courssub{Direct and Indirect Similarities}

\begin{definition}[Direct and indirect similarities]
A similarity is called \cle{direct} if it preserves the orientation of angles, and \cle{indirect} if it reverses the orientation of angles. Equivalently, writing $s = h \circ f$ with $h$ a homothety of positive ratio:
\begin{itemize}
\item $s$ is direct when $f$ is a direct isometry (translation or rotation);
\item $s$ is indirect when $f$ is an indirect isometry (reflection or glide reflection).
\end{itemize}
\end{definition}

Concretely: a direct similarity sends a triangle labelled anticlockwise $A, B, C$ onto a triangle labelled anticlockwise $A', B', C'$; an indirect similarity flips the cyclic order, like the image seen in a mirror.

\begin{exemplebox}[Classifying the familiar transformations]
\begin{itemize}
\item \textbf{Direct:} translations ($k=1$), rotations ($k=1$), homotheties of any ratio $\lambda \neq 0$ ($k = |\lambda|$), spiral similarities.
\item \textbf{Indirect:} reflections and glide reflections ($k = 1$), and any composite of a reflection or glide reflection with a homothety ($k > 0$).
\end{itemize}
\end{exemplebox}

\courssub{Spiral Similarities}

Among direct similarities, the most important new objects are the \cle{spiral similarities}, obtained by composing a rotation and a homothety about the \emph{same} centre.

\begin{definition}[Spiral similarity]
A \cle{spiral similarity} is the composite $s = h \circ r$ (equivalently $r \circ h$) of a rotation $r$ of angle $\theta$ and a homothety $h$ of positive ratio $k$, both having the same centre $S$. It is completely characterised by three elements: its \cle{centre} $S$, its \cle{angle} $\theta$ and its \cle{ratio} $k$. For every point $M \neq S$ with image $M' = s(M)$:
\[
SM' = k \cdot SM \qquad \text{and} \qquad \widehat{MSM'} = \theta.
\]
\end{definition}

\begin{attention}[The ratio is always positive]
The ratio of a similarity is by definition a \emph{positive} number. A homothety of \emph{negative} ratio $-k$ about $S$ is nothing other than the homothety of ratio $k$ composed with the half-turn of centre $S$, i.e.\ a rotation of angle $180^\circ$. So the ``minus sign'' is absorbed into the rotation part: a spiral similarity of ratio $k$ and angle $\theta$ with a negative homothety is simply a spiral similarity of ratio $k$ and angle $\theta + 180^\circ$.
\end{attention}

\begin{popfigure}
\begin{center}
\begin{tikzpicture}[scale=1.0, >=stealth]
\coordinate (S) at (0,0);
\coordinate (A) at (1.1,0.4);
\coordinate (B) at (1.7,1.5);
\coordinate (C) at (0.5,1.6);
\coordinate (A1) at (0.18,1.16);
\coordinate (B1) at (-0.35,2.21);
\coordinate (C1) at (-0.83,0.80);
\coordinate (A2) at (0.36,2.32);
\coordinate (B2) at (-0.70,4.42);
\coordinate (C2) at (-1.66,1.60);
\fill[popink] (S) circle (1.6pt) node[below right=1pt] {$S$};
\draw[thick, popblue] (A)--(B)--(C)--cycle;
\fill[popblue] (A) circle (1.3pt) node[below right] {$A$};
\fill[popblue] (B) circle (1.3pt) node[right] {$B$};
\fill[popblue] (C) circle (1.3pt) node[above left] {$C$};
\draw[thick, poporange, dashed] (A1)--(B1)--(C1)--cycle;
\fill[poporange] (A1) circle (1.3pt) node[below right=0pt] {$A_1$};
\fill[poporange] (B1) circle (1.3pt) node[right=1pt] {$B_1$};
\fill[poporange] (C1) circle (1.3pt) node[below left=0pt] {$C_1$};
\draw[thick, poppurple] (A2)--(B2)--(C2)--cycle;
\fill[poppurple] (A2) circle (1.3pt) node[above right=0pt] {$A'$};
\fill[poppurple] (B2) circle (1.3pt) node[above left=0pt] {$B'$};
\fill[poppurple] (C2) circle (1.3pt) node[above left=0pt] {$C'$};
\draw[popmuted, thin] (S)--(A) (S)--(A2) (S)--(B2);
\draw[->, popgreen, thick] (1.5,0.55) arc[start angle=20, end angle=81, radius=1.6];
\node[popgreen] at (1.75,1.35) {$\theta$};
\draw[->, popdark, thick] (0.30,1.93) -- (0.42,2.72);
\node[popdark] at (0.85,2.75) {$\times k$};
\node[popblue, align=center] at (2.6,0.9) {$ABC$};
\node[poporange, align=center] at (-1.6,2.6) {$A_1B_1C_1 = r(ABC)$};
\node[poppurple, align=center] at (-2.3,4.3) {$A'B'C' = h(A_1B_1C_1)$};
\end{tikzpicture}
\end{center}
\legende{A spiral similarity of centre $S$, angle $\theta$ and ratio $k = 2$: the triangle $ABC$ is first rotated about $S$ (triangle $A_1B_1C_1$), then enlarged from $S$ by the factor $k$ (triangle $A'B'C'$). For each vertex, $SM' = k\,SM$ and $\widehat{MSM'} = \theta$.}
\end{popfigure}

\courssub{Determination of a Direct Similarity}

\begin{propriete}[Two point--image pairs determine a direct similarity]
Let $A, B$ be two distinct points and $A', B'$ two distinct points. There exists a \emph{unique} direct similarity $s$ such that $s(A) = A'$ and $s(B) = B'$. Its ratio is necessarily
\[
k = \frac{A'B'}{AB}.
\]
\end{propriete}

\textbf{Justification of uniqueness.} If two direct similarities $s_1, s_2$ send $A$ to $A'$ and $B$ to $B'$, then $s_2^{-1} \circ s_1$ is a direct similarity fixing both $A$ and $B$. Its ratio must be $1$ (since $AB$ is preserved), so it is a direct isometry fixing two points: a translation or rotation fixing two distinct points, hence the identity. Therefore $s_1 = s_2$. \hfill$\blacksquare$

\begin{methode}[Finding the centre of a spiral similarity (geometric construction)]
Given $A \mapsto A'$ and $B \mapsto B'$ (with the segments not related by a mere translation or homothety), to construct the centre $S$ of the spiral similarity:
\begin{enumerate}
\item Compute the ratio $k = \dfrac{A'B'}{AB}$ and read the angle $\theta = \widehat{(AB, A'B')}$ between the two segments.
\item The centre $S$ must see the segment $[AA']$ under the angle $\theta$: it lies on an arc of circle through $A$ and $A'$ (locus of points $M$ with $\widehat{AMA'} = \theta$).
\item Likewise $S$ sees $[BB']$ under the same angle $\theta$: it lies on the corresponding arc through $B$ and $B'$.
\item The intersection of the two arcs (other than any trivial common point) is the centre $S$.
\end{enumerate}
A useful shortcut: $S$ also lies at the intersection of the circles through $A, B$ and the intersection point of lines $(AB)$ and $(A'B')$, and through $A', B'$ and the same point.
\end{methode}

\begin{methode}[Finding the centre and parameters using complex numbers]
In an orthonormal frame identified with $\mathbb{C}$, a direct similarity has the form
\[
z' = a z + b \qquad \text{with } a = k e^{i\theta},\ k>0.
\]
Given $z_A \mapsto z_{A'}$ and $z_B \mapsto z_{B'}$:
\begin{enumerate}
\item Compute $a = \dfrac{z_{B'} - z_{A'}}{z_B - z_A}$. Then $k = |a|$ and $\theta = \arg(a)$.
\item Compute $b = z_{A'} - a z_A$.
\item The centre (fixed point) is $z_S = \dfrac{b}{1-a}$ (provided $a \neq 1$, i.e. not a translation).
\end{enumerate}
\end{methode}

\begin{exoresolu}[Identifying and using a spiral similarity]
The plane is referred to an orthonormal frame. A direct similarity $s$ sends $A(0, 0)$ to $A'(2, 1)$ and $B(4, 0)$ to $B'(2, 5)$.
\begin{enumerate}
\item Determine the ratio and the angle of $s$.
\item Give the complex expression of $s$ and find its centre.
\end{enumerate}
\tcblower
\textbf{Solution.}
\begin{enumerate}
\item $AB = 4$ and $A'B' = \sqrt{0^2 + 4^2} = 4$, so $k = \dfrac{A'B'}{AB} = 1$. The vector $\vec{AB} = (4, 0)$ points along the $x$-axis while $\vec{A'B'} = (0, 4)$ points along the $y$-axis: the angle is $\theta = \dfrac{\pi}{2}$. Since $k = 1$, $s$ is in fact a \emph{rotation} of angle $\dfrac{\pi}{2}$.
\item A direct similarity has complex form $z' = a z + b$ with $a = k e^{i\theta}$. Here $a = e^{i\pi/2} = i$. From $s(A) = A'$: $b = 2 + i$. Check with $B$: $z' = 4i + 2 + i = 2 + 5i$ $\checkmark$. The centre (fixed point) satisfies $z = i z + 2 + i$, so
\[
z = \frac{2 + i}{1 - i} = \frac{(2+i)(1+i)}{(1-i)(1+i)} = \frac{1 + 3i}{2}.
\]
The centre is $S\left(\dfrac{1}{2}, \dfrac{3}{2}\right)$.
\end{enumerate}
\end{exoresolu}

\begin{exoresolu}[Spiral similarity with $k \neq 1$]
A direct similarity $s$ maps $A(1, 2)$ to $A'(3, 0)$ and $B(3, 2)$ to $B'(7, 4)$. Find its ratio, angle, complex form and centre.
\tcblower
\textbf{Solution.}
In complex form: $z_A = 1+2i$, $z_{A'} = 3$, $z_B = 3+2i$, $z_{B'} = 7+4i$.
\[
a = \frac{z_{B'} - z_{A'}}{z_B - z_A} = \frac{(7+4i) - 3}{(3+2i) - (1+2i)} = \frac{4+4i}{2} = 2+2i = 2\sqrt{2}\, e^{i\pi/4}.
\]
Hence $k = 2\sqrt{2}$, $\theta = \dfrac{\pi}{4}$.
$b = z_{A'} - a z_A = 3 - (2+2i)(1+2i) = 3 - (2+4i+2i-4) = 3 - (-2+6i) = 5 - 6i$.
Complex form: $z' = (2+2i)z + 5 - 6i$.
Centre: $z_S = \dfrac{b}{1-a} = \dfrac{5-6i}{1-(2+2i)} = \dfrac{5-6i}{-1-2i} = \dfrac{(5-6i)(-1+2i)}{(-1-2i)(-1+2i)} = \dfrac{-5+10i+6i+12}{5} = \dfrac{7+16i}{5}$.
Thus $S\left(\dfrac{7}{5}, \dfrac{16}{5}\right)$.
\end{exoresolu}

\courssub{Analytical Representation of Similarities}

\begin{propriete}[Analytical form of similarities]
In an orthonormal coordinate system (or complex plane), every similarity can be written in one of the two forms:
\begin{itemize}
\item \textbf{Direct similarity:} $z' = a z + b$ with $a \in \mathbb{C}^*$, $b \in \mathbb{C}$. Ratio $k = |a|$, angle $\theta = \arg(a)$.
\item \textbf{Indirect similarity:} $z' = a \bar{z} + b$ with $a \in \mathbb{C}^*$, $b \in \mathbb{C}$. Ratio $k = |a|$.
\end{itemize}
In vector notation (with $\vec{OM} = \vec{x}$):
\begin{itemize}
\item Direct: $\vec{x}' = k\, R_\theta \vec{x} + \vec{b}$ where $R_\theta$ is the rotation matrix.
\item Indirect: $\vec{x}' = k\, S_\alpha \vec{x} + \vec{b}$ where $S_\alpha$ is the reflection matrix across a line through the origin making angle $\alpha/2$ with the $x$-axis.
\end{itemize}
\end{propriete}

\begin{exemplebox}[Special cases in analytical form]
\begin{itemize}
\item \textbf{Translation:} $z' = z + b$ ($a=1$, $k=1$, $\theta=0$).
\item \textbf{Rotation about $c$:} $z' = e^{i\theta}(z-c) + c = e^{i\theta}z + c(1-e^{i\theta})$.
\item \textbf{Homothety of centre $c$, ratio $\lambda$:} $z' = \lambda(z-c) + c = \lambda z + c(1-\lambda)$. If $\lambda>0$, direct ($k=\lambda$, $\theta=0$); if $\lambda<0$, direct ($k=|\lambda|$, $\theta=\pi$).
\item \textbf{Reflection across line through origin with angle $\alpha$:} $z' = e^{2i\alpha}\bar{z}$ ($k=1$, indirect).
\item \textbf{Glide reflection:} $z' = e^{2i\alpha}\bar{z} + b$ with $b$ parallel to the reflection line.
\end{itemize}
\end{exemplebox}

\courssub{The Family Tree of Plane Transformations}

\begin{popfigure}
\begin{center}
\begin{tikzpicture}[
  box/.style={draw=popink, thick, rounded corners=3pt, fill=popblueL, align=center, text width=3.1cm, minimum height=0.85cm, font=\small},
  ibox/.style={draw=popgreen, thick, rounded corners=3pt, fill=popgreenL, align=center, text width=2.9cm, minimum height=0.8cm, font=\small},
  >=stealth]
\node[box, fill=poppink!25, text width=4.2cm] (sim) at (0,0) {\textbf{Similarities}\\ $M'N' = k\,MN$,\; $k > 0$};
\node[box, align=center] (dir) at (-3.6,-1.8){\textbf{Direct}\\ preserve orientation};
\node[box, align=center] (ind) at (3.6,-1.8){\textbf{Indirect}\\ reverse orientation};
\node[ibox, align=center] (iso) at (-5.6,-3.6){\textbf{Isometries} $k=1$\\ translations, rotations};
\node[ibox, align=center] (hom) at (-1.6,-3.6){\textbf{Homotheties}\\ ratio $\lambda \neq 0$};
\node[ibox, align=center] (spi) at (-3.6,-5.2){\textbf{Spiral similarities}\\ rotation $\circ$ homothety, same centre};
\node[ibox] (refl) at (3.6,-3.6) {\textbf{Reflections, glide reflections} $\circ$ homothety};
\draw[->, thick, popink] (sim) -- (dir);
\draw[->, thick, popink] (sim) -- (ind);
\draw[->, thick, popgreen] (dir) -- (iso);
\draw[->, thick, popgreen] (dir) -- (hom);
\draw[->, thick, popgreen] (dir) -- (spi);
\draw[->, thick, popgreen] (ind) -- (refl);
\end{tikzpicture}
\end{center}
\legende{Family tree of plane transformations. Every similarity is direct or indirect; the direct family contains the isometries, the homotheties and their generic composites, the spiral similarities.}
\end{popfigure}

\begin{savaistu}[Similarities everywhere]
When a cartographer at the National Institute of Cartography in Yaound\'e produces a map of the Littoral region at scale $1:50\,000$, he is applying a similarity of ratio $k = \dfrac{1}{50000}$: every angle on the ground (road junctions, river bends) is exactly preserved on the map, which is why bearings read on the map are reliable. Spiral similarities also appear in nature: the growth of a snail shell is modelled by iterating a spiral similarity — each new chamber is the previous one rotated and enlarged by a fixed factor, producing the famous logarithmic spiral.
\end{savaistu}

\begin{exercice}[Recognising and determining similarities]
\begin{enumerate}
\item Classify each transformation as direct or indirect and give its ratio: (a) the reflection across the line $y = x$; (b) the homothety of centre $O$ and ratio $-3$; (c) the composite of a rotation of angle $\dfrac{\pi}{3}$ with a homothety of ratio $2$ about the same centre.
\item $A(1, 1)$, $B(3, 1)$ have images $A'(0, 2)$, $B'(0, 6)$ under a direct similarity $s$. Find the ratio, the angle, the complex form $z' = az + b$ and the centre of $s$.
\item A similarity of ratio $k = 3$ transforms a triangle of area $5\ \mathrm{cm}^2$. What is the area of the image? What is the image of its inscribed circle of radius $1.2\ \mathrm{cm}$?
\item A direct similarity $s$ has ratio $k = \frac{1}{2}$ and angle $\theta = -\frac{\pi}{3}$. It maps $A(2, 0)$ to $A'$. Find $A'$ if the centre of $s$ is the origin. What is the complex form of $s$?
\item The vertices of a triangle are $P(1, 0)$, $Q(0, 2)$, $R(-1, 0)$. A similarity of ratio $2$ maps this triangle to $P'Q'R'$ with $P'(3, 1)$, $Q'(1, 5)$. Determine whether the similarity is direct or indirect, find the image $R'$, and give the transformation in complex form.
\end{enumerate}
\end{exercice}

\begin{corrige}[Exercise n\textdegree~1]
\begin{enumerate}
\item (a) Indirect, $k = 1$. (b) Direct (a homothety preserves orientation), $k = |-3| = 3$; the negative sign corresponds to an extra half-turn. (c) Direct spiral similarity, $k = 2$, angle $\dfrac{\pi}{3}$.
\item $AB = 2$, $A'B' = 4$, so $k = 2$. $\vec{AB} = (2,0)$, $\vec{A'B'} = (0,4)$: angle $\dfrac{\pi}{2}$. Hence $a = 2i$; from $s(A) = A'$: $2i(1+i) + b = 2i$, giving $b = 2$. So $z' = 2iz + 2$. Fixed point: $z = \dfrac{2}{1 - 2i} = \dfrac{2 + 4i}{5}$; centre $S\left(\dfrac{2}{5}, \dfrac{4}{5}\right)$.
\item Area multiplied by $k^2 = 9$: $45\ \mathrm{cm}^2$. The inscribed circle maps to a circle of radius $3 \times 1.2 = 3.6\ \mathrm{cm}$.
\item Centre at origin $\Rightarrow$ $z' = a z$ with $a = \frac{1}{2} e^{-i\pi/3} = \frac{1}{2}\left(\frac{1}{2} - i\frac{\sqrt{3}}{2}\right) = \frac{1}{4} - i\frac{\sqrt{3}}{4}$. $z_A = 2$, so $z_{A'} = 2a = \frac{1}{2} - i\frac{\sqrt{3}}{2}$, i.e. $A'\left(\frac{1}{2}, -\frac{\sqrt{3}}{2}\right)$.
\item $PQ = \sqrt{5}$, $P'Q' = \sqrt{(3-1)^2 + (1-5)^2} = \sqrt{20} = 2\sqrt{5}$, so $k=2$ (matches). Vector $\vec{PQ} = (-1, 2)$, $\vec{P'Q'} = (-2, 4) = 2\vec{PQ}$. The angle is $0$, so the similarity is a homothety (direct). $R' = P' + 2\vec{PR} = (3,1) + 2(-2,0) = (-1, 1)$. Complex form: $a=2$, $b = z_{P'} - 2z_P = (3+i) - 2 = 1+i$. So $z' = 2z + 1 + i$.
\end{enumerate}
\end{corrige}

\begin{aretenir}[Key points of the section]
\begin{itemize}
\item A similarity of ratio $k > 0$ multiplies \emph{all} distances by $k$; it preserves angles and shapes, and multiplies areas by $k^2$.
\item Every similarity is the composite of a homothety of ratio $k$ and an isometry.
\item Direct similarities preserve orientation (translations, rotations, homotheties, spiral similarities); indirect ones reverse it (reflections and glide reflections composed with homotheties).
\item A spiral similarity is characterised by its centre, its angle and its ratio; the ratio is always positive.
\item A direct similarity is uniquely determined by the images of two distinct points, with $k = \dfrac{A'B'}{AB}$.
\item Analytical forms: direct $z' = az+b$ ($a=k e^{i\theta}$), indirect $z' = a\bar{z}+b$ ($k=|a|$).
\end{itemize}
\end{aretenir}

\coursec{Analytical Study of Similarities and Applications}

In the previous section we defined a \cle{plane similarity} geometrically: a transformation that multiplies all distances by a fixed positive ratio $k$ and preserves angles (in magnitude). We saw that a similarity is \emph{direct} when it preserves orientation and \emph{indirect} when it reverses it. This geometric description is beautiful, but for computation it is clumsy. The great power of the Argand diagram is that it turns every similarity into a one-line formula. In this section we give the complex (analytical) form of similarities, learn to read off their characteristic elements from the formula, study composition, and finish with GCE-style synthesis problems.

\courssub{The Complex Form of a Similarity}

\textbf{Observation.} Consider the map $S$ of the plane into itself defined, in the Argand diagram, by
\[ z' = az + b, \qquad a \in \mathbb{C},\ a \neq 0,\ b \in \mathbb{C}. \]
Take two points $M_1(z_1)$ and $M_2(z_2)$ with images $M_1'(z_1')$ and $M_2'(z_2')$. Then
\[ |z_1' - z_2'| = |a(z_1 - z_2)| = |a|\,|z_1 - z_2|, \]
so $M_1'M_2' = |a|\,M_1M_2$: all distances are multiplied by the same positive constant $|a|$. Moreover,
\[ \arg\!\left(\frac{z_2'-z_1'}{z_2-z_1}\right) = \arg(a) \pmod{2\pi}, \]
so every directed angle is increased by the same amount $\arg(a)$: angles are preserved in size \emph{and} orientation. This is exactly a direct similarity of ratio $|a|$ and angle $\arg(a)$.

\begin{propriete}[Complex form of a similarity — proof examinable]
Every \cle{direct plane similarity} has a unique complex expression
\[ z' = az + b, \qquad a \neq 0, \]
with \textbf{ratio} $k = |a|$ and \textbf{angle} $\theta = \arg(a)$. Every \cle{indirect similarity} has a unique complex expression
\[ z' = a\bar{z} + b, \qquad a \neq 0, \]
with ratio $|a|$ (angles are preserved in magnitude but reversed in orientation). Conversely, every map of one of these two forms is a similarity.
\end{propriete}

\textbf{Proof (direct case).} Let $S$ be a direct similarity of ratio $k$. Pick any point $A$ with affix $\alpha$ and image $A'(\alpha')$. For any $M(z)$ with image $M'(z')$, the similarity preserves angles and multiplies distances by $k$, so
\[ \frac{z' - \alpha'}{z - \alpha} \ \text{has modulus } k \text{ and constant argument } \theta, \]
i.e.\ $z' - \alpha' = a(z - \alpha)$ with $a = k e^{i\theta}$. Setting $b = \alpha' - a\alpha$ gives $z' = az + b$. Uniqueness follows because $a = \dfrac{z_1' - z_2'}{z_1 - z_2}$ is forced by any two distinct points and their images. \hfill$\blacksquare$

\begin{propriete}[Reading off the geometric elements — direct case]
Given a direct similarity $S:\ z' = az + b$ with $a \neq 1$, the map has a unique \textbf{fixed point} $\Omega$ (the \cle{centre}), whose affix is
\[ \omega = \frac{b}{1 - a}. \]
Then $S$ is the composite (in either order) of the homothety $h(\Omega, |a|)$ and the rotation $r(\Omega, \arg a)$: a \cle{spiral similarity} centred at $\Omega$. If $a = 1$, then $z' = z + b$ is the \textbf{translation} by the vector of affix $b$.
\end{propriete}

\textbf{Derivation.} A fixed point satisfies $\omega = a\omega + b$, i.e.\ $(1-a)\omega = b$. If $a \neq 1$ this gives $\omega = \dfrac{b}{1-a}$ uniquely. Subtracting $\omega = a\omega + b$ from $z' = az + b$ yields $z' - \omega = a(z - \omega)$, which says precisely: rotate about $\Omega$ through $\arg a$ and dilate from $\Omega$ by factor $|a|$. \hfill$\blacksquare$

\begin{attention}[Ratio versus coefficient]
Do not confuse the \textbf{ratio} of the similarity, which is $|a|$ and is \emph{always a positive real number}, with the complex coefficient $a$ itself. If $a = -2$, the ratio is $2$, not $-2$; the negative sign contributes an angle of $\pi$. Likewise the angle is $\arg(a)$, not $a$.
\end{attention}

\begin{popfigure}
\begin{tikzpicture}[scale=1.5, >=stealth]
\draw[->, popline] (-0.5,0) -- (4.6,0) node[below left] {Re};
\draw[->, popline] (0,-0.5) -- (0,3.6) node[below left] {Im};
\coordinate (O) at (0,0);
\coordinate (A) at (0.8,0.4);
\coordinate (B) at (2.0,0.6);
\coordinate (C) at (1.1,1.3);
\coordinate (Ap) at (2.6,1.2);
\coordinate (Bp) at (3.7,2.1);
\coordinate (Cp) at (2.3,2.9);
\fill[popblueL] (A)--(B)--(C)--cycle;
\draw[popblue, thick] (A)--(B)--(C)--cycle;
\fill[poppink!30] (Ap)--(Bp)--(Cp)--cycle;
\draw[poppink, thick] (Ap)--(Bp)--(Cp)--cycle;
\foreach \p/\n in {A/A, B/B, C/C} \fill[popblue] (\p) circle (0.035) node[below left, popblue] {\small $\n$};
\foreach \p/\n in {Ap/A', Bp/B', Cp/C'} \fill[poppink] (\p) circle (0.035) node[above right, poppink] {\small $\n'$};
\draw[->, poporange, thick, dashed] (A) -- (Ap);
\draw[->, poporange, thick, dashed] (C) -- (Cp);
\draw[popgreen, thick, ->] (0.9,0.15) arc (8:52:0.95);
\node[popgreen] at (1.45,0.62) {\small $\theta=\arg a$};
\node[popdark, align=center] at (3.6,0.5) {\small $z'=az+b$\\ \small ratio $|a|$};
\end{tikzpicture}
\legende{A direct similarity $z' = az + b$ sends triangle $ABC$ to the enlarged, rotated triangle $A'B'C'$: sides are multiplied by $|a|$ and turned through $\arg(a)$.}
\end{popfigure}

\courssub{Case-by-Case Identification}

\begin{methode}[Identifying a transformation given by $z' = az + b$ or $z' = a\bar{z} + b$]
\begin{enumerate}
\item Compute $|a|$ and $\arg(a)$ (for $a \neq 0$).
\item \textbf{Direct case ($z' = az + b$):}
  \begin{itemize}
  \item If $a = 1$: \textbf{translation} by the vector of affix $b$.
  \item If $a \in \mathbb{R}$, $a \neq 1$: \textbf{homothety} of ratio $a$ (direct if $a>0$, ratio $|a|$ with a half-turn if $a<0$), centre $\omega = \dfrac{b}{1-a}$.
  \item If $|a| = 1$, $a \neq 1$: \textbf{rotation} of angle $\arg(a)$, centre $\omega = \dfrac{b}{1-a}$.
  \item Otherwise ($|a| \neq 1$, $a \notin \mathbb{R}$): \textbf{spiral similarity}, centre $\omega = \dfrac{b}{1-a}$, ratio $|a|$, angle $\arg(a)$.
  \end{itemize}
\item \textbf{Indirect case ($z' = a\bar{z} + b$):}
  \begin{itemize}
  \item If $|a| = 1$: \textbf{indirect isometry} — reflection (if fixed points exist, i.e. line of fixed points) or glide reflection (no fixed points).
  \item If $|a| \neq 1$: \textbf{indirect similarity} (homothety followed by reflection, or spiral similarity with reflection).
  \end{itemize}
\end{enumerate}
\end{methode}

\begin{exemplebox}[Four quick identifications — direct similarities]
\begin{itemize}
\item $z' = z + (2 - 3i)$: $a = 1$, translation by $\vec{u}(2, -3)$.
\item $z' = 3z - 4 + 2i$: $a = 3 \in \mathbb{R}$, homothety of ratio $3$, centre $\omega = \dfrac{-4+2i}{1-3} = 2 - i$.
\item $z' = iz + 1$: $|a| = 1$, $\arg a = \tfrac{\pi}{2}$, rotation through $90^\circ$ about $\omega = \dfrac{1}{1-i} = \dfrac{1+i}{2}$.
\item $z' = (1+i)z$: $|a| = \sqrt{2}$, $\arg a = \tfrac{\pi}{4}$, spiral similarity centred at $O$, ratio $\sqrt{2}$, angle $45^\circ$.
\end{itemize}
\end{exemplebox}

\begin{exemplebox}[Indirect similarities]
\begin{itemize}
\item $z' = \bar{z}$: reflection in the real axis (fixed points: real axis).
\item $z' = \bar{z} + 2$: glide reflection (reflection in real axis + translation by $2$ parallel to axis); no fixed points.
\item $z' = 2\bar{z}$: indirect similarity, ratio $2$, centre $O$ (homothety $2$ + reflection in real axis).
\end{itemize}
\end{exemplebox}

\begin{exoresolu}[Full identification, GCE style]
A transformation $S$ of the plane is given by $z' = (1 - i\sqrt{3})\,z + 2i\sqrt{3}$. Determine the nature of $S$ and its characteristic elements.
\tcblower
\textbf{Solution.} Here $a = 1 - i\sqrt{3}$ and $b = 2i\sqrt{3}$.
\[ |a| = \sqrt{1 + 3} = 2, \qquad \arg(a) = -\arctan\!\left(\frac{\sqrt{3}}{1}\right) = -\frac{\pi}{3} \pmod{2\pi}. \]
Since $|a| = 2 \neq 1$ and $a \notin \mathbb{R}$, $S$ is a \textbf{spiral similarity} (direct similarity which is neither a homothety nor an isometry). Its centre is
\[ \omega = \frac{b}{1-a} = \frac{2i\sqrt{3}}{1 - (1 - i\sqrt{3})} = \frac{2i\sqrt{3}}{i\sqrt{3}} = 2. \]
\textbf{Conclusion:} $S$ is the direct similarity of centre $\Omega(2, 0)$, ratio $k = 2$ and angle $\theta = -\dfrac{\pi}{3}$; equivalently $S = h(\Omega, 2) \circ r\!\left(\Omega, -\tfrac{\pi}{3}\right)$.
\end{exoresolu}

\begin{popfigure}
\begin{tikzpicture}[scale=1.1, >=stealth, node distance=1.2cm and 1.5cm]
\node[draw=popdark, fill=popblueL, rounded corners, align=center, font=\small] (start) {$z' = az + b$,\ $a \neq 0$};
\node[draw=popblue, rounded corners, font=\small] (q1) [below left=of start] {$a = 1$?};
\node[draw=popblue, rounded corners, font=\small] (q2) [below=of start] {$a \in \mathbb{R}$?};
\node[draw=popblue, rounded corners, font=\small] (q3) [below right=of start] {$|a| = 1$?};
\node[draw=popgreen, fill=popgreenL, rounded corners, font=\small, align=center] (t1) [below=of q1] {Translation\\ vector $b$};
\node[draw=popgreen, fill=popgreenL, rounded corners, font=\small, align=center] (t2) [below=of q2] {Homothety\\ ratio $a$};
\node[draw=popgreen, fill=popgreenL, rounded corners, font=\small, align=center] (t3) [below=of q3] {Rotation\\ angle $\arg a$};
\node[draw=poporange, fill=poppink!30, rounded corners, font=\small, align=center] (t4) [below right=of q3] {Spiral similarity\\ ratio $|a|$, angle $\arg a$};
\node[draw=popmuted, rounded corners, font=\small, align=center] (c) [below=1.5cm of t2] {Centre $\omega = \dfrac{b}{1-a}$\\ for homothety, rotation,\\ spiral similarity};
\draw[->, popline] (start) -- (q1);
\draw[->, popline] (start) -- (q2);
\draw[->, popline] (start) -- (q3);
\draw[->, popline] (q1) -- (t1) node[midway, left, font=\scriptsize] {yes};
\draw[->, popline] (q2) -- (t2) node[midway, left, font=\scriptsize] {yes};
\draw[->, popline] (q3) -- (t3) node[midway, left, font=\scriptsize] {yes};
\draw[->, popline] (q1) -- (q2) node[midway, above, font=\scriptsize] {no};
\draw[->, popline] (q2) -- (q3) node[midway, above, font=\scriptsize] {no};
\draw[->, popline] (q3) -- (t4) node[midway, above right, font=\scriptsize] {no};
\draw[->, popline] (t2) -- (c);
\draw[->, popline] (t3) -- (c);
\draw[->, popline] (t4) -- (c);
\end{tikzpicture}
\legende{Decision chart for direct similarities ($z' = az + b$). For indirect similarities ($z' = a\bar{z} + b$): if $|a|=1$ it is a reflection or glide reflection; otherwise an indirect similarity with ratio $|a|$.}
\end{popfigure}

\courssub{Composition of Similarities via Complex Expressions}

Composing similarities is now pure algebra. If $S_1: z' = a_1 z + b_1$ and $S_2: z' = a_2 z + b_2$ (both direct), then
\[ S_2 \circ S_1:\ z \mapsto a_2(a_1 z + b_1) + b_2 = a_1 a_2\, z + (a_2 b_1 + b_2). \]

\begin{propriete}[Ratio and angle of a composite — direct similarities]
The composite of two direct similarities is a direct similarity with
\[ \textbf{ratio} = k_1 k_2 \quad (|a_1 a_2| = |a_1||a_2|), \qquad \textbf{angle} = \theta_1 + \theta_2 \quad (\arg(a_1a_2) = \arg a_1 + \arg a_2). \]
In particular the set of direct similarities is a \emph{group} under composition; the inverse of $z' = az + b$ is $z' = \dfrac{1}{a}z - \dfrac{b}{a}$, a similarity of ratio $\dfrac{1}{|a|}$ and angle $-\arg a$.
\end{propriete}

\begin{propriete}[Composition involving indirect similarities]
\begin{itemize}
\item Direct $\circ$ Indirect = Indirect; Indirect $\circ$ Direct = Indirect.
\item Indirect $\circ$ Indirect = Direct.
\item In all cases, the ratio of the composite is the product of the ratios ($|a_1||a_2|$).
\item For the angle: direct similarities add their angles; an indirect similarity contributes $-\arg(a)$ (or equivalently, the oriented angle is reversed).
\end{itemize}
\end{propriete}

\begin{exemplebox}[Composing a rotation and a homothety]
Let $r: z' = iz$ (rotation about $O$ through $\tfrac{\pi}{2}$) and $h: z' = 2z - 3$ (homothety of ratio $2$). Then
\[ h \circ r:\ z \mapsto 2(iz) - 3 = 2iz - 3, \]
a spiral similarity of ratio $2$, angle $\tfrac{\pi}{2}$, centre $\omega = \dfrac{-3}{1 - 2i} = \dfrac{-3(1+2i)}{5} = -\dfrac{3}{5} - \dfrac{6}{5}i$. Note that $r \circ h: z \mapsto i(2z - 3) = 2iz - 3i$ has the \emph{same} ratio and angle but a \emph{different} centre: composition is not commutative.
\end{exemplebox}

\courssub{Geometric Applications: Triangles and Circles}

\begin{propriete}[Conservation properties of similarities]
Every similarity of ratio $k$:
\begin{itemize}
\item preserves angles (direct ones preserve oriented angles, indirect ones reverse them);
\item maps three collinear points to three collinear points and preserves ratios of distances along a line;
\item maps a line to a line and a circle to a circle, the radius being multiplied by $k$;
\item maps a triangle to a \emph{similar} triangle (same shape, sides multiplied by $k$).
\end{itemize}
\end{propriete}

\textbf{Why circles map to circles.} A circle of centre $\Omega(\omega)$ and radius $R$ is the set $\{z : |z - \omega| = R\}$. Under $z' = az + b$, writing $z = \dfrac{z' - b}{a}$, the condition becomes $\left|\dfrac{z' - b}{a} - \omega\right| = R$, i.e.\ $|z' - (a\omega + b)| = |a|R$: a circle of centre $\Omega'(a\omega + b)$ and radius $kR$. The same holds for $z' = a\bar{z} + b$ because $|\bar{z} - \bar{\omega}| = |z - \omega|$. \hfill$\blacksquare$

\begin{exemplebox}[Similar triangles via the complex ratio]
Triangles $ABC$ and $A'B'C'$ (affixes $z_A, \dots, z_{C'}$) are \textbf{directly similar} (same orientation) if and only if
\[ \frac{z_{B'} - z_{A'}}{z_{C'} - z_{A'}} = \frac{z_B - z_A}{z_C - z_A}. \]
Indeed this single complex equality encodes both the ratio of sides (moduli) and the angle at $A$ (arguments). This is the analytic version of the classical criteria (SSS, SAS, AA) for similar triangles.
\end{exemplebox}

\courssub{Spiral Similarities in Problem-Solving}

A spiral similarity is determined by \emph{one pair of corresponding points} together with its centre, or equivalently by two segments: given segments $AB$ and $A'B'$, there is a unique direct similarity sending $A \mapsto A'$, $B \mapsto B'$, namely $z' = az + b$ with $a = \dfrac{z_{B'} - z_{A'}}{z_B - z_A}$ and $b = z_{A'} - a z_A$.

\begin{methode}[Proving collinearity, concyclicity or equal angles using similarities]
\begin{enumerate}
\item Identify (or construct) a spiral similarity $S: z' = az + b$ sending a known configuration to the configuration under study.
\item \textbf{Collinearity:} to prove $P, Q, R$ collinear, show that $P, Q$ are images under $S$ of two points whose line must contain the preimage of $R$; or show $\dfrac{z_R - z_P}{z_Q - z_P} \in \mathbb{R}$.
\item \textbf{Equal angles:} use $\arg\!\left(\dfrac{z_C - z_A}{z_B - z_A}\right)$, which is invariant under a direct similarity.
\item \textbf{Concyclicity:} four points are concyclic (or collinear) iff their cross-ratio $\dfrac{(z_A - z_C)(z_B - z_D)}{(z_A - z_D)(z_B - z_C)}$ is real; similarities preserve cross-ratios, hence concyclicity.
\end{enumerate}
\end{methode}

\begin{exoresolu}[The classic two-triangle configuration]
Triangles $OAB$ and $OA'B'$ are directly similar at the common vertex $O$, with $A'$, $B'$ the images of $A$, $B$ under the spiral similarity $S$ of centre $O$, ratio $k$, angle $\theta$. Show that the lines $AA'$ and $BB'$ make the same angle with $OA$ and $OB$ respectively, and that $\dfrac{AA'}{BB'} = \dfrac{OA}{OB}$.
\tcblower
\textbf{Solution.} Place $O$ at the origin; then $S: z' = k e^{i\theta} z$, so $z_{A'} = a z_A$, $z_{B'} = a z_B$ with $a = ke^{i\theta}$. Then
\[ z_{A'} - z_A = (a - 1) z_A, \qquad z_{B'} - z_B = (a - 1) z_B. \]
Hence $\dfrac{z_{A'} - z_A}{z_{B'} - z_B} = \dfrac{z_A}{z_B}$, a \emph{constant} complex number. Taking moduli: $\dfrac{AA'}{BB'} = \dfrac{OA}{OB}$. Taking arguments: the directed angle from $BB'$ to $AA'$ equals the angle from $OB$ to $OA$, i.e.\ the two segments $AA'$ and $BB'$ are inclined to $OA$, $OB$ in exactly the same way. In particular, if $O, A, B$ are collinear then $A, A', B, B'$ lie in a configuration where $AA' \parallel BB'$ — a standard route to proving parallelism and collinearity in GCE problems.
\end{exoresolu}

\begin{popfigure}
\begin{tikzpicture}[scale=1.6, >=stealth]
\coordinate (O) at (0,0);
\coordinate (A) at (2.4,0.3);
\coordinate (B) at (1.5,1.5);
\coordinate (Ap) at (1.5,1.9);
\coordinate (Bp) at (-0.4,2.3);
\draw[popblue, thick] (O)--(A)--(B)--cycle;
\draw[poppink, thick] (O)--(Ap)--(Bp)--cycle;
\foreach \p/\n/\pos in {O/O/below left, A/A/below right, B/B/above right, Ap/A'/above right, Bp/B'/above left}
  \fill[popdark] (\p) circle (0.03) node[\pos, font=\small] {$\n$};
\draw[->, poporange, dashed, thick] (A) -- (Ap);
\draw[->, poporange, dashed, thick] (B) -- (Bp);
\draw[popgreen, ->] (0.85,0.11) arc (7:52:0.86);
\node[popgreen, font=\small] at (1.15,0.62) {$\theta$};
\node[popdark, font=\small, align=center] at (3.1,1.9) {$z' = ke^{i\theta}z$\\ $OA' = k\,OA$};
\end{tikzpicture}
\legende{Two directly similar triangles sharing the vertex $O$: the spiral similarity $z' = ke^{i\theta}z$ sends $A \mapsto A'$ and $B \mapsto B'$. The dashed segments $AA'$, $BB'$ are in a constant ratio and constant relative inclination.}
\end{popfigure}

\begin{savaistu}[Spiral similarities and the ``pivot theorem'']
Spiral similarities are the secret engine behind many olympiad and GCE geometry facts: the centre of the spiral similarity sending segment $AB$ to $A'B'$ lies at the intersection of the circles through $A, A'$ and the intersection point of lines $AA'$, $BB'$. Napoleon's theorem, the Miquel point of a complete quadrilateral, and the properties of the nine-point circle can all be proved elegantly this way.
\end{savaistu}

\courssub{Synthesis Problems, Assessment and Remediation}

\begin{exoresolu}[GCE Paper 2 style — mixing homothety, rotation and complex numbers]
The plane is referred to an orthonormal frame $(O; \vec{u}, \vec{v})$. Let $T$ be the transformation with complex expression $z' = \left(\dfrac{1}{2} + \dfrac{\sqrt{3}}{2}i\right) z + 1 - \sqrt{3}i$.
\begin{enumerate}
\item[(a)] Show that $T$ is a rotation and determine its centre and angle.
\item[(b)] Let $h$ be the homothety of centre $O$ and ratio $2$. Write the complex expression of $S = h \circ T$.
\item[(c)] Deduce the nature and characteristic elements of $S$.
\item[(d)] The circle $\mathcal{C}$ of centre $O$ and radius $1$ is transformed by $S$ into $\mathcal{C}'$. Give the centre and radius of $\mathcal{C}'$.
\end{enumerate}
\tcblower
\textbf{Solution.}
(a) $a = \dfrac{1}{2} + \dfrac{\sqrt{3}}{2}i = e^{i\pi/3}$, so $|a| = 1$, $a \neq 1$: $T$ is a \textbf{rotation} of angle $\dfrac{\pi}{3}$. Centre: $\omega = \dfrac{b}{1-a} = \dfrac{1 - \sqrt{3}i}{1 - \tfrac12 - \tfrac{\sqrt3}{2}i} = \dfrac{1 - \sqrt{3}i}{\tfrac12 - \tfrac{\sqrt3}{2}i} = \dfrac{1 - \sqrt3 i}{e^{-i\pi/3}}$. Since $1 - \sqrt3 i = 2e^{-i\pi/3}$, we get $\omega = 2$. Centre $\Omega(2, 0)$.

(b) $h: z' = 2z$. $S = h \circ T:\ z \mapsto 2\left(e^{i\pi/3} z + 1 - \sqrt3 i\right) = (1 + \sqrt3 i)\,z + 2 - 2\sqrt3 i$.

(c) $|1 + \sqrt3 i| = 2$, $\arg(1 + \sqrt3 i) = \dfrac{\pi}{3}$. Since $|a| = 2 \neq 1$: $S$ is a \textbf{spiral similarity} of ratio $2$, angle $\dfrac{\pi}{3}$, centre $\omega_S = \dfrac{2 - 2\sqrt3 i}{1 - 1 - \sqrt3 i} = \dfrac{2 - 2\sqrt3 i}{-\sqrt3 i} = \dfrac{(2 - 2\sqrt3 i)\,i}{\sqrt3} = 2 + \dfrac{2i}{\sqrt3} = 2 + \dfrac{2\sqrt3}{3}i$.

(d) $\mathcal{C}'$ is a circle (similarities map circles to circles) of radius $2 \times 1 = 2$, centred at $S(O)$: $z_{O'} = 2 - 2\sqrt3 i$, i.e.\ $O'(2, -2\sqrt3)$.
\end{exoresolu}

\begin{attention}[Exam tip — GCE Paper 2]
In Paper 2, part (a) usually asks for the \textbf{complex form} and part (b) for the \textbf{geometric elements}. Answer both explicitly and completely: state the \emph{nature} (translation / homothety / rotation / reflection / glide reflection / spiral similarity), the \emph{centre} (affix $\omega = \dfrac{b}{1-a}$ for direct, or fixed point analysis for indirect), the \emph{ratio} $|a|$ and the \emph{angle} $\arg(a)$. A common loss of marks: giving the ratio as a complex number, or forgetting that a translation ($a = 1$) has no centre.
\end{attention}

\begin{exercice}[Assessment]
\begin{enumerate}
\item Identify each transformation and give its characteristic elements:
\begin{itemize}
\item $z' = -2z + 6$; \quad $z' = e^{-i\pi/4} z + i$; \quad $z' = (3 - 3i)z + 4i$; \quad $z' = \bar z + 2$.
\end{itemize}
\item $S_1: z' = iz + 2$ and $S_2: z' = 2z - i$. Find the complex expressions of $S_2 \circ S_1$ and $S_1 \circ S_2$, and the nature and elements of each composite. Comment.
\item $A(1)$, $B(2 + i)$, $C(i)$ are three points and $S: z' = (1 + i)z - i$. Find the images $A', B', C'$; show that $ABC$ and $A'B'C'$ are similar and give the ratio.
\item The circle $\mathcal{C}: |z - 1 + i| = 3$ is transformed by $z' = 2i\bar z + 1$. Find the centre and radius of the image circle. Is the similarity direct or indirect?
\end{enumerate}
\end{exercice}

\begin{corrige}[Exercise 1]
\textbf{1.} (i) $a = -2 \in \mathbb{R}$: homothety of ratio $-2$ (i.e.\ ratio $2$ with a half-turn), centre $\omega = \dfrac{6}{1-(-2)} = 2$. \\
(ii) $a = e^{-i\pi/4}$, $|a|=1$, $a\neq 1$: rotation of angle $-\dfrac{\pi}{4}$. Centre $\omega = \dfrac{i}{1 - e^{-i\pi/4}}$. Compute exactly:
$1 - e^{-i\pi/4} = 1 - \frac{\sqrt2}{2} + i\frac{\sqrt2}{2}$.
$\omega = \dfrac{i\left(1 - \frac{\sqrt2}{2} - i\frac{\sqrt2}{2}\right)}{\left(1-\frac{\sqrt2}{2}\right)^2 + \left(\frac{\sqrt2}{2}\right)^2} = \dfrac{\frac{\sqrt2}{2} + i\left(1-\frac{\sqrt2}{2}\right)}{2-\sqrt2} = \dfrac{1+\sqrt2}{2} + \frac{i}{2}$.
So centre $\Omega\left(\frac{1+\sqrt2}{2}, \frac12\right)$. \\
(iii) $a = 3-3i = 3\sqrt2 e^{-i\pi/4}$: $|a| = 3\sqrt2$, $\arg a = -\frac{\pi}{4}$. Spiral similarity, centre $\omega = \dfrac{4i}{1 - 3 + 3i} = \dfrac{4i}{-2 + 3i} = \dfrac{4i(-2-3i)}{13} = \dfrac{12 - 8i}{13}$, ratio $3\sqrt2$, angle $-\frac{\pi}{4}$. \\
(iv) $z' = \bar z + 2$: indirect isometry ($|a|=1$). Fixed points satisfy $z = \bar z + 2 \Rightarrow z - \bar z = 2 \Rightarrow 2i\,\mathrm{Im}(z) = 2$, impossible for real $\mathrm{Im}(z)$. Hence no fixed points: it is a \textbf{glide reflection} (reflection in the real axis followed by translation by $2$ parallel to that axis).

\textbf{2.} $S_2 \circ S_1: z \mapsto 2(iz + 2) - i = 2iz + 4 - i$: spiral similarity, ratio $2$, angle $\tfrac{\pi}{2}$, centre $\omega = \dfrac{4 - i}{1 - 2i} = \dfrac{(4-i)(1+2i)}{5} = \dfrac{6 + 7i}{5}$. \\
$S_1 \circ S_2: z \mapsto i(2z - i) + 2 = 2iz + 3$: same ratio and angle, centre $\omega = \dfrac{3}{1 - 2i} = \dfrac{3 + 6i}{5}$. \\
The composites differ: composition of similarities is \emph{not commutative} in general.

\textbf{3.} $z_{A'} = (1+i)(1) - i = 1$; $z_{B'} = (1+i)(2+i) - i = 1 + 2i$; $z_{C'} = (1+i)(i) - i = -1$. \\
Check: $\dfrac{z_{B'} - z_{A'}}{z_{C'} - z_{A'}} = \dfrac{2i}{-2} = -i$ and $\dfrac{z_B - z_A}{z_C - z_A} = \dfrac{1 + i}{-1 + i} = \dfrac{(1+i)(-1-i)}{2} = -i$. Equal, so the triangles are directly similar; ratio $|1 + i| = \sqrt2$.

\textbf{4.} $z' = 2i\bar z + 1$ is an \textbf{indirect} similarity (presence of $\bar z$) of ratio $|2i| = 2$. The centre $1 - i$ maps to $2i\overline{(1 - i)} + 1 = 2i(1 + i) + 1 = -1 + 2i$. Image circle: centre $(-1, 2)$, radius $2 \times 3 = 6$.
\end{corrige}

\begin{aretenir}[Key points of the section]
\begin{itemize}
\item Direct similarity: $z' = az + b$ ($a \neq 0$); indirect: $z' = a\bar z + b$. Ratio $= |a|$, angle $= \arg(a)$ (direct) or orientation-reversing (indirect).
\item If $a \neq 1$, unique fixed point (centre) $\omega = \dfrac{b}{1 - a}$ for direct similarities; for indirect, solve $z = a\bar{z}+b$.
\item $a = 1$: translation; $a \in \mathbb{R}$: homothety; $|a| = 1$: rotation; otherwise: spiral similarity (direct). Indirect: $|a|=1$ gives reflection/glide reflection; $|a|\neq1$ gives indirect similarity.
\item Composition multiplies ratios and adds angles (direct) or subtracts (indirect): $a_{\text{composite}} = a_1 a_2$ for direct$\circ$direct.
\item Similarities preserve angles, collinearity and concyclicity; circles map to circles with radius multiplied by $k$; triangles map to similar triangles.
\item In GCE answers, always state nature, centre, ratio and angle explicitly.
\end{itemize}
\end{aretenir}

\newpage

\coursec{Integration activity}

\courssub{Aims of the activity}

This integration activity mobilises, in a single realistic project, all the skills developed in this chapter on \cle{plane transformations}. By the end of the activity, you should be able to:

\begin{itemize}
\item \textbf{Compose two reflections} across intersecting axes and identify the resulting rotation, giving its centre and angle;
\item \textbf{Construct and compute} the images of a figure under homotheties of given ratios, analytically and geometrically;
\item \textbf{Write the complex expression} $z' = az + b$ of a direct plane similarity (spiral similarity) and determine its characteristic elements: centre, ratio and angle;
\item \textbf{Use the area-scaling property} $A' = k^{2}A$ of a similarity of ratio $k$ to solve a concrete estimation problem;
\item \textbf{Justify} each transformation rigorously by its characteristic elements, in the style required at the GCE Advanced Level examination;
\item Work collaboratively, present constructions on grid paper, and defend your reasoning orally.
\end{itemize}

\begin{aretenir}[Skills mobilised]
Composition of isometries (reflections, rotations, translations), glide reflections, homotheties and their composition, complex representation of direct similarities $z' = az + b$, ratio and angle of a similarity, effect of a similarity on lengths ($\times |a|$) and on areas ($\times |a|^{2}$).
\end{aretenir}

\courssub{Situation}

\textbf{Designing a Logo for a Yaoundé Printing Press.}\par
\medskip
\emph{Mvogo \& Sons Press}, a small printing company located near the Mokolo market in Yaoundé, wants a new logo. The designer proposes a pattern built from a single basic triangular motif, repeated by rotation around a central point $O$, and then produced in three sizes: a small version for \textbf{business cards}, a medium version for \textbf{posters}, and a large version for \textbf{billboards}.

The press works with a graphic designer who uses a coordinate plane. The unit on both axes is $1\ \mathrm{cm}$ on the business-card artwork. The basic motif is the triangle $T = ABC$ with
\[ A(1\,;\,0), \qquad B(3\,;\,0), \qquad C(2\,;\,1.5), \]
and the centre of the pattern is the origin $O(0\,;\,0)$.

To generate the repeated pattern, the designer reflects the motif successively across two axes passing through $O$: the $x$-axis $(\Delta_1)$ and the line $(\Delta_2)$ of equation $y = x$. The motif is first reflected across $(\Delta_1)$, and the image is then reflected across $(\Delta_2)$. The press owner wants to understand which single transformation achieves the same result, so that the pattern can be repeated around $O$.

For the three sizes, the artwork is enlarged from $O$ using homotheties of centre $O$ and ratios
\[ k_1 = 1 \quad (\text{business card}), \qquad k_2 = 2.5 \quad (\text{poster}), \qquad k_3 = 6 \quad (\text{billboard}). \]

Finally, for an advertising animation, the designer wants to map the small motif $ABC$ directly onto a \emph{rotated and enlarged} copy $A''B''C''$, where $A''$ and $B''$ are the images of $A$ and $B$ under the direct similarity $s$ of complex expression
\[ z' = az + b, \qquad a = 1 + i\sqrt{3}, \qquad b = 2. \]
The complex plane is referred to the orthonormal frame $(O\,;\,\vec{u},\vec{v})$.

The printing cost of the coloured motif is estimated at $2500$ FCFA per square decimetre of inked surface. The owner asks: \emph{what will the ink for one motif cost on each of the three supports?}

\begin{popfigure}
\begin{center}
\begin{tikzpicture}[scale=1.1]
\draw[popline, very thin, step=0.5] (-2.2,-2.2) grid (3.6,3.6);
\draw[->, popdark, thick] (-2.2,0) -- (3.6,0) node[below left] {$x$};
\draw[->, popdark, thick] (0,-2.2) -- (0,3.6) node[below left] {$y$};
\draw[popmuted, dashed] (-2.2,-2.2) -- (3.4,3.4) node[below right, pos=0.95] {$(\Delta_2):\ y=x$};
\node[below left, popdark] at (0,0) {$O$};
% motif T
\draw[popblue, very thick, fill=popblueL] (1,0) -- (3,0) -- (2,1.5) -- cycle;
\node[popblue, below] at (1,0) {$A$};
\node[popblue, below] at (3,0) {$B$};
\node[popblue, above] at (2,1.5) {$C$};
% image under r1
\draw[poporange, very thick, fill=popgold!25] (1,0) -- (3,0) -- (2,-1.5) -- cycle;
\node[poporange, above] at (1,0) {$A_1$};
\node[poporange, above] at (3,0) {$B_1$};
\node[poporange, below] at (2,-1.5) {$C_1$};
% image under r2 o r1
\draw[poppurple, very thick, fill=poppink!30] (0,1) -- (0,3) -- (-1.5,2) -- cycle;
\node[poppurple, right] at (0,1) {$A'$};
\node[poppurple, right] at (0,3) {$B'$};
\node[poppurple, left] at (-1.5,2) {$C'$};
% angle between axes
\draw[popgreen, thick, ->] (1.6,0) arc (0:45:1.6);
\node[popgreen] at (1.9,0.55) {$45^{\circ}$};
\end{tikzpicture}
\end{center}
\legende{The motif $T=ABC$ (blue), its image under $r_1$ (orange) and its image under $f=r_2\circ r_1$ (purple): a quarter-turn about $O$.}
\end{popfigure}

\courssub{Tasks}

Work in groups of three or four. Present all constructions on grid paper and justify every transformation by its \textbf{characteristic elements}.

\begin{enumerate}
\item \textbf{Placing the motif.} Draw the coordinate plane (unit $1\ \mathrm{cm}$), plot the triangle $T = ABC$, and compute its area $\mathcal{A}$ in $\mathrm{cm}^2$.
\item \textbf{Composing two reflections.} Let $r_1$ be the reflection across $(\Delta_1)$ (the $x$-axis) and $r_2$ the reflection across $(\Delta_2)$ (the line $y = x$).
\begin{itemize}
\item Determine the coordinates of $A_1 = r_1(A)$, $B_1 = r_1(B)$, $C_1 = r_1(C)$, then those of $A' = r_2(A_1)$, $B' = r_2(B_1)$, $C' = r_2(C_1)$.
\item State the theorem on the composition of two reflections across \textbf{intersecting axes}. Identify the single transformation $f = r_2 \circ r_1$, giving its centre and its angle (measure the angle from $(\Delta_1)$ to $(\Delta_2)$).
\item Verify numerically that $A'$ is indeed the image of $A$ under this rotation.
\item Deduce the transformation that generates the whole repeated pattern around $O$, and state how many copies of the motif complete one full turn.
\end{itemize}
\item \textbf{The three sizes.} Let $h_k$ denote the homothety of centre $O$ and ratio $k$.
\begin{itemize}
\item Give the analytical expression of $h_k$, i.e. the coordinates of $M'(x'\,;\,y')$, image of $M(x\,;\,y)$.
\item Compute the coordinates of the vertices of the poster motif ($k_2 = 2.5$) and of the billboard motif ($k_3 = 6$).
\item State the property giving the image of a triangle under a homothety, and the effect on side lengths.
\end{itemize}
\item \textbf{The spiral similarity.} Consider $s:\ z' = (1 + i\sqrt{3})\,z + 2$.
\begin{itemize}
\item Compute the modulus and an argument of $a = 1 + i\sqrt{3}$. Deduce the ratio and the angle of $s$.
\item Show that $s$ has a unique fixed point $\Omega$, and compute its affix $\omega$. Conclude on the nature and the characteristic elements of $s$.
\item Compute the affixes of $A'' = s(A)$ and $B'' = s(B)$, then the distance $A''B''$. Verify that $A''B'' = |a| \cdot AB$.
\end{itemize}
\item \textbf{Ink cost estimation.}
\begin{itemize}
\item State the area-scaling property of a similarity of ratio $k$.
\item Compute the area of one motif on each support, in $\mathrm{cm}^2$ then in $\mathrm{dm}^2$ ($1\ \mathrm{dm}^2 = 100\ \mathrm{cm}^2$).
\item Estimate the ink cost of one motif on each support at $2500$ FCFA per $\mathrm{dm}^2$.
\end{itemize}
\item \textbf{Synthesis.} Each group presents its grid-paper constructions and completes a summary table: transformation, definition used, characteristic elements, image of the motif, effect on lengths and areas.
\end{enumerate}

\courssub{Model answer}

\textbf{1. Placing the motif.}\par
The triangle $ABC$ has base $AB$ on the $x$-axis with $AB = 3 - 1 = 2$ cm, and the height from $C$ to the line $(AB)$ is the ordinate of $C$, namely $1.5$ cm. Hence
\[ \mathcal{A} = \frac{1}{2} \times AB \times 1.5 = \frac{1}{2} \times 2 \times 1.5 = 1.5\ \mathrm{cm}^2. \]

\textbf{2. Composing two reflections.}\par
The reflection $r_1$ across the $x$-axis maps $(x\,;\,y) \mapsto (x\,;\,-y)$:
\[ A_1(1\,;\,0), \qquad B_1(3\,;\,0), \qquad C_1(2\,;\,-1.5). \]
The reflection $r_2$ across the line $y = x$ maps $(x\,;\,y) \mapsto (y\,;\,x)$:
\[ A'(0\,;\,1), \qquad B'(0\,;\,3), \qquad C'(-1.5\,;\,2). \]
\textbf{Theorem (composition of two reflections, intersecting axes):} if $(\Delta_1)$ and $(\Delta_2)$ intersect at $O$, then $r_2 \circ r_1$ is the \cle{rotation} of centre $O$ whose angle is twice the directed angle from the first axis to the second:
\[ r_2 \circ r_1 = \mathrm{Rot}\big(O\,;\, 2\,\widehat{(\Delta_1, \Delta_2)}\big). \]
Here $\widehat{(\Delta_1,\Delta_2)} = 45^{\circ}$, so $f = r_2 \circ r_1$ is the rotation of centre $O$ and angle $2 \times 45^{\circ} = 90^{\circ}$.\par
\emph{Verification:} the rotation of centre $O$ and angle $90^{\circ}$ maps $(x\,;\,y) \mapsto (-y\,;\,x)$. Applied to $A(1\,;\,0)$ it gives $(0\,;\,1) = A'$; applied to $B(3\,;\,0)$ it gives $(0\,;\,3) = B'$; applied to $C(2\,;\,1.5)$ it gives $(-1.5\,;\,2) = C'$. $\blacksquare$\par
The pattern is therefore generated by successive rotations of angle $90^{\circ}$ about $O$; since $\dfrac{360^{\circ}}{90^{\circ}} = 4$, exactly \textbf{four copies} of the motif complete one full turn.

\textbf{3. The three sizes.}\par
The homothety $h_k$ of centre $O$ and ratio $k$ is defined by $\overrightarrow{OM'} = k\,\overrightarrow{OM}$; its analytical expression is therefore
\[ M(x\,;\,y) \longmapsto M'(kx\,;\,ky). \]
For the \textbf{poster} ($k_2 = 2.5$): $A_2(2.5\,;\,0)$, $B_2(7.5\,;\,0)$, $C_2(5\,;\,3.75)$.\par
For the \textbf{billboard} ($k_3 = 6$): $A_3(6\,;\,0)$, $B_3(18\,;\,0)$, $C_3(12\,;\,9)$.\par
\textbf{Property:} a homothety maps a triangle onto a triangle whose corresponding sides are parallel, and it multiplies all lengths by $|k|$ (here $k > 0$, so lengths are multiplied by $k$).

\textbf{4. The spiral similarity.}\par
For $a = 1 + i\sqrt{3}$: $|a| = \sqrt{1^{2} + (\sqrt{3})^{2}} = \sqrt{4} = 2$, and writing $a = 2\left(\dfrac{1}{2} + i\dfrac{\sqrt{3}}{2}\right)$ gives $\arg(a) = \dfrac{\pi}{3}\ [2\pi]$. The ratio of $s$ is therefore $k = |a| = 2$ and its angle is $\theta = \arg(a) = \dfrac{\pi}{3}$.\par
\textbf{Fixed point:} a point of affix $z$ is fixed by $s$ when $z = az + b$, i.e. $(1-a)z = b$. Since $a = 1 + i\sqrt{3} \neq 1$, there is a \textbf{unique} fixed point $\Omega$, of affix
\[ \omega = \frac{b}{1 - a} = \frac{2}{1 - (1 + i\sqrt{3})} = \frac{2}{-i\sqrt{3}} = \frac{2i}{\sqrt{3}} = \frac{2i\sqrt{3}}{3}. \]
Thus $s$ is the \cle{direct similarity} (spiral similarity) of centre $\Omega\left(0\,;\,\dfrac{2\sqrt{3}}{3}\right)$, ratio $2$ and angle $\dfrac{\pi}{3}$: it is the composition, in either order, of the homothety $h\!\left(\Omega\,;\,2\right)$ and the rotation $\mathrm{Rot}\!\left(\Omega\,;\,\dfrac{\pi}{3}\right)$.\par
\textbf{Images:} $z_A = 1$, so $z_{A''} = (1 + i\sqrt{3})(1) + 2 = 3 + i\sqrt{3}$; $z_B = 3$, so $z_{B''} = 3(1 + i\sqrt{3}) + 2 = 5 + 3i\sqrt{3}$. Then
\[ A''B'' = |z_{B''} - z_{A''}| = |2 + 2i\sqrt{3}| = \sqrt{2^{2} + (2\sqrt{3})^{2}} = \sqrt{4 + 12} = 4 = 2 \times AB, \]
which confirms that a direct similarity of ratio $|a|$ multiplies all distances by $|a|$. $\blacksquare$

\textbf{5. Ink cost estimation.}\par
\textbf{Property:} a similarity of ratio $k$ multiplies areas by $k^{2}$. With $\mathcal{A} = 1.5\ \mathrm{cm}^2$:
\begin{center}
\begin{tabularx}{\linewidth}{|l|X|X|X|}
\hline
\rowcolor{popblueL} \textbf{Support} & \textbf{Ratio} $k$ & \textbf{Area} $= 1.5k^{2}\ \mathrm{cm}^2$ & \textbf{Cost} (FCFA) \\
\hline
Business card & $1$ & $1.5\ \mathrm{cm}^2 = 0.015\ \mathrm{dm}^2$ & $0.015 \times 2500 = 37.5$ \\
\hline
Poster & $2.5$ & $9.375\ \mathrm{cm}^2 = 0.09375\ \mathrm{dm}^2$ & $\approx 234.4$ \\
\hline
Billboard & $6$ & $54\ \mathrm{cm}^2 = 0.54\ \mathrm{dm}^2$ & $0.54 \times 2500 = 1350$ \\
\hline
\end{tabularx}
\end{center}

\textbf{6. Synthesis.}\par
\begin{center}
\begin{tabularx}{\linewidth}{|l|X|X|}
\hline
\rowcolor{popblueL} \textbf{Transformation} & \textbf{Characteristic elements} & \textbf{Effect} \\
\hline
$r_2 \circ r_1$ & Rotation, centre $O$, angle $2 \times 45^{\circ} = 90^{\circ}$ & Preserves lengths and areas \\
\hline
$h(O\,;\,k)$ & Centre $O$, ratio $k \in \{1\,;\,2.5\,;\,6\}$ & Lengths $\times k$, areas $\times k^{2}$ \\
\hline
$s:\ z' = az + b$ & Centre $\Omega$, ratio $|a| = 2$, angle $\arg a = \frac{\pi}{3}$ & Lengths $\times 2$, areas $\times 4$ \\
\hline
\end{tabularx}
\end{center}

\begin{attention}[Common mistakes to avoid]
\begin{itemize}
\item The angle of the rotation $r_2 \circ r_1$ is \textbf{twice} the angle between the axes, measured from the \emph{first} axis to the \emph{second}: the order of composition matters ($r_2 \circ r_1 \neq r_1 \circ r_2$ in general).
\item The centre of $z' = az + b$ is \emph{not} the point of affix $b$: it is the fixed point of affix $\omega = \dfrac{b}{1-a}$.
\item Areas scale by $k^{2}$, not by $k$: the billboard motif is $36$ times, not $6$ times, the area of the card motif.
\item Do not forget unit conversions ($\mathrm{cm}^2 \to \mathrm{dm}^2$) before computing the cost.
\end{itemize}
\end{attention}

\begin{exercice}[Extension: a glide reflection for the logo]
The press owner now wants a motif mapped onto a shifted copy by a glide reflection. Let $t$ be the translation of vector $\vec{u}(4\,;\,0)$ and $r$ the reflection across the $x$-axis.
\begin{enumerate}
\item Compute the coordinates of the images of $A$, $B$, $C$ under $g = r \circ t$, then under $g' = t \circ r$.
\item What do you observe? Explain, using the fact that $\vec{u}$ is parallel to the axis of $r$.
\item Justify that $g$ is an \textbf{indirect isometry} and prove that it has \textbf{no fixed point}.
\end{enumerate}
\end{exercice}

\begin{corrige}[Extension: a glide reflection for the logo]
\textbf{1.} Under $t$: $(x\,;\,y) \mapsto (x+4\,;\,y)$; under $r$: $(x\,;\,y) \mapsto (x\,;\,-y)$.\par
For $g = r \circ t$ (translate first, then reflect): $(x\,;\,y) \mapsto (x+4\,;\,-y)$, giving
\[ A \mapsto (5\,;\,0), \qquad B \mapsto (7\,;\,0), \qquad C \mapsto (6\,;\,-1.5). \]
For $g' = t \circ r$ (reflect first, then translate): $(x\,;\,y) \mapsto (x+4\,;\,-y)$ as well, giving the same images.\par
\textbf{2.} We observe that $g = g'$: the translation and the reflection \textbf{commute}. This is because the vector $\vec{u}$ is parallel to the axis of reflection: reflecting across the $x$-axis does not change a horizontal displacement. The composite $g = r \circ t$ with $\vec{u}$ parallel to the axis is precisely a \cle{glide reflection}.\par
\textbf{3.} The reflection $r$ is an indirect isometry and $t$ is a direct isometry, so their composition is an \textbf{indirect isometry} (it reverses orientation). A fixed point would satisfy $(x+4\,;\,-y) = (x\,;\,y)$, i.e. $x + 4 = x$, which is impossible. Hence $g$ has \textbf{no fixed point}, as expected for a glide reflection with a non-zero translation vector.
\end{corrige}

\newpage

\coursec{Methods and worked examples}

\courssub{Composing two reflections (orthogonal symmetries)}

Recall that the \cle{orthogonal symmetry} (reflection) $S_d$ in a line $d$ sends a point $M$ to the point $M'$ such that $d$ is the perpendicular bisector of $[MM']$ (and fixes every point of $d$). What happens when we apply two reflections one after the other? The answer depends only on the relative position of the two axes, and it produces the two fundamental \emph{direct} isometries: translations and rotations.

\begin{propriete}[Composite of two reflections — proof not examinable]
Let $d_1$ and $d_2$ be two lines of the plane.
\begin{itemize}
\item If $d_1 \parallel d_2$, then $S_{d_2} \circ S_{d_1}$ is the \cle{translation} of vector $\vec{u} = 2\overrightarrow{H_1H_2}$, where $H_1 \in d_1$ and $H_2 \in d_2$ are the feet of a common perpendicular. In particular $\|\vec{u}\| = 2\,\mathrm{dist}(d_1, d_2)$.
\item If $d_1$ and $d_2$ intersect at $O$, then $S_{d_2} \circ S_{d_1}$ is the \cle{rotation} of centre $O$ and angle $2\theta$, where $\theta = (d_1, d_2)$ is the oriented angle from $d_1$ to $d_2$.
\end{itemize}
\end{propriete}

\begin{popfigure}
\begin{tikzpicture}[scale=0.9]
\draw[popblue, thick] (0,-0.4) -- (0,2.6) node[above] {$d_1$};
\draw[poporange, thick] (1.5,-0.4) -- (1.5,2.6) node[above] {$d_2$};
\fill (0.6,0.7) circle (1.5pt) node[left] {$M$};
\fill (2.4,0.7) circle (1.5pt) node[right] {$M_1$};
\fill (0.6,0.7) circle (0pt);
\fill (0.6+3.0-2.4+2.4,0.7) circle (0pt);
\fill (0.6,0.7) ++(0,0) circle (0pt);
\coordinate (M) at (0.6,0.7);
\coordinate (M1) at (2.4,0.7);
\coordinate (M2) at (0.6,0.7);
\fill (0.6,1.8) circle (1.5pt) node[left] {$N$};
\fill (2.4,1.8) circle (1.5pt) node[right] {$N_1$};
\fill (0.6,1.8) ++(3.0,0) circle (1.5pt) node[right] {$N'$};
\fill (0.6,0.7) ++(3.0,0) circle (1.5pt) node[right] {$M'$};
\draw[->, popdark, thick] (0.6,0.7) -- (3.6,0.7);
\draw[->, popdark, thick] (0.6,1.8) -- (3.6,1.8);
\draw[<->, popmuted] (0,-0.25) -- (1.5,-0.25) node[midway, below] {$a$};
\draw[->, popgreen, thick] (4.1,0.4) -- (7.1,0.4) node[midway, above] {$\vec{u}$, $\|\vec{u}\|=2a$};
\end{tikzpicture}
\legende{Two reflections in parallel axes $d_1, d_2$ give a translation of vector twice the gap between the axes.}
\end{popfigure}

\begin{methode}[Identifying the composite of two reflections]
To determine the nature of $S_{d_2} \circ S_{d_1}$ where $d_1, d_2$ are lines in the plane:
\begin{enumerate}
\item Draw or visualise the two axes $d_1$ and $d_2$. Decide whether they are \cle{parallel} or \cle{intersecting}.
\item \textbf{Parallel axes} (distance $a$ apart): the composite is a \cle{translation} of vector $\vec{u} = 2\overrightarrow{H_1H_2}$, where $H_1 \in d_1$, $H_2 \in d_2$ are feet of a common perpendicular. Hence $\|\vec{u}\| = 2a$, directed from $d_1$ towards $d_2$.
\item \textbf{Intersecting axes} at $O$, with angle $\theta$ from $d_1$ to $d_2$: the composite is the \cle{rotation} of centre $O$ and angle $2\theta$.
\item Remember the order: $S_{d_2} \circ S_{d_1}$ means $S_{d_1}$ is applied \emph{first}. Reversing the order reverses the translation vector or the rotation angle.
\item To verify, test the composite on one or two well-chosen points (e.g.\ a point on $d_2$ is fixed by $S_{d_2}$, so only $S_{d_1}$ moves it).
\end{enumerate}
\end{methode}

\begin{exoresolu}[Composite of reflections in two parallel lines]
In the plane with an orthonormal frame $(O;\vec{i},\vec{j})$, let $d_1$ be the line of equation $x = 1$ and $d_2$ the line of equation $x = 4$. Determine the nature and characteristic elements of $f = S_{d_2} \circ S_{d_1}$, and find the image of $A(2,3)$ under $f$.
\tcblower
\textbf{Solution.}

\textbf{Step 1 — Nature of the axes.} Both $d_1$ and $d_2$ are vertical lines, hence parallel. The distance between them is $a = 4 - 1 = 3$.

\textbf{Step 2 — Apply the theorem.} The composite of two reflections in parallel axes is a translation of vector $\vec{u} = 2\overrightarrow{H_1H_2}$, where $H_1, H_2$ are the intersections of a common perpendicular with $d_1, d_2$. Take the horizontal perpendicular: $H_1(1,0) \in d_1$, $H_2(4,0) \in d_2$. Then
\[ \vec{u} = 2\overrightarrow{H_1H_2} = 2(4-1)\,\vec{i} = 6\vec{i}. \]
So $f$ is the \textbf{translation of vector $\vec{u} = 6\vec{i}$}, i.e.\ $M(x,y) \mapsto M'(x+6, y)$.

\textbf{Step 3 — Direct verification.} For $M(x,y)$: $S_{d_1}$ sends $M$ to $M_1$ with the same ordinate and abscissa $x_1$ satisfying $\frac{x+x_1}{2} = 1$, so $x_1 = 2 - x$. Then $S_{d_2}$ sends $M_1$ to $M'$ with $\frac{x_1 + x'}{2} = 4$, giving $x' = 8 - x_1 = 8 - (2-x) = x + 6$. Confirmed: $f$ is the translation by $6\vec{i}$.

\textbf{Step 4 — Image of $A(2,3)$.} $A' = f(A) = (2+6,\, 3) = (8,3)$.
\end{exoresolu}

\begin{exoresolu}[Composite of reflections in two intersecting lines]
Let $d_1$ and $d_2$ be two lines through the point $O$, such that the oriented angle from $d_1$ to $d_2$ is $\dfrac{\pi}{6}$. Show that $f = S_{d_2} \circ S_{d_1}$ is a rotation, give its elements, and deduce the image of a point $M \neq O$ with $OM = 4$ cm when $M \in d_1$.
\tcblower
\textbf{Solution.}

\textbf{Step 1 — Apply the theorem on intersecting axes.} The composite of two reflections in axes intersecting at $O$ is the rotation of centre $O$ whose angle is twice the angle from the first axis to the second:
\[ f = r\left(O,\; 2 \times \tfrac{\pi}{6}\right) = r\left(O,\; \tfrac{\pi}{3}\right). \]

\textbf{Step 2 — Justification.} The point $O$ lies on both axes, so $S_{d_1}(O) = O$ and $S_{d_2}(O) = O$: $O$ is fixed by $f$. Each reflection preserves distances, so $f$ is an isometry fixing $O$; being a composite of two indirect isometries, $f$ is direct, hence a rotation about $O$. For the angle, let $\alpha$ and $\beta$ be the angles that $d_1$ and $d_2$ make with a fixed reference direction, so $\beta - \alpha = \theta = \frac{\pi}{6}$. A reflection in a line through $O$ making angle $\alpha$ sends a direction of angle $\varphi$ to a direction of angle $2\alpha - \varphi$. Hence, if $\overrightarrow{OM}$ has angle $\varphi$, then $\overrightarrow{OM_1}$ (with $M_1 = S_{d_1}(M)$) has angle $2\alpha - \varphi$, and $\overrightarrow{OM'}$ has angle
\[ 2\beta - (2\alpha - \varphi) = \varphi + 2(\beta - \alpha) = \varphi + 2\theta. \]
Therefore $\left(\overrightarrow{OM}, \overrightarrow{OM'}\right) = 2\theta = \dfrac{\pi}{3}$, as claimed.

\textbf{Step 3 — Image of $M \in d_1$ with $OM = 4$ cm.} Since $M \in d_1$, $S_{d_1}(M) = M$, so $M' = S_{d_2}(M)$: $M'$ is the mirror image of $M$ in $d_2$. Equivalently, $M' = r(O,\frac{\pi}{3})(M)$, so $OM' = OM = 4$ cm and $\left(\overrightarrow{OM}, \overrightarrow{OM'}\right) = \frac{\pi}{3}$. The triangle $OMM'$ is isosceles at $O$ with vertex angle $\frac{\pi}{3}$, hence \textbf{equilateral}: $MM' = 4$ cm.
\end{exoresolu}

\courssub{Composing translations, rotations and glide reflections}

Translations and rotations are the \emph{direct} isometries; composing them can never produce a reflection. The results below are obtained either by Chasles' relation (for translations) or by decomposing each rotation into two reflections with a cleverly chosen common axis.

\begin{methode}[Composing translations and rotations]
\begin{enumerate}
\item \textbf{Two translations:} $t_{\vec{u}} \circ t_{\vec{v}} = t_{\vec{u}+\vec{v}}$ (Chasles). Always a translation, commutative.
\item \textbf{Two rotations of the same centre $O$:} $r(O,\alpha) \circ r(O,\beta) = r(O, \alpha+\beta)$. If $\alpha + \beta \equiv 0 \pmod{2\pi}$, the composite is the identity.
\item \textbf{Two rotations of distinct centres $A, B$:} the composite $r(B,\beta) \circ r(A,\alpha)$ is a rotation of angle $\alpha + \beta$ if $\alpha + \beta \not\equiv 0 \pmod{2\pi}$, and a \textbf{translation} if $\alpha + \beta \equiv 0 \pmod{2\pi}$. To find the new centre, decompose each rotation into two reflections sharing a common axis (usually the line $(AB)$).
\item \textbf{Translation $\circ$ reflection} (or reflection $\circ$ translation): in general a \cle{glide reflection}; it reduces to a pure reflection exactly when the translation vector is perpendicular to the axis.
\item Reflex: to locate the centre of a composite rotation, look for a fixed point: solve $f(M) = M$.
\end{enumerate}
\end{methode}

\begin{exoresolu}[Composite of two rotations with distinct centres]
Let $A$ and $B$ be two distinct points. Let $r_A = r\left(A, \dfrac{\pi}{2}\right)$ and $r_B = r\left(B, \dfrac{\pi}{2}\right)$. Determine the nature of $f = r_B \circ r_A$ and construct its centre when $AB = 2$ cm.
\tcblower
\textbf{Solution.}

\textbf{Step 1 — Nature.} The sum of the angles is $\frac{\pi}{2} + \frac{\pi}{2} = \pi \not\equiv 0 \pmod{2\pi}$. Therefore $f$ is a \textbf{rotation of angle $\pi$}, i.e.\ a \cle{central symmetry} (half-turn). Its centre $\Omega$ remains to be found.

\textbf{Step 2 — Decomposition into reflections.} Write each rotation as a composite of two reflections whose axes pass through its centre, choosing the common axis $(AB)$:
\[ r_A = S_{(AB)} \circ S_{d_1}, \qquad r_B = S_{d_2} \circ S_{(AB)}, \]
where $d_1$ is the line through $A$ making angle $-\frac{\pi}{4}$ with $(AB)$, and $d_2$ the line through $B$ making angle $+\frac{\pi}{4}$ with $(AB)$ (each rotation angle is twice the angle between its axes). Then
\[ f = r_B \circ r_A = S_{d_2} \circ S_{(AB)} \circ S_{(AB)} \circ S_{d_1} = S_{d_2} \circ S_{d_1}. \]

\textbf{Step 3 — The centre.} The lines $d_1$ and $d_2$ make angles $\mp\frac{\pi}{4}$ with $(AB)$; they intersect at a point $\Omega$ such that the triangle $A\Omega B$ has angles $\frac{\pi}{4}, \frac{\pi}{4}, \frac{\pi}{2}$: it is right-angled isosceles at $\Omega$. Hence $\Omega$ is the point with $\Omega A = \Omega B$ and $(\overrightarrow{\Omega A}, \overrightarrow{\Omega B}) = \frac{\pi}{2}$ — the apex of the right isosceles triangle built on $[AB]$. With $AB = 2$ cm: $\Omega A = \Omega B = \frac{AB}{\sqrt{2}} = \sqrt{2}$ cm.

\textbf{Conclusion.} $f$ is the half-turn about the point $\Omega$ such that $A\Omega B$ is right-angled isosceles at $\Omega$.
\end{exoresolu}

\begin{exoresolu}[Recognising a glide reflection]
In an orthonormal frame, let $d$ be the $x$-axis and $\vec{u} = 4\vec{i}$. Let $f = t_{\vec{u}} \circ S_d$. Find the image of $M(x,y)$, show that $f$ has no fixed point, and express $f$ as a glide reflection (axis and vector).
\tcblower
\textbf{Solution.}

\textbf{Step 1 — Image of $M(x,y)$.} First $S_d(M) = M_1(x, -y)$. Then $t_{\vec{u}}(M_1) = M'(x+4, -y)$. Hence
\[ f : (x, y) \longmapsto (x+4,\; -y). \]

\textbf{Step 2 — Fixed points.} $M$ is fixed iff $x + 4 = x$ and $-y = y$. The first equation gives $4 = 0$: impossible. So $f$ has \textbf{no fixed point}.

\textbf{Step 3 — Nature.} Since $\vec{u} = 4\vec{i}$ is \emph{parallel} to the axis $d$ (the $x$-axis), the decomposition cannot be simplified to a pure reflection: $f$ is a \textbf{glide reflection} with axis $d$ (the line $y = 0$) and vector $\vec{u} = 4\vec{i}$.

\textbf{Step 4 — Consistency check.} $f \circ f = t_{\vec{u}} \circ S_d \circ t_{\vec{u}} \circ S_d = t_{2\vec{u}}$ (since $S_d \circ t_{\vec{u}} \circ S_d = t_{\vec{u}}$ when $\vec{u} \parallel d$). Indeed $f(f(x,y)) = f(x+4, -y) = (x+8, y)$, a translation by $8\vec{i}$: the square of a glide reflection is always a translation. $\blacksquare$
\end{exoresolu}

\courssub{Classifying a plane isometry from given data}

An \cle{isometry} of the plane is a transformation that preserves distances: $M'N' = MN$ for all points $M, N$. Isometries are of two kinds: \cle{direct} isometries, which preserve the orientation of angles (translations and rotations), and \cle{indirect} isometries, which reverse it (reflections and glide reflections). The classification theorem states that these four types are the \emph{only} plane isometries; the fixed points tell us which type we are facing.

\begin{methode}[Recognition table for isometries]
An isometry is \cle{direct} (preserves orientation: translations, rotations) or \cle{indirect} (reflections, glide reflections). Use fixed points:
\begin{enumerate}
\item Compute (or test) the fixed points of $f$.
\item \textbf{Three non-collinear fixed points} $\Rightarrow$ identity. \textbf{Exactly a line of fixed points} $\Rightarrow$ reflection in that line.
\item \textbf{Exactly one fixed point} $\Rightarrow$ rotation about that point (find the angle from one point and its image).
\item \textbf{No fixed point and direct} $\Rightarrow$ translation (vector $= \overrightarrow{MM'}$ for any $M$).
\item \textbf{No fixed point and indirect} $\Rightarrow$ glide reflection: find the axis as the line through the midpoints of the segments $[MM']$, and the vector from the displacement along the axis.
\item Orientation test: a direct isometry preserves the sign of angles; analytically, the matrix $\begin{pmatrix} a & -b \\ b & a \end{pmatrix}$ (with $a^2+b^2=1$) is direct, while $\begin{pmatrix} a & b \\ b & -a \end{pmatrix}$ is indirect.
\end{enumerate}
\end{methode}

\begin{exoresolu}[Identifying an isometry given analytically]
The plane is referred to an orthonormal frame $(O;\vec{i},\vec{j})$. A transformation $f$ maps $M(x,y)$ to $M'(x',y')$ with
\[ x' = \tfrac{1}{2}x - \tfrac{\sqrt{3}}{2}y + 1, \qquad y' = \tfrac{\sqrt{3}}{2}x + \tfrac{1}{2}y - 2. \]
Show that $f$ is an isometry, determine its nature and its characteristic elements.
\tcblower
\textbf{Solution.}

\textbf{Step 1 — Isometry and orientation.} The linear part has matrix $\begin{pmatrix} \frac12 & -\frac{\sqrt3}{2} \\[2pt] \frac{\sqrt3}{2} & \frac12 \end{pmatrix}$. Its columns are orthonormal: $\left(\frac12\right)^2 + \left(\frac{\sqrt3}{2}\right)^2 = 1$ and the determinant is $\frac14 + \frac34 = 1 > 0$. Hence $f$ is a \textbf{direct isometry}: a translation or a rotation.

\textbf{Step 2 — Fixed points.} Solve $f(M) = M$:
\[ \begin{cases} \frac12 x - \frac{\sqrt3}{2} y + 1 = x \\ \frac{\sqrt3}{2} x + \frac12 y - 2 = y \end{cases} \Longleftrightarrow \begin{cases} -\frac12 x - \frac{\sqrt3}{2} y = -1 \\ \frac{\sqrt3}{2} x - \frac12 y = 2 \end{cases} \Longleftrightarrow \begin{cases} x + \sqrt{3}\, y = 2 \\ \sqrt{3}\, x - y = 4. \end{cases} \]
Multiply the second equation by $\sqrt{3}$: $3x - \sqrt{3}y = 4\sqrt{3}$. Adding to the first: $4x = 2 + 4\sqrt{3}$, so $x = \frac{1 + 2\sqrt{3}}{2}$. Then $y = \sqrt{3}x - 4 = \frac{\sqrt{3} + 6}{2} - 4 = \frac{\sqrt{3} - 2}{2}$.

There is exactly \textbf{one fixed point} $\Omega\left(\dfrac{1+2\sqrt3}{2}, \dfrac{\sqrt3 - 2}{2}\right)$, so $f$ is a \textbf{rotation of centre $\Omega$}.

\textbf{Step 3 — Angle.} The linear part is the matrix of the rotation of angle $\theta$ with $\cos\theta = \frac12$, $\sin\theta = \frac{\sqrt3}{2}$, i.e.\ $\theta = \dfrac{\pi}{3}$.

\textbf{Conclusion.} $f = r\left(\Omega, \dfrac{\pi}{3}\right)$ with $\Omega$ as above.
\end{exoresolu}

\courssub{Constructing and using a homothety}

Enlarging a photograph, drawing a plan of a school campus to scale, projecting a shadow: all these situations suggest a transformation that \emph{multiplies all distances by the same factor} from a fixed point. This is the idea of a homothety.

\begin{definition}[Homothety]
Let $\Omega$ be a point and $k$ a non-zero real number. The \cle{homothety} of centre $\Omega$ and ratio $k$, denoted $h(\Omega, k)$, is the transformation sending every point $M$ to the point $M'$ defined by
\[ \overrightarrow{\Omega M'} = k\,\overrightarrow{\Omega M}. \]
If $k > 0$ the homothety is called \cle{direct} (positive); if $k < 0$ it is called \cle{indirect} (negative), and $M'$ lies on the ray opposite to $[\Omega M)$.
\end{definition}

\begin{popfigure}
\begin{tikzpicture}[scale=1.0]
\fill (0,0) circle (1.5pt) node[below left] {$\Omega$};
\fill (2,0.8) circle (1.5pt) node[above] {$M$};
\fill (3,1.2) circle (1.5pt) node[above] {$M'$};
\fill (-1.5,-0.6) circle (1.5pt) node[below] {$N'$};
\fill (1,0.4) circle (1.5pt) node[above right] {$N$};
\draw[popmuted] (-2.2,-0.88) -- (3.6,1.44);
\draw[->, popblue, thick] (0,0) -- (2,0.8);
\draw[->, poporange, thick] (0,0) -- (3,1.2);
\draw[->, popgreen, thick] (0,0) -- (-1.5,-0.6);
\node[popdark] at (4.6,1.0) {$k=\tfrac32$ : $\overrightarrow{\Omega M'}=\tfrac32\overrightarrow{\Omega M}$};
\node[popdark] at (4.6,-0.6) {$k=-\tfrac32$ : $\overrightarrow{\Omega N'}=-\tfrac32\overrightarrow{\Omega N}$};
\end{tikzpicture}
\legende{Homothety of centre $\Omega$: positive ratio keeps the direction, negative ratio reverses it.}
\end{popfigure}

\begin{methode}[Homothety $h(\Omega, k)$: construction and properties]
\begin{enumerate}
\item \textbf{Image of a point:} $M' = h(M) \iff \overrightarrow{\Omega M'} = k\,\overrightarrow{\Omega M}$. If $k > 0$, $M'$ is on the ray $[\Omega M)$; if $k < 0$, on the opposite ray (indirect/negative homothety).
\item \textbf{Image of a line} $d$: a line $d'$ \emph{parallel} to $d$ (if $\Omega \notin d$); $d$ is globally invariant if $\Omega \in d$.
\item \textbf{Image of a circle} $\mathcal{C}(O, R)$: the circle of centre $O' = h(O)$ and radius $|k|R$.
\item \textbf{Distances} are multiplied by $|k|$; \textbf{areas} by $k^2$; angles and alignment are preserved.
\item To \textbf{find the centre} of a homothety mapping $A \mapsto A'$, $B \mapsto B'$: the centre lies at the intersection of the lines $(AA')$ and $(BB')$; the ratio is $k = \dfrac{\overrightarrow{\Omega A'}}{\overrightarrow{\Omega A}}$.
\item $k = 1$ gives the identity; $k = -1$ gives the central symmetry about $\Omega$.
\end{enumerate}
\end{methode}

\begin{exoresolu}[Determining a homothety from two pairs of points]
In an orthonormal frame, $A(1,2)$, $B(4,2)$ have images $A'(3,6)$, $B'(9,6)$ under a homothety $h$. Find the centre and ratio of $h$, then the image of the circle of centre $C(2,0)$ and radius $1$.
\tcblower
\textbf{Solution.}

\textbf{Step 1 — The centre.} The centre $\Omega$ lies on $(AA')$ and on $(BB')$.
Line $(AA')$: through $A(1,2)$ and $A'(3,6)$, direction $(2,4) \parallel (1,2)$, equation $y = 2x$.
Line $(BB')$: through $B(4,2)$ and $B'(9,6)$, direction $(5,4)$, equation $y - 2 = \frac45(x - 4)$.

Intersecting: $2x - 2 = \frac45(x-4) \Rightarrow 10x - 10 = 4x - 16 \Rightarrow 6x = -6 \Rightarrow x = -1$, $y = -2$. Hence $\Omega(-1, -2)$.

\textbf{Step 2 — The ratio.} $\overrightarrow{\Omega A} = (1-(-1), 2-(-2)) = (2,4)$ and $\overrightarrow{\Omega A'} = (3-(-1), 6-(-2)) = (4,8) = 2(2,4)$. So $k = 2$. Check with $B$: $\overrightarrow{\Omega B} = (5,4)$, $2\overrightarrow{\Omega B} = (10,8) = \overrightarrow{\Omega B'}$. \checkmark

Thus $h = h(\Omega, 2)$ with $\Omega(-1,-2)$.

\textbf{Step 3 — Image of the circle.} The image is the circle of centre $C' = h(C)$ and radius $|k| \times 1 = 2$.
\[ \overrightarrow{\Omega C'} = 2\overrightarrow{\Omega C} = 2(2-(-1), 0-(-2)) = 2(3,2) = (6,4), \]
so $C' = \Omega + (6,4) = (5, 2)$. The image circle has centre $(5,2)$ and radius $2$: equation $(x-5)^2 + (y-2)^2 = 4$.
\end{exoresolu}

\courssub{Analytical and complex representation of homothety}

\begin{methode}[Formulas for $h(\Omega, k)$]
\begin{enumerate}
\item \textbf{Vector form:} $\overrightarrow{\Omega M'} = k\,\overrightarrow{\Omega M}$.
\item \textbf{Analytical form} (frame $(O;\vec{i},\vec{j})$, $\Omega(a,b)$):
\[ \begin{cases} x' = kx + (1-k)a \\ y' = ky + (1-k)b. \end{cases} \]
\item \textbf{Complex form} ($\Omega$ of affix $\omega$, $M$ of affix $z$): $z' = kz + (1-k)\omega$ with $k \in \mathbb{R}^*$.
\item Conversely, any map $z \mapsto kz + p$ with $k \in \mathbb{R} \setminus \{0,1\}$ is a homothety of ratio $k$; its centre has affix $\omega = \dfrac{p}{1-k}$ (the unique fixed point).
\item Reflex: to identify the transformation $z \mapsto az + b$, first check whether $a \in \mathbb{R}$ (homothety/translation) or $|a| = 1$ (rotation/translation).
\end{enumerate}
\end{methode}

\begin{exoresolu}[Identifying a homothety from its complex expression]
The plane is equipped with an orthonormal frame; to each point $M$ we associate its affix $z$. Let $f$ be the map defined by
\[ z' = -3z + 8 - 4i. \]
Determine the nature of $f$, its centre and its ratio. Find the image of the point $E$ of affix $1 + i$.
\tcblower
\textbf{Solution.}

\textbf{Step 1 — Nature.} The map has the form $z' = az + b$ with $a = -3 \in \mathbb{R}^* \setminus \{1\}$. Hence $f$ is a \textbf{homothety of ratio $k = -3$} (an indirect, i.e.\ negative, homothety).

\textbf{Step 2 — Centre.} The centre $\Omega$ is the unique fixed point: $\omega = -3\omega + 8 - 4i$, so $4\omega = 8 - 4i$ and
\[ \omega = 2 - i. \]
Check with the formula $\omega = \dfrac{b}{1-a} = \dfrac{8-4i}{1-(-3)} = \dfrac{8-4i}{4} = 2 - i$. \checkmark

\textbf{Step 3 — Standard form.} $f = h(\Omega, -3)$, i.e.\ $z' - \omega = -3(z - \omega)$. Indeed $-3(z - 2 + i) + 2 - i = -3z + 6 - 3i + 2 - i = -3z + 8 - 4i$. \checkmark

\textbf{Step 4 — Image of $E(1+i)$.} $z'_E = -3(1+i) + 8 - 4i = -3 - 3i + 8 - 4i = 5 - 7i$. So $E'$ has affix $5 - 7i$.

Verification: $\overrightarrow{\Omega E}$ has affix $(1+i) - (2-i) = -1 + 2i$; $-3(-1+2i) = 3 - 6i$; $\Omega + (3-6i) = 5 - 7i$. \checkmark
\end{exoresolu}

\courssub{Composing two homotheties}

\begin{methode}[Composite of two homotheties]
Let $h_1 = h(\Omega_1, k_1)$ and $h_2 = h(\Omega_2, k_2)$.
\begin{enumerate}
\item \textbf{Same centre $\Omega$:} $h(\Omega, k_2) \circ h(\Omega, k_1) = h(\Omega, k_1 k_2)$. The homotheties of centre $\Omega$ commute and form a group.
\item \textbf{Distinct centres:} the composite $h_2 \circ h_1$ multiplies distances by $|k_1 k_2|$ and sends every line to a parallel line.
\begin{itemize}
\item If $k_1 k_2 \neq 1$: it is a homothety of ratio $k_1 k_2$ whose centre $\Omega$ is \textbf{aligned with $\Omega_1$ and $\Omega_2$}; more precisely $\overrightarrow{\Omega_2 \Omega} = \dfrac{k_2(1-k_1)}{1 - k_1k_2}\,\overrightarrow{\Omega_2\Omega_1}$. In practice: solve $f(M) = M$.
\item If $k_1 k_2 = 1$: it is a \textbf{translation} of vector $(1 - k_2)\,\overrightarrow{\Omega_1 \Omega_2}$. In practice, compute the image of one point: $\vec{u} = \overrightarrow{M\,f(M)}$.
\end{itemize}
\item Strategy: write the complex or analytical expressions and compose them; then identify.
\end{enumerate}
\end{methode}

\begin{exoresolu}[Composite of two homotheties with distinct centres]
Let $h_1 = h(A, 2)$ and $h_2 = h(B, \tfrac12)$ where $A$ and $B$ have affixes $a = 0$ and $b = 3$ in an orthonormal frame. Determine the nature and elements of $f = h_2 \circ h_1$, then of $g = h_1 \circ h_2$. Comment.
\tcblower
\textbf{Solution.}

\textbf{Step 1 — Complex expressions.} $h_1 : z \mapsto 2z$ (centre $0$, ratio $2$). $h_2 : z \mapsto \frac12 z + (1-\frac12)\cdot 3 = \frac12 z + \frac32$.

\textbf{Step 2 — Composite $f = h_2 \circ h_1$.}
\[ f(z) = h_2(2z) = \tfrac12 (2z) + \tfrac32 = z + \tfrac32. \]
Since $k_1 k_2 = 2 \times \frac12 = 1$, $f$ is a \textbf{translation} of vector affix $\frac32$, i.e.\ $\vec{u} = \frac32\vec{i}$. Note $\vec{u} = (1-k_2)\overrightarrow{AB} = \frac12 \overrightarrow{AB}$: consistent, since $\overrightarrow{AB}$ has affix $3$.

\textbf{Step 3 — Composite $g = h_1 \circ h_2$.}
\[ g(z) = h_1\!\left(\tfrac12 z + \tfrac32\right) = 2\left(\tfrac12 z + \tfrac32\right) = z + 3. \]
So $g$ is the \textbf{translation} of vector affix $3$, i.e.\ $\vec{v} = \overrightarrow{AB}$.

\textbf{Step 4 — Comment.} $f \neq g$: composition of homotheties with distinct centres is \textbf{not commutative}. Both are translations (because the product of ratios is $1$), but with different vectors: $\vec{v} = 2\vec{u}$. This is a classic GCE trap: always compute the composite in the order requested.
\end{exoresolu}

\courssub{Plane similarities: recognition and applications}

Homotheties multiply distances by $|k|$; rotations and translations preserve them. Composing such maps yields transformations that multiply all distances by a fixed positive factor while preserving angles: these are the \cle{plane similarities}, the most general "shape-preserving" transformations of the plane.

\begin{methode}[Studying a plane similarity]
\begin{enumerate}
\item A \cle{similarity} of ratio $k > 0$ multiplies all distances by $k$: $M'N' = k \cdot MN$. It preserves angles, alignment, parallelism, midpoints, and sends circles to circles.
\item \textbf{Complex test:} $z' = az + b$ ($a \neq 0$) is a \cle{direct similarity}: ratio $|a|$, angle $\arg(a)$. $z' = a\bar{z} + b$ is an \cle{indirect similarity}.
\item Identification of $z' = az + b$:
\begin{itemize}
\item $a = 1$: translation of vector affix $b$.
\item $|a| = 1$, $a \neq 1$: rotation of angle $\arg a$ about the fixed point $\omega = \frac{b}{1-a}$.
\item $a \in \mathbb{R}\setminus\{0,1\}$: homothety of ratio $a$, centre $\omega = \frac{b}{1-a}$.
\item General case ($|a| \neq 1$, $a \notin \mathbb{R}$): direct similarity $= h(\Omega, |a|) \circ r(\Omega, \arg a)$, centre $\Omega$ of affix $\frac{b}{1-a}$.
\end{itemize}
\item A \cle{spiral similarity} is exactly the composite $h \circ r$ of a homothety and a rotation of the same centre: it is the general direct similarity with a fixed point.
\item Applications: two circles are always homothetic (the two homothety centres lie on the line of centres); similar triangles are linked by one similarity.
\end{enumerate}
\end{methode}

\begin{exoresolu}[Identifying and decomposing a direct similarity]
Let $s$ be the plane transformation given in complex form by
\[ z' = (1 + i)\,z + 2 - i. \]
\begin{enumerate}
\item Show that $s$ is a direct similarity; give its ratio and angle.
\item Determine its fixed point $\Omega$.
\item Write $s$ as the composite of a homothety and a rotation of centre $\Omega$.
\item Find the image of the circle $\mathcal{C}$ of centre $O$ (affix $0$) and radius $2$.
\end{enumerate}
\tcblower
\textbf{Solution.}

\textbf{Step 1 — Nature, ratio, angle.} $s$ has the form $z' = az + b$ with $a = 1 + i \neq 0$: it is a \textbf{direct similarity}. Ratio: $k = |a| = |1+i| = \sqrt{2}$. Angle: $\arg(a) = \arg(1+i) = \dfrac{\pi}{4}$. Since $|a| \neq 1$ and $a \notin \mathbb{R}$, $s$ is neither a rotation nor a homothety: it is a genuine spiral similarity.

\textbf{Step 2 — Fixed point.} Solve $\omega = (1+i)\omega + 2 - i$:
\[ \omega - (1+i)\omega = 2 - i \;\Longleftrightarrow\; -i\omega = 2 - i \;\Longleftrightarrow\; \omega = \frac{2-i}{-i} = \frac{(2-i)\,i}{(-i)(i)} = \frac{2i + 1}{1} = 1 + 2i. \]
Hence $\Omega$ has affix $\omega = 1 + 2i$.

\textbf{Step 3 — Decomposition.} Any direct similarity with fixed point $\Omega$ writes $z' - \omega = a(z - \omega)$, i.e.
\[ s = h(\Omega, \sqrt{2}) \circ r\!\left(\Omega, \tfrac{\pi}{4}\right) = r\!\left(\Omega, \tfrac{\pi}{4}\right) \circ h(\Omega, \sqrt{2}), \]
the two maps commuting since they share the same centre. Check: $z' = \omega + (1+i)(z - \omega) = (1+i)z + \omega(1 - 1 - i) = (1+i)z - i(1+2i) = (1+i)z + 2 - i$. \checkmark

\textbf{Step 4 — Image of the circle $\mathcal{C}(O, 2)$.} A similarity of ratio $\sqrt{2}$ sends a circle to a circle whose radius is multiplied by $\sqrt{2}$ and whose centre is the image of the centre. Image of $O$ (affix $0$): $z'_{O} = (1+i)\cdot 0 + 2 - i = 2 - i$. Hence $s(\mathcal{C})$ is the circle of centre $O'(2 - i)$ and radius $2\sqrt{2}$.
\end{exoresolu}

\begin{attention}[Frequent errors to avoid]
\begin{itemize}
\item Confusing the order in $f \circ g$: $g$ acts first. For reflections, $S_{d_2} \circ S_{d_1}$ has vector/angle directed from $d_1$ to $d_2$.
\item Forgetting the factor $2$: the translation vector is \emph{twice} the gap between parallel axes; the rotation angle is \emph{twice} the angle between intersecting axes.
\item Claiming two rotations of distinct centres always give a rotation: when the angles sum to $0 \pmod{2\pi}$ the composite is a translation.
\item Writing the homothety ratio as $k$ in the radius: the radius is multiplied by $|k|$, never by a negative number.
\item For homotheties with $k_1k_2 = 1$ and distinct centres, concluding "homothety of ratio 1": it is a translation, not the identity in general.
\item In $z' = az + b$, concluding "rotation" without checking $|a| = 1$, or "homothety" without checking $a \in \mathbb{R}$.
\end{itemize}
\end{attention}

\newpage

\coursec{Exercises}

\courssub{I check my knowledge}

\begin{exercice}[MCQ: nature of a composite]
For each question, choose the single correct answer and justify briefly.
\begin{enumerate}
\item The composite $s_2 \circ s_1$ of the reflections in two parallel lines $d_1$, $d_2$ with $d_1 \neq d_2$ is:
\begin{enumerate}
\item a rotation \item a translation \item a reflection \item a glide reflection
\end{enumerate}
\item The composite $r_2 \circ r_1$ of two rotations of angles $\theta_1$ and $\theta_2$ with $\theta_1 + \theta_2 \not\equiv 0 \pmod{2\pi}$ is:
\begin{enumerate}
\item always a translation \item always a rotation \item always a homothety \item a reflection
\end{enumerate}
\item The transformation of complex expression $z' = 3z - 4 + 2i$ is:
\begin{enumerate}
\item a translation \item a rotation \item a homothety of ratio $3$ \item an indirect similarity
\end{enumerate}
\item A homothety of ratio $k = -2$ multiplies areas by:
\begin{enumerate}
\item $-2$ \item $2$ \item $4$ \item $-4$
\end{enumerate}
\end{enumerate}
\end{exercice}

\begin{exercice}[True or false]
Say whether each statement is true or false, and justify your answer.
\begin{enumerate}
\item Every isometry of the plane is a translation, a rotation or a reflection.
\item The composite of two translations is a translation whose vector is the sum of the two vectors.
\item Two homotheties of distinct centres $O_1 \neq O_2$ and ratios $k_1, k_2$ with $k_1 k_2 = 1$ compose to give a translation.
\item A direct similarity of ratio $k = 1$ is an isometry.
\item A homothety of ratio $-1$ is a reflection in a line.
\end{enumerate}
\end{exercice}

\begin{exercice}[Definitions]
Give a precise definition of each of the following, stating its characteristic elements:
\begin{enumerate}
\item a glide reflection;
\item a homothety of centre $\Omega$ and ratio $k$;
\item a direct plane similarity;
\item an indirect plane similarity.
\end{enumerate}
\end{exercice}

\begin{exercice}[Direct calculations]
\begin{enumerate}
\item Let $t_{\vec{u}}$ and $t_{\vec{v}}$ be the translations of vectors $\vec{u}(2,-1)$ and $\vec{v}(3,4)$. Determine the vector of $t_{\vec{v}} \circ t_{\vec{u}}$ and the image of $M(1,5)$.
\item Find the complex expression of the homothety of centre $\Omega$ of affix $1 - i$ and ratio $k = -3$, then the image of the point of affix $2 + i$.
\item Find the complex expression of the rotation of centre $A(1,0)$ and angle $\dfrac{\pi}{3}$, and the image of $B(4,0)$.
\end{enumerate}
\end{exercice}

\courssub{I practise}

\begin{exercice}[Composing reflections in parallel lines]
$d_1$ and $d_2$ are two parallel lines $3$ cm apart, and $s_1$, $s_2$ are the reflections in $d_1$, $d_2$.
\begin{enumerate}
\item Describe precisely the transformation $s_2 \circ s_1$ (nature and characteristic elements).
\item Construct the image of a given triangle $ABC$ by $s_2 \circ s_1$.
\item Find two lines whose reflections compose to give the translation of vector $\vec{u}(4,0)$.
\end{enumerate}
\end{exercice}

\begin{exercice}[Composing two rotations]
Let $r_1 = r(A, 60^{\circ})$ and $r_2 = r(B, -60^{\circ})$ with $A \neq B$.
\begin{enumerate}
\item Prove that $r_2 \circ r_1$ is a translation.
\item Determine its vector in terms of $A$ and $B$.
\item What is the nature of $r_2 \circ r_1$ if instead $r_2 = r(B, 120^{\circ})$? Give its characteristic elements.
\end{enumerate}
\end{exercice}

\begin{exercice}[Recognising an isometry from images]
An isometry $f$ maps $A(0,0)$ to $A'(2,1)$, $B(3,0)$ to $B'(5,1)$ and $C(0,4)$ to $C'(-2,5)$.
\begin{enumerate}
\item Show that $f$ is a translation and give its vector.
\item What is the nature of $f$ if instead $C' = (6,1)$? Determine its characteristic elements.
\end{enumerate}
\end{exercice}

\begin{exercice}[Homothety from its complex expression]
Let $f$ be the transformation of the complex plane given by $z' = 3z - 4 + 2i$.
\begin{enumerate}
\item Prove that $f$ is a homothety, and find its centre and ratio.
\item Compute the image by $f$ of the circle of centre $1 + i$ and radius $2$.
\item Determine the image of the line of equation $z + \bar{z} = 4$.
\end{enumerate}
\end{exercice}

\courssub{I prepare for the GCE}

\begin{exercice}[Composing homotheties] \hfill (8 marks)
Let $h_1 = h(O_1, 2)$ and $h_2 = h\!\left(O_2, \dfrac{1}{2}\right)$ with $O_1 \neq O_2$.
\begin{enumerate}
\item Prove that $h_2 \circ h_1$ is a translation and express its vector in terms of $\overrightarrow{O_1 O_2}$. \hfill (3 marks)
\item Construct this vector on a figure. \hfill (1 mark)
\item Repeat the study with $h_2 = h(O_2, 3)$: prove that $h_2 \circ h_1$ is a homothety, give its ratio, and determine the position of its centre on the line $(O_1 O_2)$. \hfill (4 marks)
\end{enumerate}
\end{exercice}

\begin{exercice}[Centres of similitude of two circles] \hfill (10 marks)
Two circles $\mathcal{C}$ and $\mathcal{C}'$ have centres $A$, $B$ with $AB = 10$ cm and radii $2$ cm and $6$ cm respectively.
\begin{enumerate}
\item Prove that there exist exactly two homotheties mapping $\mathcal{C}$ onto $\mathcal{C}'$, and give their ratios. \hfill (3 marks)
\item Locate the two centres of similitude on the line $(AB)$ by computing their distances from $A$ and $B$. \hfill (3 marks)
\item Show that a common external tangent to the two circles passes through one of the centres of similitude. \hfill (2 marks)
\item Construct the two centres of similitude and the common tangents on a figure. \hfill (2 marks)
\end{enumerate}
\end{exercice}

\begin{exercice}[Euler line and direct similarity] \hfill (12 marks)
\textbf{Part A.} In a triangle $ABC$, let $G$ be the centroid, $O$ the circumcentre and $H$ the orthocentre.
\begin{enumerate}
\item Show that the homothety $h(G, -2)$ maps the midpoint of $[BC]$ onto $A$, and deduce that it maps $O$ onto $H$. \hfill (3 marks)
\item Deduce that $O$, $G$, $H$ are collinear and that $GH = 2\,GO$. \hfill (2 marks)
\end{enumerate}
\textbf{Part B.} A direct similarity $S$ has complex expression $z' = (1 + i\sqrt{3})z + 2 - i$.
\begin{enumerate}
\item Determine the ratio and the angle of $S$, and classify $S$. \hfill (3 marks)
\item Determine the affix of the fixed point (centre) of $S$. \hfill (2 marks)
\item Find the image of the point of affix $1$. \hfill (1 mark)
\item Prove that the composite of a glide reflection and a homothety of ratio $k \neq 1$ is an indirect similarity with exactly one fixed point, and illustrate with $z' = 2\bar{z} + 3i$. \hfill (1 mark)
\end{enumerate}
\end{exercice}

\newpage

\coursec{Solutions to the exercises}

\begin{corrige}[Exercise 1 — MCQ: nature of a composite]
\begin{enumerate}
\item \textbf{Answer: (b) a translation.} The composite of the reflections in two \emph{distinct parallel} lines $d_1 \parallel d_2$ is the translation of vector $2\vec{u}$, where $\vec{u}$ is the vector from $d_1$ to $d_2$ perpendicular to both lines. It cannot be a rotation (no fixed point, and the angle of rotation would be $0$ with distinct axes).
\item \textbf{Answer: (b) always a rotation.} The composite of two rotations of angles $\theta_1, \theta_2$ is a rotation of angle $\theta_1+\theta_2$ when $\theta_1+\theta_2 \not\equiv 0 \pmod{2\pi}$ (and a translation when the sum is $0$ modulo $2\pi$). Since the sum is non-zero here, the composite is a rotation.
\item \textbf{Answer: (c) a homothety of ratio $3$.} The expression has the form $z' = kz + b$ with $k = 3 \in \mathbb{R}$, $k \neq 1$: this is a homothety of ratio $3$ (its centre is the fixed point $\omega = \dfrac{b}{1-k} = \dfrac{-4+2i}{-2} = 2 - i$).
\item \textbf{Answer: (c) $4$.} A homothety of ratio $k$ multiplies lengths by $|k|$ and areas by $k^2$. Here $k^2 = (-2)^2 = 4$.
\end{enumerate}
\end{corrige}

\begin{corrige}[Exercise 2 — True or false]
\begin{enumerate}
\item \textbf{False.} The classification theorem states that every plane isometry is a translation, a rotation, a reflection \emph{or a glide reflection}. The glide reflection (composite of a reflection and a translation parallel to its axis) has been forgotten. Counter-example: $z' = \bar{z} + 1$ is a glide reflection, which is none of the three listed types.
\item \textbf{True.} For every point $M$, $t_{\vec{v}}\circ t_{\vec{u}}(M) = M + \vec{u} + \vec{v} = M + (\vec{u}+\vec{v})$, so the composite is the translation of vector $\vec{u}+\vec{v}$. (Translations commute: $t_{\vec{v}}\circ t_{\vec{u}} = t_{\vec{u}}\circ t_{\vec{v}}$.)
\item \textbf{True.} $h(O_2,k_2)\circ h(O_1,k_1)$ has ratio $k_1k_2 = 1$. A homothety-type composite of ratio $1$ with distinct centres is a translation; its vector is $(1-k_2)\overrightarrow{O_1O_2}$ (equivalently $(k_1-1)\overrightarrow{O_2O_1}$), which is non-zero since $k_2 \neq 1$ (otherwise $k_1k_2=1$ would force $k_1=1$ and the centres could coincide).
\item \textbf{True.} A direct similarity of ratio $k$ multiplies all distances by $k$. If $k = 1$, distances are preserved: the similarity is a (direct) isometry, i.e.\ a translation or a rotation.
\item \textbf{False.} A homothety of ratio $-1$ is the \emph{central symmetry} (reflection in a \emph{point}, the centre), i.e.\ the rotation of angle $\pi$ about its centre. It is a \emph{direct} isometry, whereas a reflection in a line is an \emph{indirect} isometry. They are never equal.
\end{enumerate}
\end{corrige}

\begin{corrige}[Exercise 3 — Definitions]
\begin{enumerate}
\item \textbf{Glide reflection.} A glide reflection is the composite $t_{\vec{u}} \circ s_d = s_d \circ t_{\vec{u}}$ of the reflection $s_d$ in a line $d$ and a translation $t_{\vec{u}}$ whose vector $\vec{u} \neq \vec{0}$ is parallel to $d$. Characteristic elements: the \emph{axis} $d$ and the \emph{vector} $\vec{u}$ (parallel to $d$). It is an indirect isometry with no fixed point.
\item \textbf{Homothety.} The homothety of centre $\Omega$ and ratio $k \in \mathbb{R}^*$ is the transformation $h(\Omega,k)$ mapping every point $M$ to the point $M'$ defined by $\overrightarrow{\Omega M'} = k\,\overrightarrow{\Omega M}$. Characteristic elements: the \emph{centre} $\Omega$ (the unique fixed point when $k \neq 1$) and the \emph{ratio} $k$.
\item \textbf{Direct plane similarity.} A direct similarity is a transformation of the plane that multiplies all distances by a fixed positive number $k$ \emph{and preserves oriented angles}. Equivalently it is the composite of a homothety of ratio $k$ with a direct isometry (translation or rotation). Characteristic elements: the \emph{ratio} $k > 0$, the \emph{angle} $\theta$, and (when $k \neq 1$) the \emph{centre} $\Omega$, its unique fixed point. Complex form: $z' = az + b$ with $a \in \mathbb{C}^*$, $|a| = k$, $\arg a = \theta$.
\item \textbf{Indirect plane similarity.} An indirect similarity multiplies all distances by a fixed $k > 0$ and \emph{reverses} oriented angles; equivalently it is the composite of a homothety with an indirect isometry (reflection or glide reflection). Complex form: $z' = a\bar{z} + b$, $a \in \mathbb{C}^*$, with ratio $|a|$.
\end{enumerate}
\end{corrige}

\begin{corrige}[Exercise 4 — Direct calculations]
\begin{enumerate}
\item The composite of two translations is the translation of the sum vector:
\[ \vec{u} + \vec{v} = (2+3,\ -1+4) = (5,3). \]
Hence $t_{\vec{v}}\circ t_{\vec{u}} = t_{(5,3)}$ and the image of $M(1,5)$ is
\[ M' = (1+5,\ 5+3) = (6,8). \]
\item A homothety of centre $\omega$ and ratio $k$ has complex expression $z' - \omega = k(z-\omega)$. With $\omega = 1-i$ and $k=-3$:
\[ z' = -3(z-(1-i)) + (1-i) = -3z + 4(1-i) = -3z + 4 - 4i. \]
Image of $z = 2+i$: $z' = -3(2+i) + 4 - 4i = -6 - 3i + 4 - 4i = -2 - 7i$.
\item The rotation of centre $A$ (affix $\omega = 1$) and angle $\theta = \dfrac{\pi}{3}$ has expression $z' - \omega = e^{i\pi/3}(z-\omega)$:
\[ z' = e^{i\pi/3}(z-1) + 1 = \left(\tfrac12 + i\tfrac{\sqrt3}{2}\right)(z-1) + 1. \]
For $B$ of affix $4$: $z' = \left(\tfrac12 + i\tfrac{\sqrt3}{2}\right)\cdot 3 + 1 = \dfrac{5}{2} + \dfrac{3\sqrt3}{2}i$, so $B' = \left(\dfrac{5}{2},\ \dfrac{3\sqrt3}{2}\right)$.
\end{enumerate}
\end{corrige}

\begin{corrige}[Exercise 5 — Composing reflections in parallel lines]
\begin{enumerate}
\item Since $d_1 \parallel d_2$ with $d_1 \neq d_2$, the composite $s_2 \circ s_1$ is a \textbf{translation}. Its vector is $\vec{w} = 2\vec{u}$, where $\vec{u}$ is the vector going perpendicularly from $d_1$ to $d_2$. As the lines are $3$ cm apart, $\|\vec{u}\| = 3$ cm, so $\|\vec{w}\| = 6$ cm, directed from $d_1$ towards $d_2$, perpendicular to both lines.
\item \textbf{Construction.} For each vertex, say $A$: drop the perpendicular from $A$ to $d_1$ and $d_2$; translate $A$ by the vector $\vec{w} = 2\vec{u}$ (i.e.\ move $6$ cm along the perpendicular direction from $d_1$ to $d_2$) to obtain $A'$. Repeat for $B$ and $C$, then join $A'B'C'$. Alternatively, reflect $ABC$ in $d_1$ to get $A_1B_1C_1$, then reflect $A_1B_1C_1$ in $d_2$.
\item We need $2\vec{u} = \vec{u}_{\text{target}} = (4,0)$, so $\vec{u} = (2,0)$: the two lines must be vertical lines $2$ units apart, the second being to the right of the first. For example $d_1: x = 0$ and $d_2: x = 2$; then $s_{d_2}\circ s_{d_1}$ is the translation of vector $(4,0)$. (Infinitely many pairs work: $x = a$ and $x = a+2$.)
\end{enumerate}
\end{corrige}

\begin{corrige}[Exercise 6 — Composing two rotations]
\begin{enumerate}
\item The angles sum to $60^{\circ} + (-60^{\circ}) = 0$. By the theorem on the composition of two rotations, when the sum of the angles is $0 \pmod{360^{\circ}}$ and the centres are distinct, the composite is a \textbf{translation}. Hence $r_2 \circ r_1$ is a translation.
\item Write the rotations as composites of reflections through the line $(AB)$: $r_1 = s_{(AB)} \circ s_{d_1}$ where $d_1$ passes through $A$ with $\big(d_1,(AB)\big) = 30^{\circ}$, and $r_2 = s_{d_2} \circ s_{(AB)}$ where $d_2$ passes through $B$ with $\big((AB),d_2\big) = -30^{\circ}$. Then
\[ r_2 \circ r_1 = s_{d_2} \circ s_{(AB)} \circ s_{(AB)} \circ s_{d_1} = s_{d_2} \circ s_{d_1}. \]
Since $\big(d_1,d_2\big) = 30^{\circ} + (-30^{\circ})\cdot(\text{sign check})$: the lines $d_1$ and $d_2$ are parallel (both make an angle of $30^{\circ}$ with $(AB)$ on appropriate sides), so $s_{d_2}\circ s_{d_1}$ is the translation of vector $2\vec{u}$ where $\vec{u}$ goes from $d_1$ to $d_2$. A direct computation with complex numbers, $z' = e^{-i\pi/3}\big(e^{i\pi/3}(z-a)+a - b\big) + b = z + (a-b)\big(1 - e^{-i\pi/3}\big)$, gives the vector of affix $(a-b)\left(1 - e^{-i\pi/3}\right)$, i.e.
\[ \vec{w} = \big(1 - e^{-i\pi/3}\big)\overrightarrow{BA}. \]
Since $1 - e^{-i\pi/3} = 1 - \tfrac12 + i\tfrac{\sqrt3}{2} = e^{i\pi/3}$, the vector is $\vec{w} = e^{i\pi/3}\overrightarrow{BA}$: it has length $AB$ and makes an angle of $60^{\circ}$ with $\overrightarrow{BA}$.
\item Now the angles sum to $60^{\circ} + 120^{\circ} = 180^{\circ} \neq 0$, so $r_2 \circ r_1$ is a \textbf{rotation of angle $180^{\circ}$}, i.e.\ a \textbf{central symmetry}. Its centre $\Omega$ is found from the fixed-point condition. Using complex expressions with $e^{i\pi/3}$ and $e^{2i\pi/3}$:
\[ z' = e^{2i\pi/3}\big(e^{i\pi/3}(z-a) + a - b\big) + b = -z + a\big(e^{2i\pi/3} - e^{i\pi}\big) + b\big(1 - e^{2i\pi/3}\big), \]
and the centre is $\omega = \dfrac{z'+z}{2}\Big|_{z=\omega}$, giving $\omega = \dfrac{a\,e^{2i\pi/3} + a + b\big(1 - e^{2i\pi/3}\big)}{2}$. Geometrically, $\Omega$ is the point such that the triangle built on $A, B, \Omega$ encodes the half-angles: $\Omega$ is the vertex of the triangle with $\widehat{\Omega A B} = 30^{\circ}$ and $\widehat{AB\Omega} = 60^{\circ}$ (half-angles at $A$ and $B$), hence a right angle at $\Omega$.
\end{enumerate}
\end{corrige}

\begin{corrige}[Exercise 7 — Recognising an isometry from images]
\begin{enumerate}
\item Compute the displacement vectors:
\[ \overrightarrow{AA'} = (2,1), \qquad \overrightarrow{BB'} = (5-3, 1-0) = (2,1), \qquad \overrightarrow{CC'} = (-2-0, 5-4) = (2,1). \]
All three displacement vectors are equal to $(2,1)$. Since an isometry mapping three non-collinear points is uniquely determined, $f$ coincides with the translation of vector $\vec{u}(2,1)$ on the whole plane. Check distances: $AB = 3 = A'B'$, $AC = 4 = A'C'$, consistent. Hence $f = t_{\vec{u}}$ with $\vec{u}(2,1)$.
\item Now $C' = (6,1)$. We still have $A'(2,1)$, $B'(5,1)$. Note $A'B' = 3 = AB$ and $A'C' = \sqrt{36+1} = \sqrt{37}$, while $AC = 4$: distances are \emph{not} preserved, so \textbf{no isometry} maps $A,B,C$ to $A',B',C'$ in this case. (Indeed $B'C' = \sqrt{1+0} = 1$ while $BC = 5$.) Therefore such an isometry $f$ does not exist; the data are inconsistent with an isometry.
\end{enumerate}
\textbf{Remark.} The examiner's trap: always verify that all pairwise distances are preserved before concluding. Three image points that fail the distance test cannot come from an isometry.
\end{corrige}

\begin{corrige}[Exercise 8 — Homothety from its complex expression]
\begin{enumerate}
\item The expression $z' = 3z - 4 + 2i$ has the form $z' = kz + b$ with $k = 3 \in \mathbb{R}\setminus\{1\}$: this is a homothety of ratio $3$. Its centre is the unique fixed point: $\omega = 3\omega - 4 + 2i \Rightarrow 2\omega = 4 - 2i \Rightarrow \omega = 2 - i$. Indeed $z' - (2-i) = 3\big(z - (2-i)\big)$. \textbf{Centre $\Omega(2,-1)$, ratio $k = 3$.}
\item A homothety of ratio $k$ maps a circle of centre $C$ and radius $R$ to the circle of centre $C' = f(C)$ and radius $|k|R$. Centre: $C$ has affix $1+i$, so $c' = 3(1+i) - 4 + 2i = -1 + 5i$. Radius: $3 \times 2 = 6$. \textbf{Image: circle of centre $(-1,5)$ and radius $6$.}
\item The line $d: z + \bar{z} = 4$ is the vertical line $x = 2$ (since $z+\bar z = 2\,\mathrm{Re}(z)$). A homothety of ratio $3$ maps a line not through the centre to a parallel line. Write $z = 2 + iy$; then $z' = 3(2+iy) - 4 + 2i = 2 + (3y+2)i$, whose real part is $2$. Wait — check: the centre $\Omega = 2-i$ has real part $2$, so the line $x = 2$ passes \emph{through the centre}! A line through the centre of a homothety is \textbf{globally invariant}: its image is the line itself, $x = 2$, i.e.\ $z + \bar{z} = 4$. (Each point is mapped to another point of the same line.)
\end{enumerate}
\end{corrige}

\begin{corrige}[Exercise 9 — Composing homotheties (8 marks)]
\begin{enumerate}
\item \textbf{(3 marks)} For any point $M$, set $M_1 = h_1(M)$ and $M' = h_2(M_1)$. Then
\[ \overrightarrow{O_1 M_1} = 2\overrightarrow{O_1 M}, \qquad \overrightarrow{O_2 M'} = \tfrac12 \overrightarrow{O_2 M_1}. \]
Compute $\overrightarrow{MM'} = \overrightarrow{MO_1} + \overrightarrow{O_1O_2} + \overrightarrow{O_2M'}$. Alternatively use affixes with $o_1, o_2$: $z_1 - o_1 = 2(z-o_1)$ so $z_1 = 2z - o_1$; then $z' - o_2 = \tfrac12(z_1 - o_2)$ gives
\[ z' = \tfrac12(2z - o_1 - o_2) + o_2 = z + \tfrac12(o_2 - o_1). \]
This is the translation of vector $\dfrac12 \overrightarrow{O_1O_2}$. Since $O_1 \neq O_2$ the vector is non-zero. \textbf{Mark scheme:} 1 mark for writing both homothety relations; 1 mark for correct composition to $z' = z + b$ (or vector form); 1 mark for identifying the translation and its vector $\frac12\overrightarrow{O_1O_2}$.
\item \textbf{(1 mark)} Draw $O_1$ and $O_2$; the vector of the translation is the vector from $O_1$ to the midpoint of $[O_1O_2]$ (half of $\overrightarrow{O_1O_2}$, same direction). 1 mark for a correct arrow of length $\frac{O_1O_2}{2}$ parallel to $(O_1O_2)$.
\item \textbf{(4 marks)} Now $h_2 = h(O_2,3)$: $z_1 = 2z - o_1$ and $z' - o_2 = 3(z_1 - o_2)$, so
\[ z' = 3(2z - o_1 - o_2) + o_2 = 6z - 3o_1 - 2o_2. \]
This has the form $z' = Kz + b$ with $K = 6 \in \mathbb{R}\setminus\{1\}$: a \textbf{homothety of ratio $k_1k_2 = 6$}. Its centre $\Omega$ is the fixed point: $\omega = 6\omega - 3o_1 - 2o_2 \Rightarrow 5\omega = 3o_1 + 2o_2$, i.e.
\[ \omega = \frac{3o_1 + 2o_2}{5}, \qquad \overrightarrow{O_1\Omega} = \frac{2}{5}\overrightarrow{O_1O_2}. \]
So $\Omega$ lies on the segment $[O_1O_2]$, at $\dfrac{2}{5}$ of the way from $O_1$ to $O_2$ (equivalently $\overrightarrow{\Omega O_2} = \frac35\overrightarrow{O_1O_2}$... precisely $O_1\Omega = \frac25 O_1O_2$). \textbf{Mark scheme:} 1 mark for the composite expression; 1 mark for ratio $6$; 1 mark for the fixed-point equation; 1 mark for the position $\overrightarrow{O_1\Omega} = \frac25\overrightarrow{O_1O_2}$ on line $(O_1O_2)$.
\end{enumerate}
\end{corrige}

\begin{corrige}[Exercise 10 — Centres of similitude of two circles (10 marks)]
\begin{enumerate}
\item \textbf{(3 marks)} A homothety mapping $\mathcal{C}(A,2)$ onto $\mathcal{C}'(B,6)$ must send the centre $A$ to the centre $B$ (centres correspond) and multiply radii by $|k|$: $|k| = \dfrac{6}{2} = 3$, so $k = 3$ or $k = -3$. For each value, the centre $\Omega$ is determined by $\overrightarrow{\Omega B} = k\overrightarrow{\Omega A}$, which has exactly one solution since $k \neq 1$. Hence \textbf{exactly two homotheties}, of ratios $\boxed{3}$ and $\boxed{-3}$.
\item \textbf{(3 marks)} Let $\Omega$ on line $(AB)$ with $\overrightarrow{\Omega B} = k\overrightarrow{\Omega A}$.
For $k = 3$: $B - \Omega = 3(A - \Omega)$ gives $2\Omega = 3A - B$, i.e.\ $\overrightarrow{A\Omega} = -\dfrac{1}{2}\overrightarrow{AB}$: $\Omega_1$ is on line $(AB)$, on the side of $A$ opposite to $B$, with $\Omega_1 A = \dfrac{AB}{2} = 5$ cm and $\Omega_1 B = 15$ cm. (External centre: $\dfrac{\Omega_1 A}{\Omega_1 B} = \dfrac{2}{6}$.)
For $k = -3$: $B - \Omega = -3(A - \Omega)$ gives $4\Omega = 3A + B$, i.e.\ $\overrightarrow{A\Omega_2} = \dfrac{1}{4}\overrightarrow{AB}$: $\Omega_2$ lies \emph{between} $A$ and $B$ with $\Omega_2 A = \dfrac{AB}{4} = 2.5$ cm and $\Omega_2 B = 7.5$ cm. (Internal centre, dividing $[AB]$ in the ratio of the radii $2:6$.)
\item \textbf{(2 marks)} Let $t$ be a common external tangent touching $\mathcal{C}$ at $T$ and $\mathcal{C}'$ at $T'$. The radii $AT$ and $BT'$ are both perpendicular to $t$, hence parallel, pointing to the same side. The homothety $h(\Omega_1, 3)$ sends $A$ to $B$ and, since $\overrightarrow{BT'} = 3\overrightarrow{AT}$ (same direction, lengths $6 = 3\times 2$), sends $T$ to $T'$; therefore $T' = h(T)$ lies on the line $(\Omega_1 T)$: the points $\Omega_1, T, T'$ are collinear, i.e.\ the tangent line at $T$ and $T'$ passes through $\Omega_1$. (Similarly the common \emph{internal} tangents pass through $\Omega_2$.)
\item \textbf{(2 marks)} Construction: draw $(AB)$, mark $\Omega_1$ with $A$ between $\Omega_1$ and $B$, $\Omega_1A = 5$ cm; mark $\Omega_2$ between $A$ and $B$ with $A\Omega_2 = 2.5$ cm. For the external tangents from $\Omega_1$: draw the circle of diameter $[\Omega_1 A]$ meeting $\mathcal{C}$ at the tangent points $T_1, T_2$; the lines $(\Omega_1 T_i)$ are the external tangents (also tangent to $\mathcal{C}'$). \textbf{Mark scheme:} 1 mark for each centre correctly placed, tangents drawn with tangent points constructed via a right angle.
\end{enumerate}
\end{corrige}

\begin{corrige}[Exercise 11 — Euler line and direct similarity (12 marks)]
\textbf{Part A.}
\begin{enumerate}
\item \textbf{(3 marks)} Let $M$ be the midpoint of $[BC]$. Since $G$ is the centroid, $\overrightarrow{GA} = -2\overrightarrow{GM}$ (the centroid divides the median in ratio $2:1$ from the vertex). Hence $h(G,-2)$ maps $M$ onto $A$. The perpendicular bisector of $[BC]$ passes through $O$ and is perpendicular to $[BC]$ at $M$; its image under $h(G,-2)$ is the line through $A$ parallel to it (a homothety maps a line to a parallel line), i.e.\ the line through $A$ perpendicular to $[BC]$ — the \emph{altitude} from $A$. Doing the same with the other sides, $O$ (intersection of the perpendicular bisectors) is mapped to the intersection of the three altitudes, which is $H$. Hence $h(G,-2)(O) = H$.
\item \textbf{(2 marks)} Since $h(G,-2)$ maps $O$ to $H$, we have $\overrightarrow{GH} = -2\overrightarrow{GO}$; the vectors are collinear, so $O, G, H$ lie on one line (the \textbf{Euler line}), and $GH = 2\,GO$ with $G$ between $O$ and $H$.
\end{enumerate}
\textbf{Part B.}
\begin{enumerate}
\item \textbf{(3 marks)} $z' = az + b$ with $a = 1 + i\sqrt3$: this is a direct similarity of ratio $k = |a| = \sqrt{1+3} = 2$ and angle $\theta = \arg a = \dfrac{\pi}{3}$. Since $k = 2 \neq 1$, $S$ is \textbf{the composite of the homothety of ratio $2$ and the rotation of angle $\frac{\pi}{3}$ about the same centre} (a direct similarity with centre, ratio $2$, angle $\frac{\pi}{3}$ — sometimes called a spiral similarity).
\item \textbf{(2 marks)} The centre $\Omega$ is the fixed point: $\omega = a\omega + b$ with $b = 2 - i$:
\[ \omega = \frac{b}{1-a} = \frac{2 - i}{1 - (1+i\sqrt3)} = \frac{2-i}{-i\sqrt3} = \frac{(2-i)\,i}{ \sqrt3} = \frac{1 + 2i}{\sqrt3} = \frac{\sqrt3}{3} + \frac{2\sqrt3}{3}i. \]
\item \textbf{(1 mark)} Image of $z = 1$: $z' = (1+i\sqrt3)(1) + 2 - i = 3 + i(\sqrt3 - 1)$.
\item \textbf{(1 mark)} A glide reflection has complex form $z' = e^{i\alpha}\bar z + c$ (indirect isometry); composing with a homothety of ratio $k \neq 1$ gives $z' = a\bar z + b$ with $|a| = k \neq 1$: an \textbf{indirect similarity}. Its fixed points satisfy $z = a\bar z + b$. Conjugating gives $\bar z = \bar a z + \bar b$. Substituting: $z = a(\bar a z + \bar b) + b = |a|^2 z + a\bar b + b$, hence $(1-|a|^2)z = b + a\bar b$. Since $|a| = k \neq 1$, this linear equation has a \textbf{unique solution} $z = \dfrac{b + a\bar b}{1-k^2}$. Illustration: $z' = 2\bar z + 3i$ is an indirect similarity of ratio $2$; fixed point: $z = 2\bar z + 3i$; conjugate $\bar z = 2z - 3i$; substitute: $z = 2(2z-3i)+3i = 4z - 3i \Rightarrow z = i$: unique fixed point $i$.
\end{enumerate}
\textbf{Examiner expectations.} Part A rewards the median relation $\overrightarrow{GA} = -2\overrightarrow{GM}$ (1 mark), the image of the perpendicular bisector being the altitude (1 mark), the conclusion $O \mapsto H$ (1 mark), then collinearity and the ratio (2 marks). Part B: correct modulus/argument (2 marks) and classification (1 mark); fixed-point computation with correct conjugate manipulation (2 marks); image computation (1 mark); the final proof requires the general form $z' = a\bar z + b$ and the conjugate substitution argument leading to a unique solution (1 mark).
\end{corrige}

\newpage

\coursec{Summary sheet}

\courssub{Chapter at a Glance}

This chapter completes your study of plane transformations begun in Lower Sixth. You already know the four basic \cle{isometries} (translation, reflection, rotation, central symmetry); here you will learn what happens when they are \emph{composed}, how every plane isometry can be \emph{classified}, and you will meet two new families of transformations: \cle{homotheties} and, more generally, \cle{plane similarities}, studied both geometrically and analytically (with coordinates and with complex numbers). The mind map below shows the route we shall follow.

\begin{popfigure}
\begin{tikzpicture}[
  mindmap,
  grow cyclic,
  every node/.style={concept, font=\small},
  root concept/.append style={concept color=popblue, font=\bfseries},
  level 1/.append style={level distance=3.6cm, sibling angle=45},
  level 1 concept/.append style={font=\scriptsize},
]
\node[root concept, align=center]{Plane\\Transformations}
  child[concept color=poporange]{ node[align=center]{Composing\\Reflections \&\\Translations} }
  child[concept color=popgreen]{ node[align=center]{Composing\\Rotations;\\Isometries} }
  child[concept color=poppurple]{ node[align=center]{Classification\\of Isometries;\\Glide Reflections} }
  child[concept color=popgold]{ node[align=center]{Homothety:\\Definition \&\\Properties} }
  child[concept color=poppink]{ node[align=center]{Analytical \&\\Complex\\Representation} }
  child[concept color=popblue]{ node[align=center]{Composition\\of Homotheties} }
  child[concept color=popgreen]{ node[align=center]{Plane\\Similarities} }
  child[concept color=poporange]{ node[align=center]{Analytical\\Study \&\\Applications} };
\end{tikzpicture}
\legende{Mind map of the chapter Plane Transformations: from composing isometries to homotheties and similarities.}
\end{popfigure}

\courssub{Summary Sheet}

\begin{bilan}
\textbf{Key definitions}\par
\begin{itemize}
\item A \cle{reflection} (orthogonal symmetry) $S_\Delta$ maps $M$ to $M'$ where $\Delta$ is the perpendicular bisector of $[MM']$; points of $\Delta$ are fixed.
\item A \cle{translation} $t_{\vec{u}}$ maps $M$ to $M'$ with $\overrightarrow{MM'}=\vec{u}$; it has no fixed point unless $\vec{u}=\vec{0}$ (the identity).
\item A \cle{rotation} $r(\Omega,\theta)$ maps $M$ to $M'$ with $\Omega M=\Omega M'$ and $\mathrm{mes}(\overrightarrow{\Omega M},\overrightarrow{\Omega M'})=\theta$; its only fixed point is $\Omega$ (unless $\theta=0\ [2\pi]$).
\item A \cle{central symmetry} is the rotation of angle $\pi$, i.e.\ the homothety of ratio $-1$.
\item A \cle{glide reflection} is the composite $S_\Delta\circ t_{\vec{u}}$ with $\vec{u}$ parallel to $\Delta$; it has no fixed point and reverses orientation.
\item A \cle{homothety} $h(\Omega,k)$, $k\neq 0$, maps $M$ to $M'$ with $\overrightarrow{\Omega M'}=k\,\overrightarrow{\Omega M}$; it is \cle{direct} if $k>0$, \cle{indirect} if $k<0$.
\item A \cle{plane similarity} of ratio $k>0$ multiplies all distances by $k$; it is \cle{direct} (preserves oriented angles) or \cle{indirect} (reverses oriented angles).
\end{itemize}

\textbf{Laws, formulas and theorems}\par
\begin{itemize}
\item $S_{\Delta_2}\circ S_{\Delta_1}$ with $\Delta_1\parallel\Delta_2$: translation by $2\,\vec{u}$, where $\vec{u}$ is the vector sending $\Delta_1$ to $\Delta_2$ perpendicularly.
\item $S_{\Delta_2}\circ S_{\Delta_1}$ with $\Delta_1\cap\Delta_2=\{\Omega\}$: rotation of centre $\Omega$ and angle $2(\Delta_1,\Delta_2)$.
\item $t_{\vec{v}}\circ t_{\vec{u}}=t_{\vec{u}+\vec{v}}$: composition of translations is commutative.
\item $r(\Omega,\theta')\circ r(\Omega,\theta)=r(\Omega,\theta+\theta')$; for distinct centres, the composite is a rotation of angle $\theta+\theta'$ unless $\theta+\theta'=0\ [2\pi]$, in which case it is a translation.
\item $S_\Delta\circ t_{\vec{u}}$ with $\vec{u}$ \emph{not} parallel to $\Delta$ is still a glide reflection: decompose $\vec{u}=\vec{u}_1+\vec{u}_2$ with $\vec{u}_1\parallel\Delta$, $\vec{u}_2\perp\Delta$, and absorb $t_{\vec{u}_2}$ by shifting the axis.
\item Classification theorem: every plane isometry is a translation, a rotation, a reflection or a glide reflection; the first two are \cle{direct} (preserve orientation), the last two are \cle{indirect} (reverse orientation).
\item Homothety $h(\Omega,k)$: $M'N'=|k|\,MN$; areas are multiplied by $k^2$; the image of a line is a parallel line; the image of a circle of radius $R$ is a circle of radius $|k|R$, the centre being the image of the centre.
\item $h(\Omega,k_2)\circ h(\Omega,k_1)=h(\Omega,k_1k_2)$; with distinct centres and $k_1k_2\neq 1$, the composite is a homothety of ratio $k_1k_2$ whose centre lies on the line of the two centres; if $k_1k_2=1$, it is a translation.
\item Complex form: homothety of centre $\Omega(a)$ and ratio $k\in\mathbb{R}^*$: $z'=kz+a(1-k)$; direct similarity: $z'=az+b$ ($|a|$ is the ratio, $\arg a$ the angle); indirect similarity: $z'=a\bar{z}+b$.
\item A direct similarity with ratio $k\neq 1$ has a unique fixed point $\Omega$ (its centre) and equals $r(\Omega,\theta)\circ h(\Omega,k)=h(\Omega,k)\circ r(\Omega,\theta)$; it is called a \cle{spiral similarity} when $k\neq 1$ and $\theta\neq 0\ [2\pi]$.
\end{itemize}

\textbf{Methods}\par
\begin{itemize}
\item To identify a composite: count the fixed points, check whether orientation is preserved or reversed, then test candidate transformations on two well-chosen points.
\item To construct $M'=h(\Omega,k)(M)$: draw the line $(\Omega M)$ and mark $M'$ with $\overrightarrow{\Omega M'}=k\overrightarrow{\Omega M}$ (same side as $M$ if $k>0$, opposite side if $k<0$).
\item To find the centre of a homothety given $A\mapsto A'$ and $B\mapsto B'$: the centre is the intersection of $(AA')$ and $(BB')$; the ratio is $k=\dfrac{\overrightarrow{\Omega A'}}{\overrightarrow{\Omega A}}$.
\item To find the centre of a spiral similarity sending $A\mapsto A'$, $B\mapsto B'$: the centre lies at the intersection of the circles through $(A,B,P)$ and $(A',B',P)$, where $P=(AB)\cap(A'B')$.
\item In analytical work: write $z'=az+b$, compute $a$ and $b$ from two point-image pairs, then read off the nature from $a$ ($a\in\mathbb{R}$: homothety or translation; $|a|=1$: isometry).
\end{itemize}

\textbf{Common mistakes}\par
\begin{itemize}
\item Forgetting the factor $2$ in the angle $2(\Delta_1,\Delta_2)$ of the composite of two reflections.
\item Believing composition is always commutative: in general $g\circ f\neq f\circ g$ (order matters; $g\circ f$ means $f$ first, then $g$).
\item Confusing the ratio $k$ with $|k|$: distances are multiplied by $|k|$, areas by $k^2$.
\item Forgetting the special cases: $k=1$ gives the identity, $k=-1$ gives the central symmetry about $\Omega$.
\item Forgetting that when $\theta+\theta'=0\ [2\pi]$ the composite of two rotations is a translation, not a rotation; similarly $k_1k_2=1$ for homotheties with distinct centres gives a translation.
\item Claiming a glide reflection has a fixed point: it has none, although its axis is globally invariant.
\end{itemize}

\textbf{GCE tips}\par
\begin{itemize}
\item Always state the \cle{characteristic elements} of a transformation: centre, ratio, angle, axis or vector, as appropriate.
\item In complex-number questions, write the form $z'=az+b$ first, then read off the nature of the transformation from $a$; this earns method marks.
\item Show constructions with dashed lines, mark equal lengths and angles, and justify each step; method marks are awarded for correct reasoning even if the final identification slips.
\item Check your answer on a simple figure (a triangle or a square) before generalising; a quick sketch often reveals a wrong ratio or a wrong angle.
\end{itemize}
\end{bilan}

\findecours

\end{document}
